% Generated mechanically from frozen input inputs/p07/ja/derham.tex.
% Do not edit this generated file; edit the bound locale source instead.

% OK, start here
%

\setcounter{chapter}{49}
\chapter{ド・ラームコホモロジー}



\phantomsection
\label{derham-section-phantom}


\section{序論}
\label{derham-section-introduction}

\noindent
本章では、まずスキームの射のド・ラーム複体を論じ、最後に、基礎体の標数が
零であるときド・ラームコホモロジーがヴェイユコホモロジー理論を定めることを証明する。




\section{ド・ラーム複体}
\label{derham-section-de-rham-complex}

\noindent
$p : X \to S$ をスキームの射とする。このとき複体
$$
\Omega^\bullet_{X/S} =
\mathcal{O}_{X/S} \to \Omega^1_{X/S} \to \Omega^2_{X/S} \to \ldots
$$
は $p^{-1}\mathcal{O}_S$-加群の複体であり、
$\Omega^i_{X/S} = \wedge^i(\Omega_{X/S})$
は次数 $i$ に置かれ、その微分は局所切断に対する規則
$\text{d}(g_0 \text{d}g_1 \wedge \ldots \wedge \text{d}g_p) =
\text{d}g_0 \wedge \text{d}g_1 \wedge \ldots \wedge \text{d}g_p$
によって定まる。
Modules, Section \ref{modules-section-de-rham-complex} を参照せよ。

\medskip\noindent
スキームの可換図式
$$
\xymatrix{
X' \ar[r]_f \ar[d] & X \ar[d] \\
S' \ar[r] & S
}
$$
が与えられると、複体の標準的な射
$f^{-1}\Omega_{X/S}^\bullet \to \Omega^\bullet_{X'/S'}$ および
$\Omega_{X/S}^\bullet \to f_*\Omega^\bullet_{X'/S'}$ がある。
Modules, Section \ref{modules-section-de-rham-complex} を参照せよ。
線形化することにより、各 $p$ に対して線形写像
$f^*\Omega^p_{X/S} \to \Omega^p_{X'/S'}$
を得る。

\medskip\noindent
特に、$f : Y \to X$ が基礎スキーム $S$ 上のスキームの射であるならば、
複体の射
$$
\Omega^\bullet_{X/S} \longrightarrow f_*\Omega^\bullet_{Y/S}
$$
がある。線形化すると、各 $p \geq 0$ に対して標準的な射
$$
\Omega^p_{X/S} \otimes_{\mathcal{O}_X} f_*\mathcal{O}_Y
\longrightarrow
f_*\Omega^p_{Y/S}
$$

を得る。

\begin{lemma}
\label{derham-lemma-base-change-de-rham}
次の図式を
$$
\xymatrix{
X' \ar[r]_f \ar[d] & X \ar[d] \\
S' \ar[r] & S
}
$$
スキームのカルテジアン図式とする。このとき、上で論じた射は同型
$f^*\Omega^p_{X/S} \to \Omega^p_{X'/S'}$
を誘導する。
\end{lemma}

\begin{proof}
Morphisms, Lemma \ref{morphisms-lemma-base-change-differentials} と、
外冪の形成が基底変換と可換であることを組み合わせればよい。
\end{proof}

\begin{lemma}
\label{derham-lemma-etale}
スキームの可換図式
$$
\xymatrix{
X' \ar[r]_f \ar[d] & X \ar[d] \\
S' \ar[r] & S
}
$$
を考える。$X' \to X$ と $S' \to S$ がエタールならば、上で論じた射は同型
$f^*\Omega^p_{X/S} \to \Omega^p_{X'/S'}$
を誘導する。
\end{lemma}

\begin{proof}
例えば Morphisms, Lemma \ref{morphisms-lemma-etale-at-point} により、
$\Omega_{S'/S} = 0$ かつ $\Omega_{X'/X} = 0$ である。したがって、
Morphisms, Lemmas
\ref{morphisms-lemma-triangle-differentials} および
\ref{morphisms-lemma-triangle-differentials-smooth}
の短完全列から $\Omega_{X'/S'} = \Omega_{X'/S} = f^*\Omega_{X/S}$ を得る。
外冪をとれば結論が従う。
\end{proof}






\section{ド・ラームコホモロジー}
\label{derham-section-de-rham-cohomology}

\noindent
$p : X \to S$ をスキームの射とする。{\it $X$ の、$S$ に関する
ド・ラームコホモロジー}を、コホモロジー群
$$
H^i_{dR}(X/S) = H^i(R\Gamma(X, \Omega^\bullet_{X/S}))
$$
によって定義する。$\Omega^\bullet_{X/S}$ は $p^{-1}\mathcal{O}_S$-加群の
複体であるから、これらのコホモロジー群は自然に
$H^0(S, \mathcal{O}_S)$ 上の加群となる。

\medskip\noindent
可換図式
$$
\xymatrix{
X' \ar[r]_f \ar[d] & X \ar[d] \\
S' \ar[r] & S
}
$$
が与えられると、Section \ref{derham-section-de-rham-complex} の標準的な射を用いて、
引き戻し写像
$$
f^* :
R\Gamma(X, \Omega^\bullet_{X/S})
\longrightarrow
R\Gamma(X', \Omega^\bullet_{X'/S'})
$$
および
$$
f^* : H^i_{dR}(X/S) \longrightarrow H^i_{dR}(X'/S')
$$
を得る。これらの引き戻しは明らかな合成則を満たす。
特に、基礎スキーム $S$ を固定して考えると、ド・ラームコホモロジーは
$S$ 上のスキームの圏上の反変関手である。

\begin{lemma}
\label{derham-lemma-de-rham-affine}
$X \to S$ を、環準同型 $R \to A$ により与えられる
アフィンスキームの射とする。このとき
$R\Gamma(X, \Omega^\bullet_{X/S}) = \Omega^\bullet_{A/R}$ は $D(R)$ において成り立ち、
$H^i_{dR}(X/S) = H^i(\Omega^\bullet_{A/R})$ である。
\end{lemma}

\begin{proof}
これは Cohomology of Schemes, Lemma
\ref{coherent-lemma-quasi-coherent-affine-cohomology-zero}
および Leray の非輪状性補題
(Derived Categories, Lemma \ref{derived-lemma-leray-acyclicity}) から従う。
\end{proof}

\begin{lemma}
\label{derham-lemma-quasi-coherence-relative}
$p : X \to S$ をスキームの射とする。$p$ が準コンパクトかつ準分離ならば、
$Rp_*\Omega^\bullet_{X/S}$ は $D_\QCoh(\mathcal{O}_S)$ の対象である。
\end{lemma}

\begin{proof}
第1頁が $E_1^{a, b} = R^bp_*\Omega^a_{X/S}$ であり、
$Rp_*\Omega^\bullet_{X/S}$ のコホモロジーに収束するスペクトル系列がある
(Derived Categories, Lemma \ref{derived-lemma-two-ss-complex-functor} を参照)。
したがって Homology, Lemma \ref{homology-lemma-first-quadrant-ss} により、
$R^bp_*\Omega^a_{X/S}$ が準連接であることを示せば十分である。
これは Cohomology of Schemes, Lemma
\ref{coherent-lemma-quasi-coherence-higher-direct-images} から従う。
\end{proof}

\begin{lemma}
\label{derham-lemma-coherence-relative}
$p : X \to S$ をスキームの固有射とし、$S$ は局所ネーターであるとする。
このとき $Rp_*\Omega^\bullet_{X/S}$ は
$D_{\textit{Coh}}(\mathcal{O}_S)$ の対象である。
\end{lemma}

\begin{proof}
この場合、Morphisms, Lemma \ref{morphisms-lemma-finite-type-differentials}
により、加群 $\Omega^i_{X/S}$ は連接である。したがって、
Lemma \ref{derham-lemma-quasi-coherence-relative} の証明と全く同じ議論を、
Cohomology of Schemes, Proposition
\ref{coherent-proposition-proper-pushforward-coherent} を用いて適用できる。
\end{proof}

\begin{lemma}
\label{derham-lemma-finite-de-Rham}
$A$ をネーター環とし、$X$ を $S = \Spec(A)$ 上の固有スキームとする。
このとき $H^i_{dR}(X/S)$ は有限生成 $A$-加群であり、これはすべての $i$ に対して成り立つ。
\end{lemma}

\begin{proof}
これは Lemma \ref{derham-lemma-coherence-relative} の特別な場合である。
\end{proof}

\begin{lemma}
\label{derham-lemma-proper-smooth-de-Rham}
$f : X \to S$ をスキームの固有かつ滑らかな射とする。このとき、
$Rf_*\Omega^p_{X/S}$（$p \geq 0$）および $Rf_*\Omega^\bullet_{X/S}$ は
$D(\mathcal{O}_S)$ の完全対象であり、その形成は任意の基底変換と可換である。
\end{lemma}

\begin{proof}
$f$ は滑らかであるから、$\Omega^p_{X/S}$ は有限局所自由
$\mathcal{O}_X$-加群である。Morphisms, Lemma
\ref{morphisms-lemma-smooth-omega-finite-locally-free} を参照せよ。
Lemma \ref{derham-lemma-base-change-de-rham} により、その形成は任意の基底変換と可換である。
したがって、$Rf_*\Omega^p_{X/S}$ は $D(\mathcal{O}_S)$ の完全対象であり、
その形成は任意の基底変換と可換である。Derived Categories of Schemes, Lemma
\ref{perfect-lemma-flat-proper-perfect-direct-image-general} を参照せよ。
これで補題の第1の主張が証明された。

\medskip\noindent
$Rf_*\Omega^\bullet_{X/S}$ が $S$ 上完全であることを示すには、$S$ 上局所的に
議論してよい。したがって $S$ は準コンパクトであると仮定してよく、このとき
$\Omega^n_{X/S}$ は十分大きい $n$ に対して零であると仮定できる。各 $p \geq 0$ に対し、
$Rf_*\sigma_{\geq p}\Omega^\bullet_{X/S}$ は $D(\mathcal{O}_S)$ の完全対象で、
その形成は任意の基底変換と可換であると主張する。上の議論から、これは
$p \gg 0$ に対して成り立つ。$p$ に対して主張が成り立つと仮定し、区別三角
$$
\sigma_{\geq p}\Omega^\bullet_{X/S} \to
\sigma_{\geq p - 1}\Omega^\bullet_{X/S} \to
\Omega^{p - 1}_{X/S}[-(p - 1)] \to
(\sigma_{\geq p}\Omega^\bullet_{X/S})[1]
$$
を $D(f^{-1}\mathcal{O}_S)$ において考える。完全関手 $Rf_*$ を適用すると、
$D(\mathcal{O}_S)$ における区別三角を得る。完全性については 2-out-of-3 性があるので
(Cohomology, Lemma \ref{cohomology-lemma-two-out-of-three-perfect})、
$Rf_*\sigma_{\geq p - 1}\Omega^\bullet_{X/S}$ は $D(\mathcal{O}_S)$ の
完全対象であると結論できる。任意の基底変換との可換性についても同様である。
\end{proof}





\section{カップ積}
\label{derham-section-cup-product}

\noindent
写像
$\Omega^p_{X/S} \times \Omega^q_{X/S} \to \Omega^{p + q}_{X/S}$
を $(\omega , \eta) \longmapsto \omega \wedge \eta$ によって定める。
Section \ref{derham-section-de-rham-complex} で与えた $\text{d}$ の公式と、
$\text{d} : \mathcal{O}_X \to \Omega_{X/S}$ に対する Leibniz 則から、
$\text{d}(\omega \wedge \eta) = \text{d}(\omega) \wedge \eta +
(-1)^{\deg(\omega)} \omega \wedge \text{d}(\eta)$ が分かる。したがって、
$\wedge$ は射
\begin{equation}
\label{derham-equation-wedge}
\wedge :
\text{Tot}(
\Omega^\bullet_{X/S} \otimes_{p^{-1}\mathcal{O}_S} \Omega^\bullet_{X/S})
\longrightarrow
\Omega^\bullet_{X/S}
\end{equation}
を $p^{-1}\mathcal{O}_S$-加群の複体の射として定める。

\medskip\noindent
Cohomology, Section \ref{cohomology-section-cup-product} のカップ積と
(\ref{derham-equation-wedge}) を組み合わせると、$H^0(S, \mathcal{O}_S)$-双線形な
カップ積写像
$$
\cup : H^i_{dR}(X/S) \times H^j_{dR}(X/S) \longrightarrow H^{i + j}_{dR}(X/S)
$$
を得る。例えば、$\omega \in \Gamma(X, \Omega^i_{X/S})$ と
$\eta \in \Gamma(X, \Omega^j_{X/S})$ が閉形式ならば、
$\omega$ と $\eta$ のド・ラームコホモロジー類のカップ積は
$\omega \wedge \eta$ のド・ラームコホモロジー類である。
Cohomology, Section \ref{cohomology-section-cup-product} の議論を参照せよ。

\medskip\noindent
可換図式
$$
\xymatrix{
X' \ar[r]_f \ar[d] & X \ar[d] \\
S' \ar[r] & S
}
$$
が与えられると、引き戻し写像
$f^* : R\Gamma(X, \Omega^\bullet_{X/S}) \to R\Gamma(X', \Omega^\bullet_{X'/S'})$
および
$f^* : H^i_{dR}(X/S) \longrightarrow H^i_{dR}(X'/S')$
は上で定義したカップ積と両立する。

\begin{lemma}
\label{derham-lemma-cup-product-graded-commutative}
$p : X \to S$ をスキームの射とする。
$H^*_{dR}(X/S)$ 上のカップ積は結合的かつ次数付き可換である。
\end{lemma}

\begin{proof}
これは Cohomology, Lemmas \ref{cohomology-lemma-cup-product-associative} および
\ref{cohomology-lemma-cup-product-commutative}
と、$\wedge$ が結合的かつ次数付き可換であることから従う。
\end{proof}

\begin{remark}
\label{derham-remark-relative-cup-product}
$p : X \to S$ をスキームの射とする。このとき $\Omega^\bullet_{X/S}$ は
微分次数付き $p^{-1}\mathcal{O}_S$-代数の層とみなせる。
Differential Graded Sheaves, Definition \ref{sdga-definition-dga} を参照せよ。
特に Differential Graded Sheaves, Section \ref{sdga-section-misc} の議論が
適用できる。例えば、任意のスキームの可換図式
$$
\xymatrix{
X \ar[d]_p \ar[r]_f & Y \ar[d]^q \\
S \ar[r]^h & T
}
$$
に対して標準的な相対カップ積
$$
\mu :
Rf_*\Omega^\bullet_{X/S}
\otimes_{q^{-1}\mathcal{O}_T}^\mathbf{L}
Rf_*\Omega^\bullet_{X/S}
\longrightarrow
Rf_*\Omega^\bullet_{X/S}
$$
が $D(Y, q^{-1}\mathcal{O}_T)$ に存在し、これは結合的であり、
コホモロジー上では上で論じたカップ積を再現する。
\end{remark}

\begin{remark}
\label{derham-remark-cup-product-as-a-map}
$f : X \to S$ をスキームの射とし、$\xi \in H_{dR}^n(X/S)$ とする。
Differential Graded Sheaves, Section \ref{sdga-section-misc} の議論により、
標準的な射
$$
\xi' : \Omega^\bullet_{X/S} \to \Omega^\bullet_{X/S}[n]
$$
が $D(f^{-1}\mathcal{O}_S)$ に存在する。これは次の性質のうち
(1) と (2) によって一意に特徴付けられる。
\begin{enumerate}
\item $\xi'$ は右微分次数付き $\Omega^\bullet_{X/S}$-加群の導来圏における
写像へ持ち上げられる。
\item $\xi'(1) = \xi$ が
$H^0(X, \Omega^\bullet_{X/S}[n]) = H^n_{dR}(X/S)$ において成り立つ。
\item 写像 $\xi'$ は $\eta \in H^m_{dR}(X/S)$ を
$\xi \cup \eta$ に送る。これは $H^{n + m}_{dR}(X/S)$ の元である。
\item $\xi'$ の構成は開集合への制限と可換である。すなわち、
$U \subset X$ が開ならば、制限 $\xi'|_U$ は像
$\xi|_U \in H^n_{dR}(U/S)$ に対応する写像である。
\item Remark \ref{derham-remark-relative-cup-product} にある任意の図式に対して、
可換図式
$$
\xymatrix{
Rf_*\Omega^\bullet_{X/S}
\otimes_{q^{-1}\mathcal{O}_T}^\mathbf{L}
Rf_*\Omega^\bullet_{X/S} \ar[d]_{\xi' \otimes \text{id}}
\ar[r]_-\mu &
Rf_*\Omega^\bullet_{X/S} \ar[d]^{\xi'} \\
Rf_*\Omega^\bullet_{X/S}[n]
\otimes_{q^{-1}\mathcal{O}_T}^\mathbf{L}
Rf_*\Omega^\bullet_{X/S}
\ar[r]^-\mu &
Rf_*\Omega^\bullet_{X/S}[n]
}
$$
を $D(Y, q^{-1}\mathcal{O}_T)$ において得る。
\end{enumerate}
\end{remark}




\section{ホッジコホモロジー}
\label{derham-section-hodge-cohomology}

\noindent
$p : X \to S$ をスキームの射とする。{\it $X$ の、$S$ に関する
ホッジコホモロジー}を、コホモロジー群
$$
H^n_{Hodge}(X/S) = \bigoplus\nolimits_{n = p + q} H^q(X, \Omega^p_{X/S})
$$
によって定義し、これを次数付き $H^0(X, \mathcal{O}_X)$-加群とみなす。
微分形式の外積と Cohomology, Section \ref{cohomology-section-cup-product}
のカップ積を組み合わせると、$H^0(X, \mathcal{O}_X)$-双線形なカップ積
$$
\cup :
H^i_{Hodge}(X/S) \times H^j_{Hodge}(X/S)
\longrightarrow
H^{i + j}_{Hodge}(X/S)
$$
を得る。もちろん、$\xi \in H^q(X, \Omega^p_{X/S})$ および
$\xi' \in H^{q'}(X, \Omega^{p'}_{X/S})$ ならば $\xi \cup \xi' \in
H^{q + q'}(X, \Omega^{p + p'}_{X/S})$ である。

\begin{lemma}
\label{derham-lemma-cup-product-hodge-graded-commutative}
$p : X \to S$ をスキームの射とする。
$H^*_{Hodge}(X/S)$ 上のカップ積は結合的かつ次数付き可換である。
\end{lemma}

\begin{proof}
証明は Lemma \ref{derham-lemma-cup-product-graded-commutative} の証明と同一である。
\end{proof}

\noindent
スキームの可換図式
$$
\xymatrix{
X' \ar[r]_f \ar[d] & X \ar[d] \\
S' \ar[r] & S
}
$$
が与えられると、引き戻し写像
$f^* : H^i_{Hodge}(X/S) \longrightarrow H^i_{Hodge}(X'/S')$
があり、これは次数付けおよび上で定義したカップ積と両立する。







\section{二つのスペクトル系列}
\label{derham-section-hodge-to-de-rham}

\noindent
$p : X \to S$ をスキームの射とする。$p^{-1}\mathcal{O}_S$-加群からなる
$X$ 上の圏は十分多くの入射対象をもつので、
$\Omega^\bullet_{X/S}$ に対する Cartan--Eilenberg 分解が存在する。
Derived Categories, Lemma \ref{derived-lemma-cartan-eilenberg} を参照せよ。
したがって Derived Categories, Lemma
\ref{derived-lemma-two-ss-complex-functor} を適用すると、いずれも
$X$ の、$S$ に関するド・ラームコホモロジーに収束する二つのスペクトル系列を得る。

\medskip\noindent
第1のものは通常、{\it ホッジ・ド・ラームスペクトル系列}と呼ばれる。
このスペクトル系列の第1頁は
$$
E_1^{p, q} = H^q(X, \Omega^p_{X/S})
$$
であり、これは $X/S$ のホッジコホモロジー群である（名称の由来でもある）。
この頁の微分 $d_1$ は、微分
$d_1^{p, q} : H^q(X, \Omega^p_{X/S}) \to H^q(X, \Omega^{p + 1}_{X/S})$
によって与えられ、これは
$\text{d} : \Omega^p_{X/S} \to \Omega^{p + 1}_{X/S}$
から誘導される。図示すると次のようになる。
$$
\xymatrix{
H^2(X, \mathcal{O}_X) \ar[r] \ar@{-->}[rrd] \ar@{..>}[rrrdd] &
H^2(X, \Omega^1_{X/S}) \ar[r] \ar@{-->}[rrd] &
H^2(X, \Omega^2_{X/S}) \ar[r] &
H^2(X, \Omega^3_{X/S}) \\
H^1(X, \mathcal{O}_X) \ar[r] \ar@{-->}[rrd] &
H^1(X, \Omega^1_{X/S}) \ar[r] \ar@{-->}[rrd] &
H^1(X, \Omega^2_{X/S}) \ar[r] &
H^1(X, \Omega^3_{X/S}) \\
H^0(X, \mathcal{O}_X) \ar[r] &
H^0(X, \Omega^1_{X/S}) \ar[r] &
H^0(X, \Omega^2_{X/S}) \ar[r] &
H^0(X, \Omega^3_{X/S})
}
$$
ここで破線矢印は $E_2$ 頁の微分の始点と終点を示し、点線矢印は
$E_3$ 頁の微分を示している。次数 $0$ を見ると、
$$
H^0_{dR}(X/S) =
\Ker(\text{d} : H^0(X, \mathcal{O}_X) \to H^0(X, \Omega^1_{X/S}))
$$
を得る。もちろん、ド・ラーム複体が次数 $0$ において
$\mathcal{O}_X \to \Omega^1_{X/S}$ から始まることからも直ちに分かる。

\medskip\noindent
第2のスペクトル系列は通常、{\it 共役スペクトル系列}と呼ばれる。
このスペクトル系列の第2頁は
$$
E_2^{p, q} = H^p(X, H^q(\Omega^\bullet_{X/S})) = H^p(X, \mathcal{H}^q)
$$
である。ここで $\mathcal{H}^q = H^q(\Omega^\bullet_{X/S})$ は第 $q$
コホモロジー層であり、これは $X/S$ のド・ラーム複体から得られる。
この頁の微分は $E_2^{p, q} \to E_2^{p + 2, q - 1}$ によって与えられる。
図示すると次のようになる。
$$
\xymatrix{
H^0(X, \mathcal{H}^2) \ar[rrd] \ar@{..>}[rrrdd] &
H^1(X, \mathcal{H}^2) \ar[rrd] &
H^2(X, \mathcal{H}^2) &
H^3(X, \mathcal{H}^2) \\
H^0(X, \mathcal{H}^1) \ar[rrd] &
H^1(X, \mathcal{H}^1) \ar[rrd] &
H^2(X, \mathcal{H}^1) &
H^3(X, \mathcal{H}^1) \\
H^0(X, \mathcal{H}^0) &
H^1(X, \mathcal{H}^0) &
H^2(X, \mathcal{H}^0) &
H^3(X, \mathcal{H}^0)
}
$$
次数 $0$ を見ると、
$$
H^0_{dR}(X/S) = H^0(X, \mathcal{H}^0)
$$
を得る。これは定義を考えれば明らかである。次数 $1$ では完全列
$$
0 \to H^1(X, \mathcal{H}^0) \to H^1_{dR}(X/S) \to
H^0(X, \mathcal{H}^1) \to H^2(X, \mathcal{H}^0) \to H^2_{dR}(X/S)
$$
を得る。$X \to S$ が滑らかで $S$ の標数が $p$ ならば、層
$\mathcal{H}^q$ は（ある微分形式の層を用いて）計算でき、共役スペクトル系列は
有力な道具となる（ここに将来の参照を挿入する）。



\section{ホッジフィルトレーション}
\label{derham-section-hodge-filtration}

\noindent
$X \to S$ をスキームの射とする。$H^n_{dR}(X/S)$ 上のホッジフィルトレーションとは、
ホッジ・ド・ラームスペクトル系列によって誘導されるフィルトレーションである
(Homology, Definition
\ref{homology-definition-filtration-cohomology-filtered-complex})。
誤解を避けるため、次のように明示的に定義する。

\begin{definition}
\label{derham-definition-hodge-filtration}
$X \to S$ をスキームの射とする。$H^n_{dR}(X/S)$ 上の
{\it ホッジフィルトレーション}とは、各項が
$$
F^pH^n_{dR}(X/S) = \Im\left(H^n(X, \sigma_{\geq p}\Omega^\bullet_{X/S})
\longrightarrow H^n_{dR}(X/S)\right)
$$
で与えられるフィルトレーションである。ここで
$\sigma_{\geq p}\Omega^\bullet_{X/S}$ は Homology, Section
\ref{homology-section-truncations} におけるものとする。
\end{definition}

\noindent
もちろん $\sigma_{\geq p}\Omega^\bullet_{X/S}$ は相対ド・ラーム複体の
部分複体であり、相対ド・ラーム複体のフィルトレーション
$$
\Omega^\bullet_{X/S} = \sigma_{\geq 0}\Omega^\bullet_{X/S} \supset
\sigma_{\geq 1}\Omega^\bullet_{X/S} \supset
\sigma_{\geq 2}\Omega^\bullet_{X/S} \supset
\sigma_{\geq 3}\Omega^\bullet_{X/S} \supset \ldots
$$
を得る。また
$\text{gr}^p(\Omega^\bullet_{X/S}) = \Omega^p_{X/S}[-p]$ である。
Cohomology, Lemma \ref{cohomology-lemma-spectral-sequence-filtered-object}
において、$\Omega^\bullet_{X/S}$ をフィルター付き層の複体とみなして構成される
スペクトル系列は、Section \ref{derham-section-hodge-to-de-rham} で構成した
ホッジ・ド・ラームスペクトル系列と同じである。これは
Cohomology, Example \ref{cohomology-example-spectral-sequence-bis} による。
さらに、外積 (\ref{derham-equation-wedge}) は
$\text{Tot}(\sigma_{\geq i}\Omega^\bullet_{X/S} \otimes_{p^{-1}\mathcal{O}_S}
\sigma_{\geq j}\Omega^\bullet_{X/S})$ を
$\sigma_{\geq i + j}\Omega^\bullet_{X/S}$ に写す。したがって可換図式
$$
\xymatrix{
H^n(X, \sigma_{\geq i}\Omega^\bullet_{X/S}))
\times 
H^m(X, \sigma_{\geq j}\Omega^\bullet_{X/S}))
\ar[r] \ar[d] &
H^{n + m}(X, \sigma_{\geq i + j}\Omega^\bullet_{X/S})) \ar[d] \\
H^n_{dR}(X/S) \times
H^m_{dR}(X/S)
\ar[r]^\cup &
H^{n + m}_{dR}(X/S)
}
$$
を得る。特に、
$$
F^iH^n_{dR}(X/S) \cup F^jH^m_{dR}(X/S) \subset F^{i + j}H^{n + m}_{dR}(X/S)
$$
が成り立つ。







\section{K\"unneth 公式}
\label{derham-section-kunneth}

\noindent
ド・ラームコホモロジーの重要な特徴の一つは、K\"unneth 公式が存在することである。

\medskip\noindent
$a : X \to S$ と $b : Y \to S$ を、同じ終域をもつスキームの射とする。$p : X \times_S Y \to X$ と $q : X \times_S Y \to Y$ を射影射とし、$f = a \circ p = b \circ q$ とおく。図は次のとおりである。
$$
\xymatrix{
& X \times_S Y \ar[ld]^p \ar[rd]_q \ar[dd]^f \\
X \ar[rd]_a & & Y \ar[ld]^b \\
& S
}
$$
この節では、$\mathcal{O}_X$-加群 $\mathcal{F}$ と
$\mathcal{O}_Y$-加群 $\mathcal{G}$ を与え、次のようにおく。
$$
\mathcal{F} \boxtimes \mathcal{G} =
p^*\mathcal{F} \otimes_{\mathcal{O}_{X \times_S Y}} q^*\mathcal{G}
$$
準連接加群上の二変数関手
$(\mathcal{F}, \mathcal{G}) \mapsto \mathcal{F} \boxtimes \mathcal{G}$
は、$S$ 上有限位数の微分作用素を備えた準連接加群上の二変数関手へ拡張される。詳細は
Morphisms, Remark \ref{morphisms-remark-base-change-differential-operators} を参照されたい。
ド・ラーム複体 $\Omega^\bullet_{X/S}$ および
$\Omega^\bullet_{Y/S}$ の微分は、位数 $1$ の
$S$ 上の微分作用素であり、Modules, Lemma
\ref{modules-lemma-differentials-relative-de-rham-complex-order-1} により、
この事実が分かる。したがって、次の複体を考えることができる。
$$
\text{Tot}(\Omega^\bullet_{X/S} \boxtimes \Omega^\bullet_{Y/S})
$$
詳細は Derived Categories of Schemes, Section
\ref{perfect-section-kunneth-complexes} の議論を参照されたい。

\begin{lemma}
\label{derham-lemma-de-rham-complex-product}
上記の状況では、次の標準同型が存在する。
$$
\text{Tot}(\Omega^\bullet_{X/S} \boxtimes \Omega^\bullet_{Y/S})
\longrightarrow
\Omega^\bullet_{X \times_S Y/S}
$$
これは $f^{-1}\mathcal{O}_S$-加群の複体の同型である。
\end{lemma}

\begin{proof}
Morphisms, Lemma \ref{morphisms-lemma-differential-product} により、
$
\Omega_{X \times_S Y/S} = p^*\Omega_{X/S} \oplus q^*\Omega_{Y/S}
$
であることが分かる。外積を取ると、
$$
\Omega^n_{X \times_S Y/S} =
\bigoplus\nolimits_{i + j = n}
p^*\Omega^i_{X/S} \otimes_{\mathcal{O}_{X \times_S Y}} q^*\Omega^j_{Y/S} =
\bigoplus\nolimits_{i + j = n}
\Omega^i_{X/S} \boxtimes \Omega^j_{Y/S}
$$
を得る。これは外積の基本的性質による。これらの同一視により、補題の矢印の左辺と右辺にある複体の各項の間の同型が定まる。これらの写像が微分と両立することの確認は省略する。
\end{proof}

\noindent
$A = \Gamma(S, \mathcal{O}_S)$ とおく。Lemma
\ref{derham-lemma-de-rham-complex-product} の結果を、Derived Categories of Schemes, Equation
(\ref{perfect-equation-de-rham-kunneth}) の写像と組み合わせると、次のカップ積を得る。
$$
R\Gamma(X, \Omega^\bullet_{X/S})
\otimes_A^\mathbf{L}
R\Gamma(Y, \Omega^\bullet_{Y/S})
\longrightarrow
R\Gamma(X \times_S Y, \Omega^\bullet_{X \times_S Y/S})
$$
コホモロジー群のレベルでは、More on Algebra, Section \ref{more-algebra-section-products-tor} の議論を用いると、標準写像
$$
H^i_{dR}(X/S) \otimes_A H^j_{dR}(Y/S)
\longrightarrow
H^{i + j}_{dR}(X \times_S Y/S),\quad
(\xi, \zeta) \longmapsto p^*\xi \cup q^*\zeta
$$
を得る。この構成は、まず引き戻しを行い、その後カップ積を取ることで進むことに注意する。

\begin{lemma}
\label{derham-lemma-kunneth-de-rham}
$X$ と $Y$ が滑らかかつ準コンパクトで、対角射がアフィンであるとし、$S = \Spec(A)$ 上で考える。すると、写像
$$
R\Gamma(X, \Omega^\bullet_{X/S})
\otimes_A^\mathbf{L}
R\Gamma(Y, \Omega^\bullet_{Y/S})
\longrightarrow
R\Gamma(X \times_S Y, \Omega^\bullet_{X \times_S Y/S})
$$
は $D(A)$ における同型である。
\end{lemma}

\begin{proof}
Morphisms, Lemma \ref{morphisms-lemma-smooth-omega-finite-locally-free} により、層
$\Omega^n_{X/S}$ と $\Omega^m_{Y/S}$ は、それぞれ $\mathcal{O}_X$ および
$\mathcal{O}_Y$ 上の有限局所自由加群である。一方、$X$ と $Y$ は
$S$ 上平坦である（Morphisms, Lemma \ref{morphisms-lemma-smooth-flat}）ので、
$\Omega^n_{X/S}$ と $\Omega^m_{Y/S}$ は $S$ 上平坦であることが分かる。
さらに、$\Omega^\bullet_{X/S}$ は局所有界複体であることに注意する。したがって、
Lemma \ref{derham-lemma-de-rham-complex-product} および
Derived Categories of Schemes, Lemma \ref{perfect-lemma-kunneth-special} により主張が従う。
\end{proof}

\noindent
カップ積には相対版、すなわち次の写像がある。
$$
Ra_*\Omega^\bullet_{X/S}
\otimes_{\mathcal{O}_S}^\mathbf{L}
Rb_*\Omega^\bullet_{Y/S}
\longrightarrow
Rf_*\Omega^\bullet_{X \times_S Y/S}
$$
これは $D(\mathcal{O}_S)$ における写像である。この構成は、
Lemma \ref{derham-lemma-de-rham-complex-product} と Derived Categories of Schemes, Equation
(\ref{perfect-equation-relative-de-rham-kunneth}) の写像を組み合わせたものである。
この構成から、この写像は次の図式で与えられる。
$$
\xymatrix{
Ra_*\Omega^\bullet_{X/S}
\otimes_{\mathcal{O}_S}^\mathbf{L}
Rb_*\Omega^\bullet_{Y/S}
\ar[d]^{\text{units of adjunction}}  \\
Rf_*(p^{-1}\Omega^\bullet_{X/S})
\otimes_{\mathcal{O}_S}^\mathbf{L}
Rf_*(q^{-1}\Omega^\bullet_{Y/S}) \ar[r] \ar[d]^{\text{relative cup product}} &
Rf_*(\Omega^\bullet_{X \times_S Y/S})
\otimes_{\mathcal{O}_S}^\mathbf{L}
Rf_*(\Omega^\bullet_{X \times_S Y/S}) \ar[d]^{\text{relative cup product}} \\
Rf_*(p^{-1}\Omega^\bullet_{X/S}
\otimes_{f^{-1}\mathcal{O}_S}^\mathbf{L}
q^{-1}\Omega^\bullet_{Y/S})
\ar[d]^{\text{from derived to usual}} \ar[r] &
Rf_*(\Omega^\bullet_{X \times_S Y/S}
\otimes_{f^{-1}\mathcal{O}_S}^\mathbf{L}
\Omega^\bullet_{X \times_S Y/S})
\ar[d]^{\text{from derived to usual}} \\
Rf_*\text{Tot}(p^{-1}\Omega^\bullet_{X/S}
\otimes_{f^{-1}\mathcal{O}_S}
q^{-1}\Omega^\bullet_{Y/S}) \ar[r] \ar[d]^{\text{canonical map}} &
Rf_*\text{Tot}(\Omega^\bullet_{X \times_S Y/S}
\otimes_{f^{-1}\mathcal{O}_S}
\Omega^\bullet_{X \times_S Y/S})
\ar[d]^{\eta \otimes \omega \mapsto \eta \wedge \omega} \\
Rf_*\text{Tot}(\Omega^\bullet_{X/S} \boxtimes \Omega^\bullet_{Y/S})
\ar@{=}[r]
&
Rf_*\Omega^\bullet_{X \times_S Y/S}
}
$$
ここで最初の矢印は、随伴の単位 $\text{id} \to Rp_* p^{-1}$
および $\text{id} \to Rq_* q^{-1}$、さらに同一視
$Rf_* p^{-1} = Ra_* Rp_* p^{-1}$ および
$Rf_* q^{-1} = Rb_* Rq_* q^{-1}$ を用いる。
第二の矢印は、Cohomology, Remark \ref{cohomology-remark-cup-product} の相対カップ積である。
第三の矢印は、複体のテンソル積の全体化への導来テンソル積の写像である。
最後の等号は Lemma \ref{derham-lemma-de-rham-complex-product} である。
この構成は、全体切断上で先に与えた構成を再現する。

\begin{lemma}
\label{derham-lemma-kunneth-de-rham-relative}
$X \to S$ と $Y \to S$ が滑らかかつ準コンパクトで、
射 $X \to X \times_S X$ および $Y \to Y \times_S Y$ がアフィンであると仮定する。
すると、相対カップ積
$$
Ra_*\Omega^\bullet_{X/S}
\otimes_{\mathcal{O}_S}^\mathbf{L}
Rb_*\Omega^\bullet_{Y/S}
\longrightarrow
Rf_*\Omega^\bullet_{X \times_S Y/S}
$$
は $D(\mathcal{O}_S)$ における同型である。
\end{lemma}

\begin{proof}
Lemma \ref{derham-lemma-kunneth-de-rham} の直ちに得られる帰結である。
\end{proof}








\section{ド・ラームコホモロジーにおける第一Chern類}
\label{derham-section-first-chern-class}

\noindent
スキームの射 $X \to S$ をとる。複体の写像
$$
\text{d}\log : \mathcal{O}_X^*[-1] \longrightarrow \Omega^\bullet_{X/S}
$$
が存在し、これは切断 $g \in \mathcal{O}_X^*(U)$ を、
$\text{d}\log(g) = g^{-1}\text{d}g$ で与えられる
$\Omega^1_{X/S}(U)$ の切断へ送る。
したがって写像
$$
\Pic(X) = H^1(X, \mathcal{O}_X^*) =
H^2(X, \mathcal{O}_X^*[-1]) \longrightarrow H^2_{dR}(X/S)
$$
を考えることができる。ここで最初の等号は
Cohomology, Lemma \ref{cohomology-lemma-h1-invertible} による。
可逆加群 $\mathcal{L}$ の同型類の像を
$c^{dR}_1(\mathcal{L}) \in H^2_{dR}(X/S)$ と記す。

\medskip\noindent
さらに写像 $\text{d}\log : \mathcal{O}_X^* \to \Omega^1_{X/S}$
を用いて、HodgeコホモロジーにおけるChern類を
$$
c_1^{Hodge} : \Pic(X) \longrightarrow H^1(X, \Omega^1_{X/S})
\subset H^2_{Hodge}(X/S)
$$
によって定義することもできる。
これらの構成は引き戻しと両立する。

\begin{lemma}
\label{derham-lemma-pullback-c1}
次の可換図式
$$
\xymatrix{
X' \ar[r]_f \ar[d] & X \ar[d] \\
S' \ar[r] & S
}
$$
がスキームの図式として与えられると、次の図式は可換である。
$$
\xymatrix{
\Pic(X') \ar[d]_{c_1^{dR}} &
\Pic(X) \ar[d]^{c_1^{dR}} \ar[l]^{f^*} \\
H^2_{dR}(X'/S') &
H^2_{dR}(X/S) \ar[l]_{f^*}
}
\quad
\xymatrix{
\Pic(X') \ar[d]_{c_1^{Hodge}} &
\Pic(X) \ar[d]^{c_1^{Hodge}} \ar[l]^{f^*} \\
H^1(X', \Omega^1_{X'/S'}) &
H^1(X, \Omega^1_{X/S}) \ar[l]_{f^*}
}
$$
\end{lemma}

\begin{proof}
省略する。
\end{proof}

\noindent
要素 $c^{dR}_1(\mathcal{L})$ を{\v C}echコホモロジーで「計算」する（{\v C}ech微分の符号規則は
Cohomology, Section
\ref{cohomology-section-cech-cohomology-of-complexes} のものと同じとする）。
すなわち、開被覆
$\mathcal{U} : X = \bigcup_{i \in I} U_i$ を選ぶ。自明化切断 $s_i$ を選び、これは
$\mathcal{L}|_{U_i}$ の切断である。これをすべての $i$ について行う。
重なり $U_{i_0i_1} = U_{i_0} \cap U_{i_1}$ 上には、可逆関数 $f_{i_0i_1}$ が存在し、
$f_{i_0i_1} = s_{i_1}|_{U_{i_0i_1}} s_{i_0}|_{U_{i_0i_1}}^{-1}$ となる
\footnote{{\v C}echの $0$-サイクル $\{a_{i_0}\}$ の微分は、
$a_{i_1} - a_{i_0}$ であり、$U_{i_0i_1}$ 上で定義される。}。
もちろん
$$
f_{i_1i_2}|_{U_{i_0i_1i_2}}
f_{i_0i_2}^{-1}|_{U_{i_0i_1i_2}}
f_{i_0i_1}|_{U_{i_0i_1i_2}} = 1
$$
が成り立つ。
$\mathcal{L}$ の $H^1(X, \mathcal{O}_X^*)$ におけるコホモロジー類は、
コサイクル $\{f_{i_0i_1}\}$ の{\v C}echコホモロジー類であり、
$\check{\mathcal{C}}^\bullet(\mathcal{U}, \mathcal{O}_X^*)$ におけるその像である。
したがって $c_1^{dR}(\mathcal{L})$ は、{\v C}echコサイクル
$\{\alpha_{i_0 \ldots i_p}\}$ に付随するコホモロジー類の像であり、
$\text{Tot}(\check{\mathcal{C}}^\bullet(\mathcal{U}, \Omega_{X/S}^\bullet))$
における次数 $2$ の類である。これは次で与えられる。
\begin{enumerate}
\item $\alpha_{i_0} = 0$ は $\Omega^2_{X/S}(U_{i_0})$ において成り立つ。
\item $\alpha_{i_0i_1} = f_{i_0i_1}^{-1}\text{d}f_{i_0i_1}$ は
$\Omega^1_{X/S}(U_{i_0i_1})$ において成り立つ。
\item $\alpha_{i_0i_1i_2} = 0$ は $\mathcal{O}_{X/S}(U_{i_0i_1i_2})$ において成り立つ。
\end{enumerate}
可逆加群 $\mathcal{L}_k$ を、$k = 1, \ldots, a$ としてとる。
これらを $U_i$ 上で自明化し、すべての $i \in I$ について、上記のように
コサイクル $f_{k, i_0i_1}$ と $\alpha_k = \{\alpha_{k, i_0 \ldots i_p}\}$ が得られる。
Cohomology, Section
\ref{cohomology-section-cech-cohomology-of-complexes} の規則を用いると、
$$
\beta = \alpha_1 \cup \alpha_2 \cup \ldots \cup \alpha_a
$$
を、以下のように記述されるコサイクル $\beta = \{\beta_{i_0 \ldots i_p}\}$ として計算できる。
\begin{enumerate}
\item $\beta_{i_0 \ldots i_p} = 0$ は
$\Omega^{2a - p}_{X/S}(U_{i_0 \ldots i_p})$ において成り立つ（ただし $p = a$ の場合を除く）。
\item $\beta_{i_0 \ldots i_a} = (-1)^{a(a - 1)/2}
\alpha_{1, i_0i_1} \wedge \alpha_{2, i_1 i_2} \wedge \ldots \wedge
\alpha_{a, i_{a - 1}i_a}$ は
$\Omega^a_{X/S}(U_{i_0 \ldots i_a})$ において成り立つ。
\end{enumerate}
したがって、これは
$c_1^{dR}(\mathcal{L}_1) \cup \ldots \cup c_1^{dR}(\mathcal{L}_a)$
を表すコサイクルである。
もちろん、同じ計算により、コサイクル
$\{\beta_{i_0 \ldots i_a}\}$ は
$\check{\mathcal{C}}^a(\mathcal{U}, \Omega_{X/S}^a))$ においてコホモロジー類
$c_1^{Hodge}(\mathcal{L}_1) \cup \ldots \cup c_1^{Hodge}(\mathcal{L}_a)$
を表す。

\begin{remark}
\label{derham-remark-truncations}
ここまでの計算を、より抽象的な言葉で言い換える。
スキームの射 $p : X \to S$ をとる。$\mathcal{L}$ を可逆
$\mathcal{O}_X$-加群とする。$\text{d}\log$ を写像
$$
\mathcal{O}_X^*[-1] \to \sigma_{\geq 1}\Omega^\bullet_{X/S}
$$
とみなすと、上と同様に $\Pic(X) = H^1(X, \mathcal{O}_X^*)$ を用いて、
コホモロジー類
$$
\gamma_1(\mathcal{L}) \in H^2(X, \sigma_{\geq 1}\Omega^\bullet_{X/S})
$$
を得る。$\gamma_1(\mathcal{L})$ の像は、写像
$\sigma_{\geq 1}\Omega^\bullet_{X/S} \to \Omega^\bullet_{X/S}$
により $c_1^{dR}(\mathcal{L})$ を回復する。
特に
$c_1^{dR}(\mathcal{L}) \in F^1H^2_{dR}(X/S)$ であることが分かる。Section
\ref{derham-section-hodge-filtration} を参照されたい。$\gamma_1(\mathcal{L})$ の像は、写像
$\sigma_{\geq 1}\Omega^\bullet_{X/S} \to \Omega^1_{X/S}[-1]$
により $c_1^{Hodge}(\mathcal{L})$ を回復する。
カップ積（Section \ref{derham-section-hodge-filtration} を参照）を取ると、
$$
\xi = \gamma_1(\mathcal{L}_1) \cup \ldots \cup \gamma_1(\mathcal{L}_a) \in
H^{2a}(X, \sigma_{\geq a}\Omega^\bullet_{X/S})
$$
を得る。Section \ref{derham-section-hodge-filtration} の可換図式から、
$\xi$ は
$c_1^{dR}(\mathcal{L}_1) \cup \ldots \cup c_1^{dR}(\mathcal{L}_a)$
へ $H^{2a}_{dR}(X/S)$ において写されることが分かる。これは写像
$\sigma_{\geq a}\Omega^\bullet_{X/S} \to \Omega^\bullet_{X/S}$ による。
また、
$c_1^{dR}(\mathcal{L}_1) \cup \ldots \cup c_1^{dR}(\mathcal{L}_a)$
は $F^a H^{2a}_{dR}(X/S)$ に含まれることが従う。同様に、
写像 $\sigma_{\geq a}\Omega^\bullet_{X/S} \to \Omega^a_{X/S}[-a]$
は $\xi$ を
$c_1^{Hodge}(\mathcal{L}_1) \cup \ldots \cup c_1^{Hodge}(\mathcal{L}_a)$
へ $H^a(X, \Omega^a_{X/S})$ において送る。
\end{remark}

\begin{remark}
\label{derham-remark-log-forms}
スキームの射 $p : X \to S$ をとる。$i > 0$ に対し
$\Omega^i_{X/S, log} \subset \Omega^i_{X/S}$ を、局所切断
$$
\text{d}\log(u_1) \wedge \ldots \wedge \text{d}\log(u_i)
$$
で生成されるアーベル部分層として定める。ただし $u_1, \ldots, u_n$ は
$\mathcal{O}_X$ の可逆局所切断である。
$i = 0$ のとき、部分層 $\Omega^0_{X/S, log} \subset \mathcal{O}_X$ は
$\mathbf{Z} \to \mathcal{O}_X$ の像である。すべての $i \geq 0$ に対して複体の写像
$$
\Omega^i_{X/S, log}[-i] \longrightarrow \Omega^\bullet_{X/S}
$$
が存在する。これは対数形式の微分が零であるためである。さらに、対数形式の外積も
再び対数形式となるので、次の双線形写像を得る。
$$
\wedge :  \Omega^i_{X/S, log} \times
\Omega^j_{X/S, log} \longrightarrow \Omega^{i + j}_{X/S, log}
$$
これは (\ref{derham-equation-wedge}) および上の写像と両立する。
$\mathcal{L}$ を可逆 $\mathcal{O}_X$-加群とする。アーベル層の写像
$\text{d}\log : \mathcal{O}_X^* \to \Omega^1_{X/S, log}$
と同一視 $\Pic(X) = H^1(X, \mathcal{O}_X^*)$ を用いると、標準的なコホモロジー類
$$
\tilde \gamma_1(\mathcal{L}) \in H^1(X, \Omega^1_{X/S, log})
$$
を得る。これらの類は次の性質をもつ。
\begin{enumerate}
\item $\tilde \gamma_1(\mathcal{L})$ について、標準写像
$\Omega^1_{X/S, log}[-1] \to \sigma_{\geq 1}\Omega^\bullet_{X/S}$
は $\tilde \gamma_1(\mathcal{L})$ を Remark
\ref{derham-remark-truncations} の類
$\gamma_1(\mathcal{L}) \in 
H^2(X, \sigma_{\geq 1}\Omega^\bullet_{X/S})$ へ送る。
\item $\tilde \gamma_1(\mathcal{L})$ について、標準写像
$\Omega^1_{X/S, log}[-1] \to \Omega^\bullet_{X/S}$
は $\tilde \gamma_1(\mathcal{L})$ を
$c_1^{dR}(\mathcal{L})$ として $H^2_{dR}(X/S)$ へ送る。
\item $\tilde \gamma_1(\mathcal{L})$ について、標準写像
$\Omega^1_{X/S, log} \to \Omega^1_{X/S}$
は $\tilde \gamma_1(\mathcal{L})$ を
$c_1^{Hodge}(\mathcal{L})$ として $H^1(X, \Omega^1_{X/S})$ へ送る。
\item これらの類の構成は引き戻しと両立する。
\item ここにさらに追加する。
\end{enumerate}
\end{remark}









\section{線束のド・ラームコホモロジー}
\label{derham-section-line-bundle}

\noindent
線束はベクトル束の特別な場合であり、ベクトル束はさらにいくつかの
付加構造を備えた錐である。これらのド・ラーム複体を適切に扱うには、
次数付き環のド・ラーム複体を論じるのが自然である。

\begin{remark}[次数付き環のド・ラーム複体]
\label{derham-remark-de-rham-complex-graded}
$G$ を、加法的に書かれ、単位元 $0$ をもつ可換モノイドとする。
$R \to A$ を環写像とし、$A$ が次数付けを備えていると仮定する。
$A = \bigoplus_{g \in G} A_g$ は $R$-加群による次数付けであり、$R$ の像は $A_0$ に入り、
$A_g \cdot A_{g'} \subset A_{g + g'}$ である。このとき微分加群も
次数付けを備える。
$$
\Omega_{A/R} = \bigoplus\nolimits_{g \in G} \Omega_{A/R, g}
$$
ここで $\Omega_{A/R, g}$ は、$R$-加群 $\Omega_{A/R}$ の部分加群であり、
$a_0 \text{d}a_1$（ただし $a_i \in A_{g_i}$）によって生成され、
$g = g_0 + g_1$ を満たす。同様に次を得る。
$$
\Omega^p_{A/R} = \bigoplus\nolimits_{g \in G} \Omega^p_{A/R, g}
$$
ここで $\Omega^p_{A/R, g}$ は、$R$-加群 $\Omega^p_{A/R}$ の部分加群であり、
$a_0 \text{d}a_1 \wedge \ldots \wedge \text{d}a_p$
（ただし $a_i \in A_{g_i}$）によって生成され、$g = g_0 + g_1 + \ldots + g_p$ を満たす。
もちろん微分は次数付けを保ち、外積は明らかな仕方で
次数付けと両立する。
\end{remark}

\noindent
$f : X \to S$ を概形の射とし、$\pi : C \to X$ を錐とする。Constructions,
Definition \ref{constructions-definition-abstract-cone} を参照されたい。
これは、$\pi$ がアフィンであり、次数付け
$\pi_*\mathcal{O}_C = \bigoplus_{n \geq 0} \mathcal{A}_n$ を備え、
$\mathcal{A}_0 = \mathcal{O}_X$ であることを意味する。
Remark \ref{derham-remark-de-rham-complex-graded} の議論を
アフィン開集合上で用いると、\footnote{この記法についてはご容赦いただきたい。}
$$
\pi_*(\Omega^\bullet_{C/S}) =
\bigoplus\nolimits_{n \geq 0} \Omega^\bullet_{C/S, n}
$$
は標準的に部分複体の直和となる。さらに因数分解
$$
\Omega^\bullet_{X/S} \to \Omega^\bullet_{C/S, 0} \to
\pi_*(\Omega^\bullet_{C/S})
$$
があり、包含 $\omega \wedge \eta \in \Omega^{p + q}_{C/S, n + m}$ は、
$\omega \in \Omega^p_{C/S, n}$ かつ $\eta \in \Omega^q_{C/S, m}$ のとき成り立つ。

\medskip\noindent
$f : X \to S$ を概形の射とし、$\pi : L \to X$ を
可逆 $\mathcal{O}_X$-加群 $\mathcal{L}$ に付随する線束とする。
これは、$\pi$ が次を満たす唯一のアフィン射であることを意味する。
$$
\pi_*\mathcal{O}_L = 
\bigoplus\nolimits_{n \geq 0} \mathcal{L}^{\otimes n}
$$
等式は $\mathcal{O}_X$-代数として成り立つ。したがって $L$ は $X$ 上の錐である。
上の議論から、標準的な直和分解
$$
\pi_*(\Omega^\bullet_{L/S}) =
\bigoplus\nolimits_{n \geq 0} \Omega^\bullet_{L/S, n}
$$
を得る。これは外積および上記の分解と両立し、
$\pi_*\mathcal{O}_L$ の分解とも両立する。また $\Omega_{X/S}$ は
部分 $\Omega_{L/S, 0}$（次数 $0$）へ写る。

\medskip\noindent
もう一つ、後で有用となる場合がある。すなわち、
補集合\footnote{概形 $L^\star$ は、$\mathbf{G}_m$-トーサーであり、
$X$ 上で $L$ に付随する。このため以下で得られる次数付けは
$\mathbf{Z}$-次数付けとなる。Groupoids,
Example \ref{groupoids-example-Gm-on-affine} および
Lemmas \ref{groupoids-lemma-complete-reducibility-Gm} と
\ref{groupoids-lemma-Gm-equivariant-module} と比較されたい。}
$L^\star \subset L$ を、零切断 $o : X \to L$ の補集合として線束 $L$ の中で考える。
局所計算により、標準同型
$$
(L^\star \to X)_*\mathcal{O}_{L^\star} =
\bigoplus\nolimits_{n \in \mathbf{Z}} \mathcal{L}^{\otimes n}
$$
が $\mathcal{O}_X$-代数として得られる。右辺は $\mathbf{Z}$-次数付きの
準連接 $\mathcal{O}_X$-代数である。Remark \ref{derham-remark-de-rham-complex-graded}
の議論をアフィン開集合上で用いると
$$
(L^\star \to X)_*(\Omega^\bullet_{L^\star/S}) =
\bigoplus\nolimits_{n \in \mathbf{Z}} \Omega^\bullet_{L^\star/S, n}
$$
を得る。これは外積および上記の分解と両立し、
$(L^\star \to X)_*\mathcal{O}_{L^\star}$ の分解とも両立する。また
$\Omega_{X/S}$ は部分 $\Omega_{L^\star/S, 0}$（次数 $0$）へ写る。
複体 $\Omega^\bullet_{L^\star/S, 0}$ は特に重要となる。

\begin{lemma}
\label{derham-lemma-the-complex-for-L-star}
上の記法のもとで、複体の短完全列
$$
0 \to \Omega^\bullet_{X/S} \to
\Omega^\bullet_{L^\star/S, 0} \to
\Omega^\bullet_{X/S}[-1] \to 0
$$
が存在する。
\end{lemma}

\begin{proof}
上で写像
$\Omega^\bullet_{X/S} \to \Omega^\bullet_{L^\star/S, 0}$ を構成した。

\medskip\noindent
次の写像を構成する。
$\text{Res} : \Omega^\bullet_{L^\star/S, 0} \to \Omega^\bullet_{X/S}[-1]$。
$U \subset X$ を開集合とし、$s \in \mathcal{L}(U)$ をとる。
さらに $s' \in \mathcal{L}^{\otimes -1}(U)$ を、
$s' s = 1$ を満たす切断とする。このとき $s$ は代数の層
$(L^\star \to X)_*\mathcal{O}_{L^\star}$ の $U$ 上の可逆切断であり、
逆元は $s' = s^{-1}$ である。そこで $1$-形式
$\text{d}\log(s) = s' \text{d}(s)$ を考える。これは、
$\Omega^1_{L^\star/S, 0}(U)$ の元である。その理由は、
$\Omega^1_{L^\star/S}$ の次数付けの構成にある。以下のアフィン上の計算は、
$1$ と $\text{d}\log(s)$ が
$\Omega^\bullet_{L^\star/S, 0}|_U$ を右加群として自由に生成し、その係数環が
$\Omega^\bullet_{X/S}|_U$ であることを示す。
したがって、$\text{Res}$ を $U$ 上で次の規則により定義できる。
$$
\text{Res}(\omega' + \text{d}\log(s) \wedge \omega) = \omega
$$
これはすべての $\omega', \omega \in \Omega^\bullet_{X/S}(U)$ に対して定める。この
写像は局所生成元 $s$ の選択に依存せず、したがって
貼り合わさって大域写像を与える。実際、別の $s$ は
$gs$ の形であり、ここで $g \in \mathcal{O}_X(U)$ は可逆である。このとき
$\text{d}\log(gs) = g^{-1}\text{d}(g) + \text{d}\log(s)$
となるので、独立性は容易に従う。
最後に、$\text{Res}$ の規則が
微分と両立することに注意する。実際、
$\text{d}(\omega' + \text{d}\log(s) \wedge \omega) =
\text{d}(\omega') - \text{d}\log(s) \wedge \text{d}(\omega)$
であり、$\Omega^\bullet_{X/S}[-1]$ 上の微分は
$\omega'$ を $-\text{d}(\omega')$ に写す。これは
Homology, Definition \ref{homology-definition-shift-cochain} の符号規約による。

\medskip\noindent
局所計算。$X$ をアフィン開集合 $U \subset X$ で覆える。
これらは $\mathcal{L}|_U \cong \mathcal{O}_U$ を満たし、さらに
アフィン開集合 $V \subset S$ に写る。$U = \Spec(A)$、$V = \Spec(R)$ と書き、
$s$ を $\mathcal{L}$ の生成元として選ぶ。このとき
$$
L^\star \times_X U = \Spec(A[s, s^{-1}])
$$
微分を計算すると
$$
\Omega^1_{A[s, s^{-1}]/R} =
A[s, s^{-1}] \otimes_A \Omega^1_{A/R} \oplus A[s, s^{-1}] \text{d}\log(s)
$$
であり、したがって外積をとると
$$
\Omega^p_{A[s, s^{-1}]/R} =
A[s, s^{-1}] \otimes_A \Omega^p_{A/R}
\oplus
A[s, s^{-1}] \text{d}\log(s) \otimes_A \Omega^{p - 1}_{A/R}
$$
次数 $0$ 部分をとると
$$
\Omega^p_{A[s, s^{-1}]/R, 0} =
\Omega^p_{A/R} \oplus \text{d}\log(s) \otimes_A \Omega^{p - 1}_{A/R}
$$
となり、補題の証明が完了する。
\end{proof}

\begin{lemma}
\label{derham-lemma-the-complex-for-L-star-gives-chern-class}
「境界」写像
$\delta : \Omega^\bullet_{X/S} \to \Omega^\bullet_{X/S}[2]$
は、$D(X, f^{-1}\mathcal{O}_S)$ において
補題 \ref{derham-lemma-the-complex-for-L-star} の短完全列から生じる。
これは Remark \ref{derham-remark-cup-product-as-a-map} の写像であり、
$\xi = c_1^{dR}(\mathcal{L})$ に対応する。
\end{lemma}

\begin{proof}
正確には、次のシフトを考える。
$$
0 \to \Omega^\bullet_{X/S}[1] \to
\Omega^\bullet_{L^\star/S, 0}[1] \to
\Omega^\bullet_{X/S} \to 0
$$
これは補題 \ref{derham-lemma-the-complex-for-L-star} の短完全列をシフトしたものである。
次数付けを備えた
$(L^\star \to X)_*\Omega^\bullet_{L^\star/S}$
の次数零部分として、$\Omega^\bullet_{L^\star/S, 0}$ は微分次数付き
$\mathcal{O}_X$-代数であり、写像
$\Omega^\bullet_{X/S} \to \Omega^\bullet_{L^\star/S, 0}$
は微分次数付き $\mathcal{O}_X$-代数の準同型である。
したがって、$\Omega^\bullet_{X/S}[1] \to
\Omega^\bullet_{L^\star/S, 0}[1]$ を右微分次数付き
$\Omega^\bullet_{X/S}$-加群の $X$ 上の写像とみなせる。写像
$\text{Res} : \Omega^\bullet_{L^\star/S, 0}[1] \to \Omega^\bullet_{X/S}$
は右微分次数付き $\Omega^\bullet_{X/S}$-加群の写像である。
実際、局所的には規則
$\text{Res}(\omega' + \text{d}\log(s) \wedge \omega) = \omega$ で定義される。補題
\ref{derham-lemma-the-complex-for-L-star} の証明を参照されたい。
したがって、
Differential Graded Sheaves, Section \ref{sdga-section-misc}
の議論により、$\delta$ は写像
$\delta' : \Omega^\bullet_{X/S} \to \Omega^\bullet_{X/S}[2]$
から生じる。この写像は、ド・ラーム複体上の右微分次数付き加群の導来圏
$D(\Omega^\bullet_{X/S}, \text{d})$ に属する。
Remark \ref{derham-remark-cup-product-as-a-map} で述べた一意性により、
$\delta(1) = c_1^{dR}(\mathcal{L})$ を示せば十分である。

\medskip\noindent
次の可換図式があることを主張する。
$$
\xymatrix{
0 \ar[r] &
\mathcal{O}_X^* \ar[r] \ar[d]_{\text{d}\log} &
E \ar[r] \ar[d] &
\underline{\mathbf{Z}} \ar[d] \ar[r] &
0 \\
0 \ar[r] &
\Omega^\bullet_{X/S}[1] \ar[r] &
\Omega^\bullet_{L^\star/S, 0}[1] \ar[r] &
\Omega^\bullet_{X/S} \ar[r] &
0
}
$$
上段はアーベル層の短完全列であり、その境界写像は
$1$ を $\mathcal{L}$ の類へ写す。この類は
$H^1(X, \mathcal{O}_X^*)$ に属する。短完全列間の写像に関する
境界写像の両立性により、この主張を示せば十分である。
$E$ を次の規則の層化として定義する。
$$
U \longmapsto \{(s, n) \mid
n \in \mathbf{Z},\ s \in \mathcal{L}^{\otimes n}(U)\text{ generator}\}
$$
群構造は $(s, n) \cdot (t, m) = (s \otimes t, n + m)$ で与える。
中央の垂直写像は $(s, n)$ を $\text{d}\log(s)$ へ写す。これにより
短完全列の写像が得られる。なぜなら、補題
\ref{derham-lemma-the-complex-for-L-star} の証明で構成した写像
$Res : \Omega^1_{L^\star/S, 0} \to \mathcal{O}_X$ は
$\text{d}\log(s)$ を $1$ へ写すからである。ただし $s$ は $\mathcal{L}$ の局所生成元とする。
上段の $1$ の境界を計算するため、局所自明化
$s_i$ を $\mathcal{L}$ の開集合 $U_i$ 上で選ぶ。これは
Section \ref{derham-section-first-chern-class} と同様である。重なり
$U_{i_0i_1} = U_{i_0} \cap U_{i_1}$
上で、可逆関数 $f_{i_0i_1}$ を
$f_{i_0i_1} = s_{i_1}|_{U_{i_0i_1}} s_{i_0}|_{U_{i_0i_1}}^{-1}$
により定める。$\mathcal{L}$ のコホモロジー類は {\v C}ech コサイクル
$\{f_{i_0i_1}\}$ によって与えられる。そして
$$
(f_{i_0i_1}, 0) = (s_{i_1}, 1)|_{U_{i_0i_1}} \cdot
(s_{i_0}, 1)|_{U_{i_0i_1}}^{-1}
$$
が $E$ の切断として成り立ち、証明が完了する。
\end{proof}

\begin{lemma}
\label{derham-lemma-push-omega-a}
上の記法のもとで次が成り立つ。
\begin{enumerate}
\item $\Omega^p_{L^\star/S, n} =
\Omega^p_{L^\star/S, 0} \otimes_{\mathcal{O}_X} \mathcal{L}^{\otimes n}$
これはすべての $n \in \mathbf{Z}$ に対して、準連接 $\mathcal{O}_X$-加群として成り立つ。
\item $\Omega^\bullet_{X/S} = \Omega^\bullet_{L/X, 0}$
これは複体として成り立つ。さらに、
\item $n > 0$ かつ $p \geq 0$ のとき
$\Omega^p_{L/X, n} = \Omega^p_{L^\star/S, n}$。
\end{enumerate}
\end{lemma}

\begin{proof}
いずれの場合も大域的に定義された標準写像があり、
省略する局所計算により同型である。
\end{proof}

\begin{lemma}
\label{derham-lemma-line-bundle-characteristic-zero}
上の状況で、射 $S \to \Spec(\mathbf{Q})$ があると仮定する。
このとき $\Omega^\bullet_{X/S} \to \pi_*\Omega^\bullet_{L/S}$ は
擬同型であり、$H_{dR}^*(X/S) = H_{dR}^*(L/S)$ である。
\end{lemma}

\begin{proof}
$R$ を $\mathbf{Q}$-代数とし、$A$ を $R$-代数とする。
アフィン局所的な主張は、写像
$$
\Omega^\bullet_{A/R} \longrightarrow \Omega^\bullet_{A[t]/R}
$$
が $R$-加群の複体の擬同型であるということである。実際、これは
ホモトピー同値であり、そのホモトピー逆は写像
$g \omega + g' \text{d}t \wedge \omega'$ を $g(0)\omega$ へ写すものとして与えられる。ただし
$g, g' \in A[t]$ かつ $\omega, \omega' \in \Omega^\bullet_{A/R}$ である。
ホモトピーは $g \omega + g' \text{d}t \wedge \omega'$ を
$(\int g') \omega'$ へ写す。ここで $\int g' \in A[t]$ は、定数項が消え、
$t$ に関する導関数が
$g'$ となる多項式である。ここでは $R$ が $\mathbf{Q}$ を含むことを用いる。実際、
$\int t^n = (1/n)t^{n + 1}$ である。
\end{proof}

\begin{example}
\label{derham-example-affine-line}
補題 \ref{derham-lemma-line-bundle-characteristic-zero} は
正標数では偽である。$\mathbf{A}^1_k = \Spec(k[x])$ の、体 $k$ 上のド・ラーム複体は直和
$$
k \oplus
\bigoplus\nolimits_{n \geq 1}
(k \cdot t^n \xrightarrow{n}
k \cdot t^{n - 1} \text{d}t)
$$
の形をしている。したがって $k$ の標数が $p > 0$ ならば、
$H^0_{dR}(\mathbf{A}^1_k/k)$ と
$H^1_{dR}(\mathbf{A}^1_k/k)$ はともに
$k$ 上無限次元であることが分かる。
\end{example}








\section{射影空間のド・ラームコホモロジー}
\label{derham-section-projective-space}

\noindent
$A$ を環とし、$n \geq 1$ とする。構造射
$\mathbf{P}^n_A \to \Spec(A)$ は相対次元 $n$ の固有かつ滑らかな射である。
これは相対次元 $n$ の滑らかな有限型射でもある。実際、
$\mathbf{P}^n_A$ は、それぞれ $\mathbf{A}^n_A$ と同型なスキームからなる
有限アフィン開被覆をもつ。Constructions, Lemma
\ref{constructions-lemma-standard-covering-projective-space} を参照されたい。
また、Constructions, Lemma
\ref{constructions-lemma-projective-space-separated} により分離的かつ普遍閉であるため、固有である。
以下、$\mathcal{O}$ と $\mathcal{O}(d)$ により、それぞれ構造層
$\mathcal{O}_{\mathbf{P}^n_A}$ と Serre 捩り
$\mathcal{O}_{\mathbf{P}^n_A}(d)$ を表す。
相対微分の層を $\Omega = \Omega_{\mathbf{P}^n_A/A}$ と書き、
その外冪を $\Omega^p$ と書く。

\begin{lemma}
\label{derham-lemma-euler-sequence}
次の短完全列が存在する。
$$
0 \to \Omega \to \mathcal{O}(-1)^{\oplus n + 1} \to \mathcal{O} \to 0
$$
\end{lemma}

\begin{proof}
これを説明するため、
$\mathbf{P}^n_A = \text{Proj}(A[T_0, \ldots, T_n])$
であることを思い出し、記号的に
$$
\mathcal{O}(-1)^{\oplus n + 1} =
\bigoplus\nolimits_{j = 0, \ldots, n} \mathcal{O}(-1) \text{d}T_j
$$
と書く。上の短完全列の第一の矢印
$$
\Omega \to
\bigoplus\nolimits_{j = 0, \ldots, n} \mathcal{O}(-1) \text{d}T_j
$$
は、上で述べた各標準開集合
$D_+(T_i) = \Spec(A[T_0/T_i, \ldots, T_n/T_i])$
上で次の規則により与えられる。
$$
\sum\nolimits_{j \not = i} g_j \text{d}(T_j/T_i)
\longmapsto
\sum\nolimits_{j \not = i} g_j/T_i \text{d}T_j
- (\sum\nolimits_{j \not = i} g_jT_j/T_i^2) \text{d}T_i
$$
これは $1/T_i$ が $\mathcal{O}(-1)$ の $D_+(T_i)$ 上の切断であるため意味をもつ。
写像
$$
\bigoplus\nolimits_{j = 0, \ldots, n} \mathcal{O}(-1) \text{d}T_j
\to
\mathcal{O}
$$
は $\text{d}T_j$ を $T_j$ へ送ることで与えられる。より正確には、
$D_+(T_i)$ 上の切断 $\sum g_j \text{d}T_j$ を
$\sum T_jg_j$ へ送る。これが短完全列を与えることの確認は省略する。
\end{proof}

\noindent
整数 $k \in \mathbf{Z}$ と準連接
$\mathcal{O}_{\mathbf{P}^n_A}$-加群 $\mathcal{F}$ に対し、通常どおり
$\mathcal{F}(k)$ と書く。これは第 $k$ Serre 捩りを $\mathcal{F}$ に施したものである。
Constructions, Definition \ref{constructions-definition-twist} を参照されたい。

\begin{lemma}
\label{derham-lemma-twisted-hodge-cohomology-projective-space}
上の状況では、コホモロジー群について次が成り立つ。
\begin{enumerate}
\item $H^q(\mathbf{P}^n_A, \Omega^p) = 0$
ただし $0 \leq p = q \leq n$ の場合を除く。
\item $0 \leq p \leq n$ に対し、$A$-加群
$H^p(\mathbf{P}^n_A, \Omega^p)$ は階数 $1$ の自由加群である。
\item $q > 0$、$k > 0$、および任意の $p$ に対して
$H^q(\mathbf{P}^n_A, \Omega^p(k)) = 0$ である。
\item ここにさらに追加する。
\end{enumerate}
\end{lemma}

\begin{proof}
以下、Cohomology of Schemes, Lemma
\ref{coherent-lemma-cohomology-projective-space-over-ring}
の結果を断りなく用いる。特に、各主張は
$H^q(\mathbf{P}^n_A, \mathcal{O}(k))$ に対して成り立つ。

\medskip\noindent
$p = 1$ の場合を証明する。Lemma \ref{derham-lemma-euler-sequence} の短完全列
$$
0 \to \Omega \to \mathcal{O}(-1)^{\oplus n + 1} \to \mathcal{O} \to 0
$$
を考える。$\mathcal{O}(-1)$ のコホモロジーはすべての次数で消えるので、
$H^q(\mathbf{P}^n_A, \Omega)$ は次数 $1$ を除いて零であり、その次数では
$1$ の境界を自由生成元とする。ここで境界を取る元は
$H^0(\mathbf{P}^n_A, \mathcal{O})$ に属する。

\medskip\noindent
$p > 1$ と仮定する。上の短完全列を、$2$ 段フィルトレーションを
$\mathcal{O}(-1)^{\oplus n + 1}$ 上に定めるものとみなす。
$\wedge^p\mathcal{O}(-1)^{\oplus n + 1}$ 上に誘導されるフィルトレーションは
次の形をしている。
$$
0 \to \Omega^p \to \wedge^p\left(\mathcal{O}(-1)^{\oplus n + 1}\right)
\to \Omega^{p - 1} \to 0
$$
$\wedge^p\mathcal{O}(-1)^{\oplus n + 1}$ は、$n + 1$ 個から $p$ 個選ぶ数だけの
$\mathcal{O}(-p)$ の直和と同型である。したがって、
そのコホモロジーはすべての次数で消える。帰納法の仮定により、
$H^q(\mathbf{P}^n_A, \Omega^p)$ は $q = p$ の場合を除いて零であり、
$H^p(\mathbf{P}^n_A, \Omega^p)$ は階数 $1$ の自由加群で、その生成元は
$H^{p - 1}(\mathbf{P}^n_A, \Omega^{p - 1})$ の生成元の境界である。

\medskip\noindent
$k > 0$ とする。$\Omega^n = \mathcal{O}(-n - 1)$ である。
これは、例えば上の短完全列で $p = n + 1$ とおくことにより分かる。したがって、
$\Omega^n(k)$ のコホモロジーは正の次数で消える。
短完全列
$$
0 \to \Omega^p(k) \to \wedge^p\left(\mathcal{O}(-1)^{\oplus n + 1}\right)(k)
\to \Omega^{p - 1}(k) \to 0
$$
と $p$ に関する{\it 降下}帰納法を用いると、$\Omega^p(k)$ の
コホモロジーが、すべての $p$ に対して正の次数で消えることを得る。
\end{proof}

\begin{lemma}
\label{derham-lemma-hodge-cohomology-projective-space}
$H^q(\mathbf{P}^n_A, \Omega^p) = 0$ である。ただし
$0 \leq p = q \leq n$ の場合を除く。$0 \leq p \leq n$ に対し、
$A$-加群 $H^p(\mathbf{P}^n_A, \Omega^p)$ は階数 $1$ の自由加群で、
基底元は $c_1^{Hodge}(\mathcal{O}(1))^p$ である。
\end{lemma}

\begin{proof}
消滅と自由性は Lemma
\ref{derham-lemma-twisted-hodge-cohomology-projective-space} による。
$p = 0$ では、$1 \in H^0(\mathbf{P}^n_A, \mathcal{O})$ が生成元であることは明らかである。

\medskip\noindent
$p = 1$ の場合を証明する。Lemma \ref{derham-lemma-euler-sequence} の短完全列
$$
0 \to \Omega \to \mathcal{O}(-1)^{\oplus n + 1} \to \mathcal{O} \to 0
$$
を考える。Lemma \ref{derham-lemma-twisted-hodge-cohomology-projective-space}
の証明で、$H^1(\mathbf{P}^n_A, \Omega)$ の生成元が
境界 $\xi$ であり、その原像が
$1 \in H^0(\mathbf{P}^n_A, \mathcal{O})$ であることを見た。
Lemma \ref{derham-lemma-euler-sequence} の証明と同様に、
$\mathcal{O}(-1)^{\oplus n + 1}$ を
$\bigoplus_{j = 0, \ldots, n} \mathcal{O}(-1)\text{d}T_j$
と同一視する。開被覆
$$
\mathcal{U} : 
\mathbf{P}^n_A =
\bigcup\nolimits_{i = 0, \ldots, n} D_{+}(T_i)
$$
を考える。大域切断 $1$（$\mathcal{O}$ の切断）の
$U_i = D_+(T_i)$ への制限は、切断 $T_i^{-1} \text{d}T_i$ に持ち上げられる。
これは $\bigoplus \mathcal{O}(-1)\text{d}T_j$ の $U_i$ 上の切断である。したがって、
$\xi$ を表すコサイクルは次で与えられる。
$$
T_{i_1}^{-1} \text{d}T_{i_1} - T_{i_0}^{-1} \text{d}T_{i_0} =
\text{d}\log(T_{i_1}/T_{i_0}) \in \Omega(U_{i_0i_1})
$$
一方、各 $i$ に対し切断 $T_i$ は $\mathcal{O}(1)$ を
$U_i$ 上で自明化する切断である。したがって、
$f_{i_0i_1} = T_{i_1}/T_{i_0} \in \mathcal{O}^*(U_{i_0i_1})$
が $\mathcal{O}(1)$ を $\Pic(\mathbf{P}^n_A)$ において表すコサイクルである。
Section \ref{derham-section-first-chern-class} を参照されたい。
したがって $c_1^{Hodge}(\mathcal{O}(1))$ は
コサイクル $\text{d}\log(T_{i_1}/T_{i_0})$ により与えられ、
これは上で $\xi$ について得たものと一致する。

\medskip\noindent
一般の $p$ について帰納法で証明する。基底の場合 $p = 0, 1$ は上で扱った。
$p > 1$ と仮定する。Lemma
\ref{derham-lemma-twisted-hodge-cohomology-projective-space} の証明で、
$H^p(\mathbf{P}^n_A, \Omega^p)$ の生成元は、次の短完全列に付随する
コホモロジー長完全列における
$c_1^{Hodge}(\mathcal{O}(1))^{p - 1}$ の境界であることを見た。
$$
0 \to \Omega^p \to \wedge^p\left(\mathcal{O}(-1)^{\oplus n + 1}\right)
\to \Omega^{p - 1} \to 0
$$
Section \ref{derham-section-first-chern-class} の計算により、
コホモロジー類 $c_1^{Hodge}(\mathcal{O}(1))^{p - 1}$ は、符号を除いて、
次の項をもつコサイクルで表される。
$$
\beta_{i_0 \ldots i_{p - 1}} =
\text{d}\log(T_{i_1}/T_{i_0}) \wedge
\text{d}\log(T_{i_2}/T_{i_1}) \wedge \ldots \wedge
\text{d}\log(T_{i_{p - 1}}/T_{i_{p - 2}})
$$
これは $\Omega^{p - 1}(U_{i_0 \ldots i_{p - 1}})$ の元である。
これらの $\beta_{i_0 \ldots i_{p - 1}}$ は、切断
$\tilde \beta_{i_0 \ldots i_{p -1}} =
T_{i_0}^{-1}\text{d}T_{i_0} \wedge \beta_{i_0 \ldots i_{p - 1}}$
へ持ち上げられる。これは
$\wedge^p(\bigoplus \mathcal{O}(-1) \text{d}T_j)$ の
$U_{i_0 \ldots i_{p - 1}}$ 上の切断である。
したがって、$H^p(\mathbf{P}^n_A, \Omega^p)$ の生成元は、
次の成分をもつコサイクルで与えられる。
\begin{align*}
\sum\nolimits_{a = 0}^p (-1)^a
\tilde \beta_{i_0 \ldots \hat{i_a} \ldots i_p}
& =
T_{i_1}^{-1}\text{d}T_{i_1} \wedge \beta_{i_1 \ldots i_p}
+ \sum\nolimits_{a = 1}^p (-1)^a
T_{i_0}^{-1}\text{d}T_{i_0} \wedge
\beta_{i_0 \ldots \hat{i_a} \ldots i_p} \\
& =
(T_{i_1}^{-1}\text{d}T_{i_1} - T_{i_0}^{-1}\text{d}T_{i_0}) \wedge
\beta_{i_1 \ldots i_p} +
T_{i_0}^{-1}\text{d}T_{i_0} \wedge \text{d}(\beta)_{i_0 \ldots i_p} \\
& =
\text{d}\log(T_{i_1}/T_{i_0}) \wedge \beta_{i_1 \ldots i_p}
\end{align*}
これを $\Omega^p$ の $U_{i_0 \ldots i_p}$ 上の切断とみなす。
符号を除けば、これは $c_1^{Hodge}(\mathcal{O}(1))^p$ を表す
コサイクルと同じであり、証明が完了する。
\end{proof}

\begin{lemma}
\label{derham-lemma-de-rham-cohomology-projective-space}
$0 \leq i \leq n$ に対し、ド・ラームコホモロジー
$H^{2i}_{dR}(\mathbf{P}^n_A/A)$ は自由 $A$-加群で階数 $1$ をもち、
基底元は $c_1^{dR}(\mathcal{O}(1))^i$ である。
それ以外のすべての次数では、$\mathbf{P}^n_A$ の $A$ 上の
ド・ラームコホモロジーは零である。
\end{lemma}

\begin{proof}
Section \ref{derham-section-hodge-to-de-rham} のホッジ・ド・ラームスペクトル系列を考える。
Lemma \ref{derham-lemma-hodge-cohomology-projective-space} で行った
$\mathbf{P}^n_A$ の $A$ 上の Hodgeコホモロジーの計算により、
このスペクトル系列は $E_1$ 頁で退化する。
したがって、$H^{2i}_{dR}(\mathbf{P}^n_A/A)$ は自由 $A$-加群で階数 $1$ をもち、
$0 \leq i \leq n$ の場合にこれが成り立ち、それ以外では零である。
また、$c_1^{dR}(\mathcal{O}(1))^i \in H^{2i}_{dR}(\mathbf{P}^n_A/A)$ であり、
これは $i = 0, \ldots, n$ に対して成り立つ。$i = n$ の場合、この元は
$c_1^{Hodge}(\mathcal{L})^n$ の、複体の写像
$$
\Omega^n_{\mathbf{P}^n_A/A}[-n]
\longrightarrow
\Omega^\bullet_{\mathbf{P}^n_A/A}
$$
による像である。
これは、例えば Remark \ref{derham-remark-truncations} の議論、または
Section \ref{derham-section-first-chern-class} におけるこれらの類を表す
コサイクルの明示的記述から従う。スペクトル系列により、誘導される写像
$$
H^n(\mathbf{P}^n_A, \Omega^n_{\mathbf{P}^n_A/A}) \longrightarrow
H^{2n}_{dR}(\mathbf{P}^n_A/A)
$$
は同型である。また $c_1^{Hodge}(\mathcal{L})^n$ はその始域の生成元なので
（Lemma \ref{derham-lemma-hodge-cohomology-projective-space}）、
$c_1^{dR}(\mathcal{L})^n$ は終域の生成元である。
カップ積の $A$-双線形性により、$c_1^{dR}(\mathcal{L})^i$ も
$H^{2i}_{dR}(\mathbf{P}^n_A/A)$ の生成元であることが従い、
これは $0 \leq i \leq n$ に対して成り立つ。
\end{proof}









\section{滑らかな射に対するスペクトル系列}
\label{derham-section-relative-spectral-sequence}

\noindent
スキームの可換図式
$$
\xymatrix{
X \ar[rr]_f \ar[rd]_p & & Y \ar[ld]^q \\
& S
}
$$
を考え、$f$ は滑らかな射であるとする。このとき、局所的に分裂する短完全列
$$
0 \to f^*\Omega_{Y/S} \to \Omega_{X/S} \to \Omega_{X/Y} \to 0
$$
を得る。これは Morphisms, Lemma
\ref{morphisms-lemma-triangle-differentials-smooth} による。
これを降下フィルトレーション $F$ とみなす。これは $\Omega_{X/S}$ 上のフィルトレーションである。
ここで $F^0\Omega_{X/S} = \Omega_{X/S}$、
$F^1\Omega_{X/S} = f^*\Omega_{Y/S}$、および
$F^2\Omega_{X/S} = 0$ である。関手 $\wedge^p$ を適用すると、
各 $p$ に対して誘導フィルトレーション
$$
\Omega^p_{X/S} = F^0\Omega^p_{X/S} \supset
F^1\Omega^p_{X/S} \supset
F^2\Omega^p_{X/S} \supset \ldots \supset F^{p + 1}\Omega^p_{X/S} = 0
$$
を得る。その逐次商は
$$
\text{gr}^k\Omega^p_{X/S} =
F^k\Omega^p_{X/S}/F^{k + 1}\Omega^p_{X/S} =
f^*\Omega^k_{Y/S} \otimes_{\mathcal{O}_X} \Omega^{p - k}_{X/Y} =
f^{-1}\Omega^k_{Y/S} \otimes_{f^{-1}\mathcal{O}_Y} \Omega^{p - k}_{X/Y}
$$
であり、$k = 0, \ldots, p$ に対して成り立つ。実際、Leibniz則を用いれば、
$F^k\Omega^\bullet_{X/S}$ が $\Omega^\bullet_{X/S}$ の部分複体であることを
読者は確認できる。このように $\Omega^\bullet_{X/S}$ はフィルター付き複体の構造をもつ。
これは次を観察することによっても分かる。
$$
F^k\Omega^\bullet_{X/S} =
\Im\left(\wedge :
\text{Tot}(
f^{-1}\sigma_{\geq k}\Omega^\bullet_{Y/S} \otimes_{p^{-1}\mathcal{O}_S}
\Omega^\bullet_{X/S})
\longrightarrow
\Omega^\bullet_{X/S}\right)
$$
は $X$ 上の複体の写像の像である。フィルター付き複体
$$
\Omega^\bullet_{X/S} = F^0\Omega^\bullet_{X/S} \supset
F^1\Omega^\bullet_{X/S} \supset F^2\Omega^\bullet_{X/S} \supset \ldots
$$
には、次の付随する次数付き部分がある。
$$
\text{gr}^k\Omega^\bullet_{X/S} = 
f^{-1}\Omega^k_{Y/S}[-k] \otimes_{f^{-1}\mathcal{O}_Y} \Omega^\bullet_{X/Y}
$$
これは上で述べたことによる。

\begin{lemma}
\label{derham-lemma-spectral-sequence-smooth}
$f : X \to Y$ を、基底スキーム $S$ 上のスキームの準コンパクト、
準分離的かつ滑らかな射とする。第一頁が
$$
E_1^{p, q} =
H^q(\Omega^p_{Y/S} \otimes_{\mathcal{O}_Y}^\mathbf{L} Rf_*\Omega^\bullet_{X/Y})
$$
である有界スペクトル系列が存在し、
$R^{p + q}f_*\Omega^\bullet_{X/S}$ に収束する。
\end{lemma}

\begin{proof}
上で導入したフィルトレーションにより、
$\Omega^\bullet_{X/S}$ をフィルター付き複体とみなす。このスペクトル系列は
Cohomology, Lemma
\ref{cohomology-lemma-relative-spectral-sequence-filtered-object}
のスペクトル系列である。Derived Categories of Schemes, Lemma
\ref{perfect-lemma-cohomology-de-rham-base-change} により
$$
Rf_*\text{gr}^k\Omega^\bullet_{X/S} =
\Omega^k_{Y/S}[-k] \otimes_{\mathcal{O}_Y}^\mathbf{L} Rf_*\Omega^\bullet_{X/Y}
$$
を得るので、結論が従う。
\end{proof}

\begin{remark}
\label{derham-remark-gauss-manin}
Lemma \ref{derham-lemma-spectral-sequence-smooth} において、コホモロジー層
$$
\mathcal{H}^q_{dR}(X/Y) = H^q(Rf_*\Omega^\bullet_{X/Y})
$$
を考える。$f$ が滑らかであることに加えて固有であり、$S$ が
$\mathbf{Q}$ 上のスキームならば、$\mathcal{H}^q_{dR}(X/Y)$ は有限局所自由である
（将来の参照をここに挿入する）。$\mathcal{H}^q_{dR}(X/Y)$ が
平坦な $\mathcal{O}_Y$-加群であることだけを仮定しても、
次を得る（短い議論は省略する）。
$$
E_1^{p, q} =
\Omega^p_{Y/S} \otimes_{\mathcal{O}_Y} \mathcal{H}^q_{dR}(X/Y)
$$
また、このスペクトル系列の微分は写像
$$
d_1^{p, q} :
\Omega^p_{Y/S} \otimes_{\mathcal{O}_Y} \mathcal{H}^q_{dR}(X/Y)
\longrightarrow
\Omega^{p + 1}_{Y/S} \otimes_{\mathcal{O}_Y} \mathcal{H}^q_{dR}(X/Y)
$$
である。特に $p = 0$ に対して写像
$d_1^{0, q} : \mathcal{H}^q_{dR}(X/Y) \to
\Omega^1_{Y/S} \otimes_{\mathcal{O}_Y} \mathcal{H}^q_{dR}(X/Y)$
を得る。これは可積分接続
$\nabla$ となる（将来の参照をここに挿入する）。また、複体
$$
\mathcal{H}^q_{dR}(X/Y) \to
\Omega^1_{Y/S} \otimes_{\mathcal{O}_Y} \mathcal{H}^q_{dR}(X/Y) \to
\Omega^2_{Y/S} \otimes_{\mathcal{O}_Y} \mathcal{H}^q_{dR}(X/Y) \to \ldots
$$
は、微分が $d_1^{\bullet, q}$ で与えられ、
$\nabla$ のド・ラーム複体である。
接続 $\nabla$ は{\it Gauss--Manin 接続}として知られる。
\end{remark}






\section{Leray--Hirsch 型定理}
\label{derham-section-leray-hirsch}

\noindent
この節では、滑らかで固有な射について、全空間側のド・ラームコホモロジーを
基底側のド・ラームコホモロジーによって表せる場合があることを証明する。

\begin{lemma}
\label{derham-lemma-relative-global-generation-on-fibres}
スキームの滑らかで固有な射 $f : X \to Y$ をとる。
$N$ と $n_1, \ldots, n_N \geq 0$ を整数とし、
$\xi_i \in H^{n_i}_{dR}(X/Y)$、$1 \leq i \leq N$ をとる。
すべての点 $y \in Y$ に対し、$\xi_1, \ldots, \xi_N$ の像が
$H^*_{dR}(X_y/y)$ において $\kappa(y)$ 上の基底をなすと仮定する。このとき写像
$$
\bigoplus\nolimits_{i = 1}^N \mathcal{O}_Y[-n_i]
\longrightarrow
Rf_*\Omega^\bullet_{X/Y}
$$
は $\xi_1, \ldots, \xi_N$ に付随し、同型である。
\end{lemma}

\begin{proof}
Lemma \ref{derham-lemma-proper-smooth-de-Rham} により、
$Rf_*\Omega^\bullet_{X/Y}$ は $D(\mathcal{O}_Y)$ の完全対象であり、
その形成は任意の基底変換と可換である。したがって補題の写像
$a : K \to L$ は $D(\mathcal{O}_Y)$ の完全対象間の写像である。
ファイバーに関する仮定により、その任意の点への導来的制限は同型である。
このとき錐 $C$ を $a$ に対して取ると、これは $D(\mathcal{O}_Y)$ の完全対象である
（Cohomology, Lemma
\ref{cohomology-lemma-two-out-of-three-perfect}）。その任意の点への
導来的制限は零である。したがって $C$ は
More on Algebra, Lemma
\ref{more-algebra-lemma-lift-perfect-from-residue-field}
により零であり、$a$ は同型である。
（代数へ移すために Derived Categories of Schemes, Lemmas
\ref{perfect-lemma-affine-compare-bounded} および
\ref{perfect-lemma-perfect-affine} も用いる。）
\end{proof}

\noindent
まず、この節の主結果を次の特別な場合に証明する。

\begin{lemma}
\label{derham-lemma-global-generation-on-fibres}
$f : X \to Y$ をスキームの滑らかで固有な射とし、基底 $S$ 上で考える。
次を仮定する。
\begin{enumerate}
\item $Y$ と $S$ はアフィンである。
\item 整数 $N$ と $n_1, \ldots, n_N \geq 0$、および
$\xi_i \in H^{n_i}_{dR}(X/S)$、$1 \leq i \leq N$ が存在し、
すべての点 $y \in Y$ に対し、$\xi_1, \ldots, \xi_N$ の像は
$H^*_{dR}(X_y/y)$ において $\kappa(y)$ 上の基底をなす。
\end{enumerate}
このとき写像
$$
\bigoplus\nolimits_{i = 1}^N H^*_{dR}(Y/S) \longrightarrow
H^*_{dR}(X/S), \quad
(a_1, \ldots, a_N) \longmapsto  \sum \xi_i \cup f^*a_i
$$
は同型である。
\end{lemma}

\begin{proof}
$Y = \Spec(A)$ および $S = \Spec(R)$ と書く。
この場合、Lemma \ref{derham-lemma-de-rham-affine} により
$\Omega^\bullet_{A/R}$ は $R\Gamma(Y, \Omega^\bullet_{Y/S})$ を計算する。
有限アフィン開被覆 $\mathcal{U} : X = \bigcup_{i \in I} U_i$ を選び、複体
$$
K^\bullet =
\text{Tot}(\check{\mathcal{C}}^\bullet(\mathcal{U}, \Omega_{X/S}^\bullet))
$$
を考える。これは Cohomology, Section
\ref{cohomology-section-cech-cohomology-of-complexes} と同様である。
この複体についていくつかの事実をまとめる。その大部分は今挙げた参照にある。
\begin{enumerate}
\item $K^\bullet$ は $R$-加群の複体であり、その各項は $A$-加群である。
\item $K^\bullet$ は $R\Gamma(X, \Omega^\bullet_{X/S})$ を
$D(R)$ において表す
（Cohomology of Schemes, Lemma
\ref{coherent-lemma-quasi-coherent-affine-cohomology-zero} および
Cohomology, Lemma \ref{cohomology-lemma-cech-complex-complex-computes}）。
\item 自然な写像 $\Omega^\bullet_{A/R} \to K^\bullet$ が
$R$-加群の複体の写像として存在する。これは各項上で
$A$-線形であり、コホモロジー上で引き戻し写像
$H^*_{dR}(Y/S) \to H^*_{dR}(X/S)$ を誘導する。
\item $K^\bullet$ は $\wedge$ と書く乗法をもち、
これにより微分次数付き $R$-代数となる。
\item $K^\bullet$ 上の乗法は $H^*_{dR}(X/S)$ 上のカップ積を誘導する
（Cohomology, Section \ref{cohomology-section-cup-product}）。
\item フィルトレーション $F$ は $\Omega^*_{X/S}$ 上のものであり、フィルトレーション
$$
K^\bullet =
F^0K^\bullet \supset F^1K^\bullet \supset F^2K^\bullet \supset \ldots
$$
を $K^\bullet$ 上に、部分複体からなるものとして誘導し、次を満たす。
\begin{enumerate}
\item $F^kK^n \subset K^n$ は $A$-部分加群である。
\item $F^kK^\bullet \wedge F^lK^\bullet \subset F^{k + l}K^\bullet$。
\item $\text{gr}^kK^\bullet$ は $A$-加群の複体である。
\item $\text{gr}^0K^\bullet =
\text{Tot}(\check{\mathcal{C}}^\bullet(\mathcal{U}, \Omega_{X/Y}^\bullet))$
であり、$R\Gamma(X, \Omega^\bullet_{X/Y})$ を $D(A)$ において表す。
\item 乗法は同型
$\Omega^k_{A/R}[-k] \otimes_A \text{gr}^0K^\bullet \to \text{gr}^kK^\bullet$
を誘導する。
\end{enumerate}
\end{enumerate}
これらの主張の詳細な証明は省略する。Lemma
\ref{derham-lemma-spectral-sequence-smooth} のスペクトル系列の構成に至る議論を参照されたい。

\medskip\noindent
各 $i = 1, \ldots, N$ に対して、コサイクル $x_i \in K^{n_i}$ を選ぶ。
これは $\xi_i$ を表す。次に複体の写像
$$
\tilde x :
M^\bullet = \bigoplus\nolimits_{i = 1, \ldots, N}
\Omega^\bullet_{A/R}[-n_i]
\longrightarrow
K^\bullet
$$
を考える。これは $\omega$ が第 $i$ 直和因子に属するとき、
それを $x_i \wedge \omega$ へ送る。
残るのは、この写像が擬同型であることを示すことだけである。
$M^\bullet$ に、次の規則によってフィルター付き複体の構造を入れる。
$$
F^kM^\bullet =
\bigoplus\nolimits_{i = 1, \ldots, N}
(\sigma_{\geq k}\Omega^\bullet_{A/R})[-n_i]
$$
この選択により、写像 $\tilde x$ はフィルター付き複体の射となる。
$\text{gr}^0M^\bullet = \bigoplus A[-n_i]$ であり、乗法は同型
$\Omega^k_{A/R}[-k] \otimes_A \text{gr}^0M^\bullet \to \text{gr}^kM^\bullet$
を誘導する。構成と Lemma \ref{derham-lemma-relative-global-generation-on-fibres}
により、
$$
\text{gr}^0\tilde x :
\text{gr}^0M^\bullet \longrightarrow
\text{gr}^0K^\bullet
$$
は $D(A)$ における同型である。したがって、すべての $k \geq 0$ に対し同型
$$
\text{gr}^k \tilde x :
\text{gr}^kM^\bullet = \Omega^k_{A/R}[-k] \otimes_A \text{gr}^0M^\bullet
\longrightarrow
\Omega^k_{A/R}[-k] \otimes_A \text{gr}^0K^\bullet =
\text{gr}^kK^\bullet
$$
を $D(A)$ において得る。実際、複体
$\text{gr}^0K^\bullet =
\text{Tot}(\check{\mathcal{C}}^\bullet(\mathcal{U}, \Omega_{X/Y}^\bullet))$
は Derived Categories of Schemes, Lemma \ref{perfect-lemma-K-flat} により
$A$-加群の複体として K-平坦である。したがって右辺のテンソル積は
導来テンソル積であり、左辺についても直接確認できる。
最後に、導来テンソル積を取る関手
$\Omega^k_{A/R}[-k] \otimes_A^\mathbf{L} -$
は $D(A)$ 上の関手であるから、同型を同型へ送る。
$k$ に関する帰納法により、
$$
\tilde x : M^\bullet/F^kM^\bullet \to K^\bullet/F^kK^\bullet
$$
は $D(R)$ における同型である。実際、短完全列
$$
0 \to F^kM^\bullet/F^{k + 1}M^\bullet \to
M^\bullet/F^{k + 1}M^\bullet \to
\text{gr}^kM^\bullet \to 0
$$
があり、$K^\bullet$ についても同様である。
各次数でフィルトレーションは有限なので、これは $\tilde x$ が擬同型であることを証明する。
\end{proof}

\begin{proposition}
\label{derham-proposition-global-generation-on-fibres}
$f : X \to Y$ をスキームの滑らかで固有な射とし、基底 $S$ 上で考える。
$N$ と $n_1, \ldots, n_N \geq 0$ を整数とし、
$\xi_i \in H^{n_i}_{dR}(X/S)$、$1 \leq i \leq N$ をとる。
すべての点 $y \in Y$ に対し、$\xi_1, \ldots, \xi_N$ の像が
$H^*_{dR}(X_y/y)$ において $\kappa(y)$ 上の基底をなすと仮定する。写像
$$
\tilde \xi = \bigoplus \tilde \xi_i[-n_i] :
\bigoplus \Omega^\bullet_{Y/S}[-n_i]
\longrightarrow
Rf_*\Omega^\bullet_{X/S}
$$
は（証明を参照）$D(Y, (Y \to S)^{-1}\mathcal{O}_S)$ における同型であり、
これに対応して写像
$$
\bigoplus\nolimits_{i = 1}^N H^*_{dR}(Y/S) \longrightarrow
H^*_{dR}(X/S), \quad
(a_1, \ldots, a_N) \longmapsto  \sum \xi_i \cup f^*a_i
$$
は同型である。
\end{proposition}

\begin{proof}
構造射を $p : X \to S$ および $q : Y \to S$ と書く。
$\xi'_i : \Omega^\bullet_{X/S} \to \Omega^\bullet_{X/S}[n_i]$ を、
$\xi_i$ に対応する Remark \ref{derham-remark-cup-product-as-a-map} の写像とする。
また、
$$
\tilde \xi_i :
\Omega^\bullet_{Y/S} \to Rf_*\Omega^\bullet_{X/S}[n_i]
$$
を、$\xi'_i$ と標準写像
$\Omega^\bullet_{Y/S} \to Rf_*\Omega^\bullet_{X/S}$ との合成とする。
等式
$$
R\Gamma(Y, Rf_*\Omega^\bullet_{X/S}) = R\Gamma(X, \Omega^\bullet_{X/S})
$$
を用いると、コホモロジー上で $\tilde \xi_i$ は写像
$\eta \mapsto \xi_i \cup f^*\eta$、すなわち
$H^m_{dR}(Y/S)$ から $H^{m + n}_{dR}(X/S)$ への写像である。
さらに、$\xi'_i$ の形成は開集合への制限と可換なので、
$\tilde \xi_i$ の形成も開集合への制限と可換である。

\medskip\noindent
したがって写像
$$
\tilde \xi = \bigoplus \tilde \xi_i[-n_i] :
\bigoplus \Omega^\bullet_{Y/S}[-n_i]
\longrightarrow
Rf_*\Omega^\bullet_{X/S}
$$
を考える。補題を証明するには、これが
$D(Y, q^{-1}\mathcal{O}_S)$ における同型であることを示せば十分である。
$\tilde \xi$ が（適切なフィルトレーションを備えた）フィルター付き複体の
写像から生じることを示せれば、Lemma
\ref{derham-lemma-spectral-sequence-smooth} のスペクトル系列を用いて証明を完了できる。
しかしそれには必要以上の手間がかかるため、代わりにアフィンの場合へ還元する
（その証明は、ある意味でスペクトル系列を用いている）。

\medskip\noindent
実際、$Y' \subset Y$ を任意の開集合とし、その逆像を
$X' \subset X$ とすると、$\tilde \xi|_{X'}$ は写像
$$
\bigoplus\nolimits_{i = 1}^N H^*_{dR}(Y'/S) \longrightarrow
H^*_{dR}(X'/S), \quad
(a_1, \ldots, a_N) \longmapsto  \sum \xi_i|_{X'} \cup f^*a_i
$$
を $Y'$ 上のコホモロジーに誘導する。上の議論を参照されたい。
したがって、各要素について命題が成り立つような $Y$ の位相の基底を
見つければ十分である（特に、この操作では写像 $\tilde \xi$ を忘れてよい）。
これにより $Y$ と $S$ がアフィンの場合へ還元される。この場合は Lemma
\ref{derham-lemma-global-generation-on-fibres} で扱われており、証明が完了する。
\end{proof}





\section{射影束公式}
\label{derham-section-projective-space-bundle-formula}

\noindent
題名が内容を言い尽くしている。

\begin{proposition}
\label{derham-proposition-projective-space-bundle-formula}
$X \to S$ をスキームの射とする。$\mathcal{E}$ を局所自由
$\mathcal{O}_X$-加群で定階数 $r$ のものとする。射
$p : P = \mathbf{P}(\mathcal{E}) \to X$ を考える。このとき写像
$$
\bigoplus\nolimits_{i = 0, \ldots, r - 1} H^*_{dR}(X/S)
\longrightarrow
H^*_{dR}(P/S)
$$
は規則
$$
(a_0, \ldots, a_{r - 1}) \longmapsto
\sum\nolimits_{i = 0, \ldots, r - 1} c_1^{dR}(\mathcal{O}_P(1))^i \cup p^*(a_i)
$$
で与えられ、同型である。
\end{proposition}

\begin{proof}
$\Spec(A) \subset X$ をアフィン開集合として選び、そこへの
$\mathcal{E}$ の制限が自明な局所自由加群
$\mathcal{O}_{\Spec(A)}^{\oplus r}$ となるようにする。このとき
$P \times_X \Spec(A) = \mathbf{P}^{r - 1}_A$ である。したがって $p$ は
固有かつ滑らかである。Section \ref{derham-section-projective-space} を参照されたい。
さらに、類 $c_1^{dR}(\mathcal{O}_P(1))^i$（$i = 0, 1, \ldots, r - 1$）を
ファイバー $X_y = \mathbf{P}^{r - 1}_y$ に制限すると、
ド・ラームコホモロジー $H^*_{dR}(X_y/y)$ を $\kappa(y)$ 上自由に生成する。
Lemma \ref{derham-lemma-de-rham-cohomology-projective-space} を参照されたい。
したがって Proposition \ref{derham-proposition-global-generation-on-fibres}
の条件を確認でき、結論を得る。
\end{proof}

\begin{remark}
\label{derham-remark-projective-space-bundle-formula}
Proposition \ref{derham-proposition-projective-space-bundle-formula} の状況では、
さらに写像
$$
\tilde \xi :
\bigoplus\nolimits_{t = 0, \ldots, r - 1}
\Omega^\bullet_{X/S}[-2t]
\longrightarrow
Rp_*\Omega^\bullet_{P/S}
$$
は $D(X, (X \to S)^{-1}\mathcal{O}_X)$ における同型である。
これは Proposition \ref{derham-proposition-global-generation-on-fibres}
の適用から直ちに従う。$t = 0$ に対応する矢印は、単に標準写像
$c_{P/X} : \Omega^\bullet_{X/S} \to Rp_*\Omega^\bullet_{P/S}$
である。Section \ref{derham-section-de-rham-complex} を参照されたい。
実際、この場合にはこの写像をさらに具体的に記述できる。すなわち、標準写像
$$
\xi' : \Omega^\bullet_{P/S} \to \Omega^\bullet_{P/S}[2]
$$
を考える。これは $c_1^{dR}(\mathcal{O}_P(1))$ に対応する
Remark \ref{derham-remark-cup-product-as-a-map} の写像である。このとき
$$
\xi'[2(t - 1)] \circ \ldots \circ \xi'[2] \circ \xi' : 
\Omega^\bullet_{P/S} \to \Omega^\bullet_{P/S}[2t]
$$
は $c_1^{dR}(\mathcal{O}_P(1))^t$ に対応する
Remark \ref{derham-remark-cup-product-as-a-map} の写像である。
Proposition \ref{derham-proposition-global-generation-on-fibres} の証明で行った
選択をたどると、次を得る。
$$
\tilde \xi|_{\Omega^\bullet_{X/S}[-2t]} =
Rp_*\xi'[-2] \circ \ldots \circ Rp_*\xi'[-2(t - 1)] \circ
Rp_*\xi'[-2t] \circ c_{P/X}[-2t]
$$
これは、同型を直和因子 $\Omega^\bullet_{X/S}[-2t]$ に制限した値である。
ここから、後で用いる次の簡単な帰結を得る。すなわち、
$$
M = \bigoplus\nolimits_{t = 1, \ldots, r - 1} \Omega^\bullet_{X/S}[-2t]
\quad\text{and}\quad
K = \bigoplus\nolimits_{t = 0, \ldots, r - 2} \Omega^\bullet_{X/S}[-2t]
$$
を、矢印 $\tilde \xi$ の始域の部分複体とみなす。
上の議論から形式的に、
$$
c_{P/X} \oplus
\tilde \xi|_M :
\Omega^\bullet_{X/S} \oplus M \longrightarrow
Rp_*\Omega^\bullet_{P/S}
$$
は同型であり、図式
$$
\xymatrix{
K \ar[d]_{\tilde \xi|_K} \ar[r]_{\text{id}} &
M[2] \ar[d]^{(\tilde \xi|_M)[2]} \\
Rp_*\Omega^\bullet_{P/S} \ar[r]^{Rp_*\xi'} &
Rp_*\Omega^\bullet_{P/S}[2]
}
$$
は可換である。ここで $\text{id} : K \to M[2]$ は、
$t$ に対応する $K$ の分解の直和因子を、
$t + 1$ に対応する $M$ の分解の直和因子と同一視する。
\end{remark}









\section{因子に沿う対数極}
\label{derham-section-divisor}

\noindent
$X \to S$ をスキームの射とし、$Y \subset X$ を有効 Cartier 因子とする。
$X$ が $Y$ に沿ってエタール局所的に $Y \times \mathbf{A}^1$ の形をしているならば、
複体の標準短完全列
$$
0 \to \Omega^\bullet_{X/S} \to
\Omega^\bullet_{X/S}(\log Y) \to
\Omega^\bullet_{Y/S}[-1] \to 0
$$
が存在する。この列は多くの良い性質をもち、この節でそれらを論じる。
この構成には、正規交差因子から始める変種もある
(\'Etale Morphisms, Definition \ref{etale-definition-strict-normal-crossings})
が、これは別の箇所で論じる（将来の参照をここに挿入する）。

\begin{definition}
\label{derham-definition-local-product}
$X \to S$ をスキームの射とし、$Y \subset X$ を有効 Cartier 因子とする。
{\it $Y \subset X$ に対する $S$ 上の対数極ド・ラーム複体が定義される}
とは、すべての $y \in Y$ と、$f \in \mathcal{O}_{X, y}$ で与えられる
$Y$ の局所方程式に対して次が成り立つことをいう。
\begin{enumerate}
\item $\mathcal{O}_{X, y} \to \Omega_{X/S, y}$, $g \mapsto g \text{d}f$
は分裂単射である。
\item $\Omega^p_{X/S, y}$ は $f$-捩れをもたない（すべての $p$ に対して）。
\end{enumerate}
\end{definition}

\noindent
容易な局所計算により、各 $y \in Y$ に対して条件 (1) と (2) を満たす
局所方程式 $f$ が一つ見つかれば十分である。

\begin{lemma}
\label{derham-lemma-log-complex}
$X \to S$ をスキームの射とし、$Y \subset X$ を有効 Cartier 因子とする。
$Y \subset X$ に対する $S$ 上の対数極ド・ラーム複体が定義されると仮定する。
複体の標準短完全列
$$
0 \to \Omega^\bullet_{X/S} \to
\Omega^\bullet_{X/S}(\log Y) \to
\Omega^\bullet_{Y/S}[-1] \to 0
$$
が存在する。
\end{lemma}

\begin{proof}
仮定により、すべての $y \in Y$ と、$f \in \mathcal{O}_{X, y}$ で与えられる
$Y$ の局所方程式に対して
$$
\Omega_{X/S, y} = \mathcal{O}_{X, y}\text{d}f \oplus M
\quad\text{および}\quad
\Omega^p_{X/S, y} = \wedge^{p - 1}(M)\text{d}f \oplus \wedge^p(M)
$$
が成り立つ。ここで $M$ は、$f$-捩れをもたない外冪 $\wedge^p(M)$ をもつ加群である。
したがって
$$
\Omega^p_{Y/S, y} = \wedge^p(M/fM) = \wedge^p(M)/f\wedge^p(M)
$$
である。以下ではこれらの事実を暗黙に用いる。
特に、層 $\Omega^p_{X/S}$ は $Y$ に台をもつ非零局所切断をもたず、
標準包含
$$
\Omega^p_{X/S} \subset \Omega^p_{X/S}(Y)
$$
がある。More on Flatness, Section \ref{flat-section-eta} を参照されたい。
$U = \Spec(A)$ を、$Y \cap U = V(f)$ が成り立つような、ある
非零因子 $f \in A$ を伴うアフィン開部分スキームとする。
$\mathcal{O}_U$-部分加群である $\Omega^p_{X/S}(Y)|_U$ のうち、
$\Omega^p_{X/S}|_U$ と $\text{d}\log(f) \wedge \Omega^{p - 1}_{X/S}$
によって生成されるものを考える。ここで
$\text{d}\log(f) = f^{-1}\text{d}(f)$ である。
これは $f$ の選択に依存しない。実際、$Y$ の $U$ 上のイデアルの別の生成元は
$uf$ に等しく、ここで $u \in A$ は単元である。
$$
\text{d}\log(uf) - \text{d}\log(f) = \text{d}\log(u) = u^{-1}\text{d}u
$$
を得る。これは $\Omega_{X/S}$ の $U$ 上の切断である。
これらの局所層は貼り合わさって準連接部分加群
$$
\Omega^p_{X/S} \subset \Omega^p_{X/S}(\log Y) \subset \Omega^p_{X/S}(Y)
$$
を与える。以下、$\Omega^p_{Y/S}$ を準連接 $\mathcal{O}_X$-加群とみなす。
一意な全射 $\mathcal{O}_X$-線形写像
$$
\text{Res} : \Omega^p_{X/S}(\log Y) \to \Omega^{p - 1}_{Y/S}
$$
が存在し、規則
$$
\text{Res}(\eta' + \text{d}\log(f) \wedge \eta) = \eta|_{Y \cap U}
$$
によって定義される。ここで $U$ は上のような任意の開集合であり、
$\eta' \in \Omega^p_{X/S}(U)$ と $\eta \in \Omega^{p - 1}_{X/S}(U)$ は任意である。
$\eta$ を $U$ 上の形式とし、その $Y \cap U$ への制限が零ならば、
$\eta = \text{d}f \wedge \eta' + f\eta''$ と書ける。ここで
$\eta'$ と $\eta''$ は $U$ 上の形式である。したがって短完全列
$$
0 \to \Omega^p_{X/S} \to \Omega^p_{X/S}(\log Y) \to \Omega^{p - 1}_{Y/S} \to 0
$$
を、すべての $p$ に対して得る。さらに微分
$\Omega^p_{X/S}(\log Y) \to \Omega^{p + 1}_{X/S}(\log Y)$
を定義しなければならない。部分層 $\Omega^p_{X/S}$ 上では、
$X$ の $S$ 上のド・ラーム複体の微分を用いる。最後に
$\text{d}(\text{d}\log(f) \wedge \eta) = -\text{d}\log(f) \wedge \text{d}\eta$.
と定める。この符号は $\Omega^\bullet_{X/S}$ 上の Leibniz 則によって強制される。
また、Homology, Definition \ref{homology-definition-shift-cochain}
の符号規約により $\Omega^\bullet_{Y/S}[-1]$ 上の微分は
$-\text{d}_{Y/S}$ であるから、これとも両立する。
このようにして、補題で述べた複体の短完全列を得る。
\end{proof}

\begin{definition}
\label{derham-definition-log-complex}
$X \to S$ をスキームの射とし、$Y \subset X$ を有効 Cartier 因子とする。
$Y \subset X$ に対する $S$ 上の対数極ド・ラーム複体が定義されていると仮定する。
このとき複体
$$
\Omega^\bullet_{X/S}(\log Y)
$$
Lemma \ref{derham-lemma-log-complex} で構成されたものを、
{\it $Y \subset X$ に対する $S$ 上の対数極ド・ラーム複体}と呼ぶ。
\end{definition}

\noindent
この複体は多くの良い性質をもつ。

\begin{lemma}
\label{derham-lemma-multiplication-log}
$p : X \to S$ をスキームの射とし、$Y \subset X$ を有効 Cartier 因子とする。
$Y \subset X$ に対する $S$ 上の対数極ド・ラーム複体が定義されていると仮定する。
\begin{enumerate}
\item 写像
$\wedge : \Omega^p_{X/S} \times \Omega^q_{X/S} \to \Omega^{p + q}_{X/S}$
は一意的に $\mathcal{O}_X$-双線形写像へ拡張される。
$$
\wedge : \Omega^p_{X/S}(\log Y) \times \Omega^q_{X/S}(\log Y)
\to \Omega^{p + q}_{X/S}(\log Y)
$$
これは Leibniz 則
$
\text{d}(\omega \wedge \eta) = \text{d}(\omega) \wedge \eta +
(-1)^{\deg(\omega)} \omega \wedge \text{d}(\eta)$,
\item (1) の積に関して写像
$\Omega^\bullet_{X/S} \to \Omega^\bullet_{X/S}(\log(Y)$
は微分次数付き $\mathcal{O}_S$-代数の準同型である。
\item (1) の写像により $\Omega^p_{X/S}(\log Y) =
\wedge^p(\Omega^1_{X/S}(\log Y))$ であり、
\item 写像
$\text{Res} : \Omega^\bullet_{X/S}(\log Y) \to \Omega^\bullet_{Y/S}[-1]$
は
$$
\text{Res}(\omega \wedge \eta) = \text{Res}(\omega) \wedge \eta|_Y
$$
を満たす。ここで $\omega$ は $\Omega^p_{X/S}(\log Y)$ の局所切断、$\eta$ は
$\Omega^q_{X/S}$ の局所切断である。
\end{enumerate}
\end{lemma}

\begin{proof}
これは Lemma \ref{derham-lemma-log-complex} の証明における複体の局所構成から、
直接計算することで従う。詳細は省略する。
\end{proof}

\noindent
スキームの可換図式
$$
\xymatrix{
X' \ar[r]_f \ar[d] & X \ar[d] \\
S' \ar[r] & S
}
$$
を考える。$Y \subset X$ を有効 Cartier 因子とし、その引き戻し
$Y' = f^*Y$ が定義されているとする（Divisors, Definition
\ref{divisors-definition-pullback-effective-Cartier-divisor}).
さらに、$Y \subset X$ に対する $S$ 上の対数極ド・ラーム複体と、
$Y' \subset X'$ に対する $S'$ 上の対数極ド・ラーム複体が定義されていると仮定する。
このとき複体の短完全列の写像
$$
\xymatrix{
0 \ar[r] &
f^{-1}\Omega^\bullet_{X/S} \ar[r] \ar[d] &
f^{-1}\Omega^\bullet_{X/S}(\log Y) \ar[r] \ar[d] &
f^{-1}\Omega^\bullet_{Y/S}[-1] \ar[r] \ar[d] &
0 \\
0 \ar[r] &
\Omega^\bullet_{X'/S'} \ar[r] &
\Omega^\bullet_{X'/S'}(\log Y') \ar[r] &
\Omega^\bullet_{Y'/S'}[-1] \ar[r] &
0
}
$$
線形化すると、すべての $p$ に対して線形写像
$f^*\Omega^p_{X/S}(\log Y) \to \Omega^p_{X'/S'}(\log Y')$.

\begin{lemma}
\label{derham-lemma-gysin-via-log-complex}
$f : X \to S$ をスキームの射とし、$Y \subset X$ を有効 Cartier 因子とする。
$Y \subset X$ に対する $S$ 上の対数極ド・ラーム複体が定義されていると仮定する。
次を記す。
$$
\delta : \Omega^\bullet_{Y/S} \to \Omega^\bullet_{X/S}[2]
$$
$D(X, f^{-1}\mathcal{O}_S)$ における、Lemma \ref{derham-lemma-log-complex} の
短完全列から得られる「境界」写像である。次を記す。
$$
\xi' : \Omega^\bullet_{X/S} \to \Omega^\bullet_{X/S}[2]
$$
$D(X, f^{-1}\mathcal{O}_S)$ における、Remark \ref{derham-remark-cup-product-as-a-map} に
対応する写像であり、$\xi = c_1^{dR}(\mathcal{O}_X(-Y))$ に対応する。次を記す。
$$
\zeta' : \Omega^\bullet_{Y/S} \to \Omega^\bullet_{Y/S}[2]
$$
$D(Y, f|_Y^{-1}\mathcal{O}_S)$ における、Remark \ref{derham-remark-cup-product-as-a-map} に
対応する写像であり、$\zeta = c_1^{dR}(\mathcal{O}_X(-Y)|_Y)$ に対応する。
すると次の図式
$$
\xymatrix{
\Omega^\bullet_{X/S} \ar[d]_{\xi'} \ar[r] &
\Omega^\bullet_{Y/S} \ar[d]^{\zeta'} \ar[ld]_\delta \\
\Omega^\bullet_{X/S}[2] \ar[r] &
\Omega^\bullet_{Y/S}[2]
}
$$
$D(X, f^{-1}\mathcal{O}_S)$ において可換である。
\end{lemma}

\begin{proof}
より正確には、$\delta$ を次のシフトした短完全列に対応する境界写像として定義する。
$$
0 \to \Omega^\bullet_{X/S}[1] \to
\Omega^\bullet_{X/S}(\log Y)[1] \to
\Omega^\bullet_{Y/S} \to 0
$$
各三角形が可換であることを示せば十分である。$\mathcal{L} = \mathcal{O}_X(-Y)$ とおく。
$\pi : L \to X$ を、$\pi_*\mathcal{O}_L = \bigoplus_{n \geq 0} \mathcal{L}^{\otimes n}$
を満たす線束として記す。

\medskip\noindent
左上の三角形の可換性を示す。Lemma \ref{derham-lemma-the-complex-for-L-star-gives-chern-class} により、
$\xi'$ は Lemma \ref{derham-lemma-the-complex-for-L-star} で与えられる三角形の境界写像である。
函手性により、次の短完全列の射が存在することを示せば十分である。
$$
\xymatrix{
0 \ar[r] &
\Omega^\bullet_{X/S}[1] \ar[r] \ar[d] &
\Omega^\bullet_{L^\star/S, 0}[1] \ar[r] \ar[d] &
\Omega^\bullet_{X/S} \ar[r] \ar[d] &
0 \\
0 \ar[r] &
\Omega^\bullet_{X/S}[1] \ar[r] &
\Omega^\bullet_{X/S}(\log Y)[1] \ar[r] &
\Omega^\bullet_{Y/S} \ar[r] &
0
}
$$
左と右の垂直矢印は明らかなものである。中央の垂直矢印を次の規則で定義できる。
$$
\omega' + \text{d}\log(s) \wedge \omega \longmapsto
\omega' + \text{d}\log(f) \wedge \omega
$$
$\omega', \omega$ は $\Omega^\bullet_{X/S}$ の局所切断であり、$s$ は $\mathcal{L}$ の局所生成元、
$f \in \mathcal{O}_X(-Y)$ は $Y$ の $X$ におけるイデアル層に対応する切断である。
Lemmas \ref{derham-lemma-the-complex-for-L-star} と \ref{derham-lemma-log-complex} における写像の構成は
完全に一致するので、これでよい。

\medskip\noindent
右下の三角形の可換性を示す。$\overline{L}$ を $L$ の $Y$ への制限とする。
Lemma \ref{derham-lemma-the-complex-for-L-star-gives-chern-class} により、$\zeta'$ は
線束 $\overline{L}$ を $Y$ 上で用いた Lemma \ref{derham-lemma-the-complex-for-L-star} の
三角形の境界写像である。函手性により、次の短完全列の射が存在することを示せば十分である。
$$
\xymatrix{
0 \ar[r] &
\Omega^\bullet_{X/S}[1] \ar[r] \ar[d] &
\Omega^\bullet_{X/S}(\log Y)[1] \ar[r] \ar[d] &
\Omega^\bullet_{Y/S} \ar[r] \ar[d] &
0 \\
0 \ar[r] &
\Omega^\bullet_{Y/S}[1] \ar[r] &
\Omega^\bullet_{\overline{L}^\star/S, 0}[1] \ar[r] &
\Omega^\bullet_{Y/S} \ar[r] &
0 \\
}
$$
左と右の垂直矢印は明らかなものである。中央の垂直矢印を次の規則で定義できる。
$$
\omega' + \text{d}\log(f) \wedge \omega \longmapsto
\omega'|_Y + \text{d}\log(\overline{s}) \wedge \omega|_Y
$$
$\omega', \omega$ は $\Omega^\bullet_{X/S}$ の局所切断であり、$f$ は
$\mathcal{O}_X(-Y)$ の局所生成元で、$X$ 上の関数とみなす。また $\overline{s}$ は $f|_Y$ を
$\mathcal{L}|_Y = \mathcal{O}_X(-Y)|_Y$ の切断とみなしたものである。
Lemmas \ref{derham-lemma-the-complex-for-L-star} と \ref{derham-lemma-log-complex} における写像の構成は
完全に一致するので、これでよい。
\end{proof}

\begin{lemma}
\label{derham-lemma-log-complex-consequence}
$X \to S$ をスキームの射とし、$Y \subset X$ を有効 Cartier 因子とする。
$Y \subset X$ に対する $S$ 上の対数極ド・ラーム複体が定義されていると仮定する。
$b \in H^m_{dR}(X/S)$ を、$Y$ への制限が零であるド・ラームコホモロジー類とする。
このとき
$c_1^{dR}(\mathcal{O}_X(Y)) \cup b = 0$ in $H^{m + 2}_{dR}(X/S)$.
\end{lemma}

\begin{proof}
これは Lemma \ref{derham-lemma-gysin-via-log-complex} から直ちに従う。実際、
$$
c_1^{dR}(\mathcal{O}_X(Y)) \cup b =
-c_1^{dR}(\mathcal{O}_X(-Y)) \cup b = -\xi'(b) = -\delta(b|_Y) = 0
$$
である。第二の等式については Remark \ref{derham-remark-cup-product-as-a-map} を参照。
\end{proof}

\begin{lemma}
\label{derham-lemma-check-log-smooth}
$X \to T \to S$ をスキームの射とし、$Y \subset X$ を有効 Cartier 因子とする。
$X \to T$ と $Y \to T$ がともに滑らかならば、$Y \subset X$ に対する $S$ 上の
対数極ド・ラーム複体が定義される。
\end{lemma}

\begin{proof}
$y \in Y$ を一点とする。
More on Morphisms, Lemma \ref{more-morphisms-lemma-etale-local-structure} により、
整数 $0 \geq m$ と可換図式
$$
\xymatrix{
Y \ar[d] &
V \ar[l] \ar[d] \ar[r] &
\mathbf{A}^m_T
\ar[d]^{(a_1, \ldots, a_m) \mapsto (a_1, \ldots, a_m, 0)} \\
X &
U \ar[l] \ar[r]^-\pi &
\mathbf{A}^{m + 1}_T
}
$$
が存在する。ここで $U \subset X$ は開であり、$V = Y \cap U$、
$\pi$ はエタール、$V = \pi^{-1}(\mathbf{A}^m_T)$、$y \in V$ である。
$z \in \mathbf{A}^m_T$ を $y$ の像とする。このとき
$$
\Omega^p_{X/S, y} = \Omega^p_{\mathbf{A}^{m + 1}_T/S, z}
\otimes_{\mathcal{O}_{\mathbf{A}^{m + 1}_T, z}} \mathcal{O}_{X, x}
$$
Lemma \ref{derham-lemma-etale} による。$x_1, \ldots, x_{m + 1}$ を
$\mathbf{A}^{m + 1}_T$ の座標関数とする。
Definition \ref{derham-definition-local-product} の条件 (1) と (2) は局所座標の選択に依存しないので、
$f$ が平坦な局所環準同型による $x_{m + 1}$ の像である場合に条件 (1) と (2) を
確認すれば十分である。
$\mathcal{O}_{\mathbf{A}^{m + 1}_T, z} \to \mathcal{O}_{X, x}$.
以上により、$\mathbf{A}^m_T \subset \mathbf{A}^{m + 1}_T$ と点 $z$ について
条件 (1) と (2) を確認すれば十分であることが分かる。この場合を証明するため、
$S = \Spec(A)$ と $T = \Spec(B)$ がアフィンであると仮定してよい。環写像
$A \to B$ を射 $T \to S$ に対応するものとし、$P = B[x_1, \ldots, x_{m + 1}]$ とおく。
すると $\mathbf{A}^{m + 1}_T = \Spec(P)$ である。このとき
$$
\Omega_{P/A} = \Omega_{B/A} \otimes_B P \oplus
\bigoplus\nolimits_{j = 1, \ldots, m} P \text{d}x_j \oplus
P \text{d}x_{m + 1}
$$
したがって写像 $P \to \Omega_{P/A}$、$g \mapsto g \text{d}x_{m + 1}$ は分裂単射であり、
$x_{m + 1}$ は $\Omega^p_{P/A}$ 上の非零因子であり、これはすべての $p \geq 0$ に対して成り立つ。
$z$ に対応する素イデアルで局所化すれば証明が終わる。
\end{proof}

\begin{remark}
\label{derham-remark-check-log-completion-1}
$S$ を局所ネーター・スキームとする。$X$ を $S$ 上局所有限型とし、$Y \subset X$ を
有効 Cartier 因子とする。写像
$$
\mathcal{O}_{X, y}^\wedge \longrightarrow \mathcal{O}_{Y, y}^\wedge
$$
がすべての $y \in Y$ に対して切断をもつならば、$Y \subset X$ に対する $S$ 上の
対数極ド・ラーム複体が定義される。この結果が必要になった場合には、ここで正確な
主張を定式化し、証明を追加する。
\end{remark}

\begin{remark}
\label{derham-remark-check-log-completion-2}
$S$ を局所ネーター・スキームとする。$X$ を $S$ 上局所有限型とし、$Y \subset X$ を
有効 Cartier 因子とする。すべての $y \in Y$ に対して、$S$ 上のスキームの図式
$$
X \xleftarrow{\varphi} U \xrightarrow{\psi} V
$$
で、$\varphi$ がエタール、かつ
$\psi|_{\varphi^{-1}(Y)} : \varphi^{-1}(Y) \to V$ がエタールとなるものを取れるならば、
$Y \subset X$ に対する $S$ 上の対数極ド・ラーム複体が定義される。特別な場合として、
対 $(X, Y)$ がエタール局所的に $(V \times \mathbf{A}^1, V \times \{0\})$ の形をしている
場合がある。この結果が必要になった場合には、ここで正確な主張を定式化し、証明を追加する。
\end{remark}

















\section{計算}
\label{derham-section-calculations}

\noindent
この節では、標準的な爆破に対するいくつかのホッジコホモロジー群と
ド・ラームコホモロジー群を計算する。

\medskip\noindent
環 $R$ を固定し、$S = \Spec(R)$ とおく。整数 $0 \leq m$ と $1 \leq n$ を固定する。
閉埋め込み
$$
Z = \mathbf{A}^m_S \longrightarrow \mathbf{A}^{m + n}_S = X,\quad
(a_1, \ldots, a_m) \mapsto (a_1, \ldots, a_m, 0, \ldots 0).
$$
$L$ を $X$ の、閉部分スキーム $Z$ に沿う爆破として考える。次のように書く。
$$
X =
\mathbf{A}^{m + n}_S =
\Spec(A)
\quad\text{ただし}\quad
A = R[x_1, \ldots, x_m, y_1, \ldots, y_n]
$$
$A = R[x_1, \ldots, x_m, y_1, \ldots, y_n]$ を次数付き $R$-代数とみなし、
$\deg(x_i) = 0$、$\deg(y_j) = 1$ と定める。この次数付けにより
$$
P =
\text{Proj}(A) =
\mathbf{A}^m_S \times_S \mathbf{P}^{n - 1}_S =
Z \times_S \mathbf{P}^{n - 1}_S =
\mathbf{P}^{n - 1}_Z
$$
$Z$ を $X$ において定めるイデアルは $A_+$ である。したがって $L$ は Rees 代数
$$
A \oplus A_+ \oplus (A_+)^2 \oplus \ldots =
\bigoplus\nolimits_{d \geq 0} A_{\geq d}
$$
したがって $L$ は More on Morphisms, Section
\ref{more-morphisms-section-proj-spec} でより一般に研究した現象の一例である。
以下では、そこで得た観察を断りなく用いる。特に、可換図式
$$
\xymatrix{
P \ar[r]_0 \ar[d]_p &
L \ar[r]_-\pi \ar[d]^b &
P \ar[d]^p \\
Z \ar[r]^i &
X \ar[r] &
Z
}
$$
であり、$\pi : L \to P$ は
$P = Z \times_S \mathbf{P}^{n - 1}_S$ 上の線束である。零切断 $0$ の像
$E = 0(P) \subset L$ は爆破 $b$ の例外因子である。

\begin{lemma}
\label{derham-lemma-comparison}
任意の $a \geq 0$ に対して、次が成り立つ。
\begin{enumerate}
\item 写像 $\Omega^a_{X/S} \to b_*\Omega^a_{L/S}$ は同型である。
\item 写像 $\Omega^a_{Z/S} \to p_*\Omega^a_{P/S}$ は同型であり、
\item 写像 $Rb_*\Omega^a_{L/S} \to i_*Rp_*\Omega^a_{P/S}$ は、次数 $\geq 1$ の
コホモロジー層上で同型である。
\end{enumerate}
\end{lemma}

\begin{proof}
まず (2) を証明する。$P = Z \times_S \mathbf{P}^{n - 1}_S$ なので、
次を得る。
$$
\Omega^a_{P/S} = \bigoplus\nolimits_{a = r + s}
\text{pr}_1^*\Omega^r_{Z/S} \otimes
\text{pr}_2^*\Omega^s_{\mathbf{P}^{n - 1}_S/S}
$$
$p = \text{pr}_1$ であることを想起し、射影公式
(Cohomology, Lemma \ref{cohomology-lemma-projection-formula}) により、次を得る。
$$
p_*\Omega^a_{P/S} = \bigoplus\nolimits_{a = r + s}
\Omega^r_{Z/S} \otimes
\text{pr}_{1, *}\text{pr}_2^*\Omega^s_{\mathbf{P}^{n - 1}_S/S}
$$
Section \ref{derham-section-projective-space} の計算、特に Lemma
\ref{derham-lemma-hodge-cohomology-projective-space} の証明により、
$\text{pr}_{1, *}\text{pr}_2^*\Omega^s_{\mathbf{P}^{n - 1}_S/S} = 0$ である。
ただし $s = 0$ のときは $\text{pr}_{1, *}\mathcal{O}_P = \mathcal{O}_Z$ となる。
これで (2) が証明された。

\medskip\noindent
Section \ref{derham-section-line-bundle} の結果、特に Lemma \ref{derham-lemma-push-omega-a} により、
$\pi_*\Omega^a_{L/S} = \Omega^a_{P/S} \oplus
\bigoplus_{k \geq 1} \Omega^a_{L/S, k}$.
上の図式における合成 $\pi \circ 0$ は $P$ 上の恒等射である。
(3) を証明するには、$\Omega^a_{L/S, k}$ の高次コホモロジーが $k > 0$ で消滅することを
示せば十分である。Lemmas \ref{derham-lemma-the-complex-for-L-star} と
\ref{derham-lemma-push-omega-a} により、短完全列
$$
0 \to \Omega^a_{P/S} \otimes \mathcal{O}_P(k)
\to \Omega^a_{L/S, k} \to
\Omega^{a - 1}_{P/S} \otimes \mathcal{O}_P(k) \to 0
$$
ここで $\Omega^{a - 1}_{P/S} = 0$ （$a = 0$ の場合）である。$P = Z \times_S \mathbf{P}^{n - 1}_S$ なので、次を得る。
$$
\Omega^a_{P/S} = \bigoplus\nolimits_{i + j = a}
\Omega^i_{Z/S} \boxtimes \Omega^j_{\mathbf{P}^{n - 1}_S/S}
$$
Lemma \ref{derham-lemma-de-rham-complex-product} による。$\Omega^i_{Z/S}$ は有限ランク自由なので、
$\mathcal{O}_Z \boxtimes \Omega^j_{\mathbf{P}^{n - 1}_S/S}(k)$ の高次コホモロジーが
$k > 0$ で零であることを示せば十分である。これは Lemma
\ref{derham-lemma-twisted-hodge-cohomology-projective-space} を
$P = Z \times_S \mathbf{P}^{n - 1}_S = \mathbf{P}^{n - 1}_Z$ に適用すれば従い、
(3) の証明が完了する。

\medskip\noindent
残っているのは (1) の証明である。$n = 1$ ならば有効 Cartier 因子を爆破しているので、
$b$ は同型であり (1) が成り立つ。$n > 1$ のとき、合成
$$
\Gamma(X, \Omega^a_{X/S})
\to
\Gamma(L, \Omega^a_{L/S})
\to
\Gamma(L \setminus E, \Omega^a_{L/S})
=
\Gamma(X \setminus Z, \Omega^a_{X/S})
$$
これは $\Omega^a_{X/S}$ が有限自由であるため同型である（細部は省略する）。したがって
(1) が失敗し得る唯一の方法は、$\Gamma(L, \Omega^a_{L/S})$ の中に $E = 0(P)$ の外で消える
非零元が存在することである。$L$ は $P$ 上の線束で、
零切断 $0 : P \to L$ をもつので、線束上の微分形式の層に零切断の外で恒等的に消える
非零切断がないことを示せば十分である。これは（望ましくは）局所計算によって、または
$\Omega_{L/S, k} \subset \Omega_{L^\star/S, k}$ を用いて（上記の参照を見よ）分かる。
\end{proof}

\noindent
ここで読者は次の節へ進んでもよい。

\begin{lemma}
\label{derham-lemma-comparison-bis}
$a \geq 0$ に対して標準的な写像
$$
b^*\Omega^a_{X/S} \longrightarrow
\Omega^a_{L/S} \longrightarrow
b^*\Omega^a_{X/S} \otimes_{\mathcal{O}_L} \mathcal{O}_L((n - 1)E)
$$
が存在し、その合成は包含
$\mathcal{O}_L \subset \mathcal{O}_L((n - 1)E)$ によって誘導される。
\end{lemma}

\begin{proof}
表示式の第一の矢印は Section \ref{derham-section-de-rham-complex} で
論じた。第二の矢印を得るには、$\Omega^a_{L/S}$ の局所切断を
$b^*\Omega^a_{X/S}$ の「有理型切断」とみなしたとき、それが
$n - 1$ 以下の位数の極を $E$ に沿ってもつことを示せばよい。
これを見るため、$L$ のアフィン局所チャート上で考える。すなわち、
$L$ がアフィン爆発代数 $A[\frac{I}{y_i}]$ のスペクトルで覆われ、ここで
$I = A_{+}$ は $y_1, \ldots, y_n$ で生成されるイデアルであることを思い出す。
Algebra, Section \ref{algebra-section-blow-up} および Divisors, Lemma
\ref{divisors-lemma-blowing-up-affine} を参照せよ。
対称性により、$i = 1$ に対応するチャートだけを考えれば十分である。すると
$$
A[\frac{I}{y_1}] = R[x_1, \ldots, x_m, y_1, t_2, \ldots, t_n]
$$
となる。ここで $t_i = y_i/y_1$ である（More on Algebra, Lemma
\ref{more-algebra-lemma-blowup-regular-sequence} を参照）。したがって、
対応するアフィン開集合上で $\Omega^1_{L/S}$ は次の元により自由に生成される。
$\text{d}x_1, \ldots, \text{d}x_m$, $\text{d}y_1$、および
$\text{d}t_1, \ldots, \text{d}t_n$ である。もちろん、これらの生成元の最初の
$m + 1$ 個は $b^*\Omega^1_{X/S}$ から来る。残りの $n - 1$ 個については、
$$
\text{d}t_j =
\text{d}\frac{y_j}{y_1} =
\frac{1}{y_1}\text{d}y_j - \frac{y_j}{y_1^2}\text{d}y_1 =
\frac{\text{d}y_j - t_j \text{d}y_1}{y_1}
$$
を得る。これは位数 $1$ の極を $E$ に沿ってもつ。このチャートでは $E$ は $y_1$ で
切り出される。これらの元の $a$ 個の楔積がこのチャート上の
$\Omega^a_{L/S}$ の基底を与え、関与するものは高々 $n - 1$ 個の $\text{d}t_j$ であるから、
証明は完了する。
\end{proof}

\begin{lemma}
\label{derham-lemma-blowup-twist-same-cohomology}
$E = 0(P)$ を爆破 $b$ の例外因子とする。任意の局所自由
$\mathcal{O}_X$-加群 $\mathcal{E}$ と $0 \leq i \leq n - 1$ に対して、写像
$$
\mathcal{E}
\longrightarrow
Rb_*(b^*\mathcal{E} \otimes_{\mathcal{O}_L} \mathcal{O}_L(iE))
$$
は $D(\mathcal{O}_X)$ における同型である。
\end{lemma}

\begin{proof}
射影公式により、これは $\mathcal{E} = \mathcal{O}_X$ の場合に示せば十分である。
Cohomology, Lemma \ref{cohomology-lemma-projection-formula} を参照せよ。
$X$ はアフィンなので、写像
$$
H^0(X, \mathcal{O}_X) \to
H^0(L, \mathcal{O}_L) \to
H^0(L, \mathcal{O}_L(iE))
$$
が同型であり、$H^j(X, \mathcal{O}_L(iE)) = 0$ が
$j > 0$ および $0 \leq i \leq n - 1$ に対して成り立つことを示せばよい。
Cohomology of Schemes, Lemma
\ref{coherent-lemma-quasi-coherence-higher-direct-images-application}.
また $\pi$ はアフィンなので、$\pi_*$ を取った後に大域切断と
コホモロジーを計算できる。Cohomology of Schemes, Lemma
\ref{coherent-lemma-relative-affine-cohomology} を参照せよ。$n = 1$ ならば、
$L \to X$ は同型で $i = 0$ だから、第一の主張が成り立つ。
$n > 1$ のとき、次の合成を考える。
$$
H^0(X, \mathcal{O}_X) \to H^0(L, \mathcal{O}_L) \to
H^0(L, \mathcal{O}_L(iE)) \to H^0(L \setminus E, \mathcal{O}_L) =
H^0(X \setminus Z, \mathcal{O}_X)
$$
この場合、$H^0(X \setminus Z, \mathcal{O}_X) = H^0(X, \mathcal{O}_X)$ が成り立つ。
これは $Z$ の余次元が $n \geq 2$ であり、$X$ における余次元だからである（細部は省略する）。
したがって第一の主張が成り立つ。第二の主張については、
$\mathcal{O}_L(E) = \mathcal{O}_L(-1)$ であることを思い出す。Divisors, Lemma
\ref{divisors-lemma-blowing-up-gives-effective-Cartier-divisor}.
したがって
$$
\pi_*\mathcal{O}_L(iE) =
\pi_*\mathcal{O}_L(-i) =
\bigoplus\nolimits_{k \geq -i} \mathcal{O}_P(k)
$$
となる。これは More on Morphisms, Section
\ref{more-morphisms-section-proj-spec} で論じたとおりである。したがって、
$P = \mathbf{P}^{n - 1}_Z$ 上の構造層の捩りのコホモロジーが消滅することから結論を得る。
Cohomology of Schemes, Lemma
\ref{coherent-lemma-cohomology-projective-space-over-ring}.
\end{proof}























\section{爆破とド・ラームコホモロジー}
\label{derham-section-blowing-up}

\noindent
基底スキーム $S$、滑らかな射 $X \to S$、および閉部分スキーム $Z \subset X$
（$S$ 上でも滑らか）を固定する。爆破を $b : X' \to X$ と書く。これは $X$ を $Z$ に沿って行ったものである。
例外因子を $E \subset X'$ と書く。図式
\begin{equation}
\label{derham-equation-blowup}
\vcenter{
\xymatrix{
E \ar[r]_j \ar[d]_p & X' \ar[d]^b \\
Z \ar[r]^i & X
}
}
\end{equation}
この節の目的は、写像
$b^* : H_{dR}^*(X/S) \longrightarrow H_{dR}^*(X'/S)$
が単射であることを証明することである（実際にはさらに多くのことが言える）。

\begin{lemma}
\label{derham-lemma-comparison-general}
More on Morphisms, Lemma \ref{more-morphisms-lemma-blowup} と同じ記号を用いる。
$a \geq 0$ に対して
\begin{enumerate}
\item 写像
$\Omega^a_{X/S} \to b_*\Omega^a_{X'/S}$ は同型である。
\item 写像 $\Omega^a_{Z/S} \to p_*\Omega^a_{E/S}$ は同型である。
\item 写像 $Rb_*\Omega^a_{X'/S} \to i_*Rp_*\Omega^a_{E/S}$ は、次数 $\geq 1$ の
コホモロジー層上で同型である。
\end{enumerate}
\end{lemma}

\begin{proof}
全射エタール射 $\epsilon : X_1 \to X$ を取る。More on Morphisms, Lemma
\ref{more-morphisms-lemma-blowup} で考えた対象の基底変換を
$i_1 : Z_1 \to X_1$, $b_1 : X'_1 \to X_1$, $E_1 \subset X'_1$,
および $p_1 : E_1 \to Z_1$ と書く。$i_1$ は $S$ 上滑らかなスキームの閉埋め込みであり、
Divisors, Lemma \ref{divisors-lemma-flat-base-change-blowing-up} により
$b_1$ は中心 $Z_1$ をもつ爆破であることに注意する。
$b_1$, $p_1$, $i_1$ に対して (1), (2), (3) を証明できると仮定する。このとき
Lemma \ref{derham-lemma-etale} により、$\epsilon$ による (1), (2), (3) の写像の引き戻しは同型である。
$\epsilon$ は全射平坦射なので結論を得る。
したがって、More on Morphisms, Lemma
\ref{more-morphisms-lemma-etale-local-structure} によりエタール局所的に考えれば、
Section \ref{derham-section-calculations} で扱った状況にあると仮定できる。この場合、この補題は
Lemma \ref{derham-lemma-comparison} と同じである。
\end{proof}

\begin{lemma}
\label{derham-lemma-distinguished-triangle-blowup}
More on Morphisms, Lemma \ref{more-morphisms-lemma-blowup} と同じ記号を用い、
構造射を $f : X \to S$ と書く。このとき標準的な識別三角
$$
\Omega^\bullet_{X/S} \to
Rb_*(\Omega^\bullet_{X'/S}) \oplus i_*\Omega^\bullet_{Z/S} \to
i_*Rp_*(\Omega^\bullet_{E/S}) \to
\Omega^\bullet_{X/S}[1]
$$
が $D(X, f^{-1}\mathcal{O}_S)$ に存在する。ここで、4つの写像
$$
\begin{matrix}
\Omega^\bullet_{X/S} & \to & Rb_*(\Omega^\bullet_{X'/S}), \\
\Omega^\bullet_{X/S} & \to & i_*\Omega^\bullet_{Z/S}, \\
Rb_*(\Omega^\bullet_{X'/S}) & \to & i_*Rp_*(\Omega^\bullet_{E/S}), \\
i_*\Omega^\bullet_{Z/S} & \to & i_*Rp_*(\Omega^\bullet_{E/S})
\end{matrix}
$$
は（Section \ref{derham-section-de-rham-complex} の）標準的な写像である。
ただし、そのうち一つは符号を反転したものとする。
\end{lemma}

\begin{proof}
識別三角
$$
C \to Rb_*\Omega^\bullet_{X'/S} \oplus i_*\Omega^\bullet_{Z/S}
\to i_*Rp_*\Omega^\bullet_{E/S} \to C[1]
$$
を $D(X, f^{-1}\mathcal{O}_S)$ において選ぶ。標準写像と両立する形で
$\Omega^\bullet_{X/S}$ が $C$ と同型であることを示せば十分である。三角圏の公理により、
識別三角の写像
$$
\xymatrix{
C' \ar[r] \ar[d] &
b_*\Omega^\bullet_{X'/S} \oplus i_*\Omega^\bullet_{Z/S} \ar[r] \ar[d] &
i_*p_*\Omega^\bullet_{E/S} \ar[r] \ar[d] &
C'[1] \ar[d] \\
C \ar[r] &
Rb_*\Omega^\bullet_{X'/S} \oplus i_*\Omega^\bullet_{Z/S} \ar[r] &
i_*Rp_*\Omega^\bullet_{E/S} \ar[r] &
C[1]
}
$$
が存在する。Lemma \ref{derham-lemma-comparison-general} の (3) および
Derived Categories, Proposition \ref{derived-proposition-9} により、
$C' \to C$ は同型である。Lemma \ref{derham-lemma-comparison-general} の (2) により、
写像 $i_*\Omega^\bullet_{Z/S} \to i_*p_*\Omega^\bullet_{E/S}$ も同型である。
したがって導来圏では $C' = b_*\Omega^\bullet_{X'/S}$ である。最後に
Lemma \ref{derham-lemma-comparison-general} の (1) を用いると、これは
$\Omega^\bullet_{X/S}$ に等しい。標準写像との両立性の検証は省略する。
\end{proof}

\begin{proposition}
\label{derham-proposition-blowup-split}
More on Morphisms, Lemma \ref{more-morphisms-lemma-blowup} と同じ記号を用いると、
写像 $\Omega^\bullet_{X/S} \to Rb_*\Omega^\bullet_{X'/S}$ は
$D(X, (X \to S)^{-1}\mathcal{O}_S)$ において分裂をもつ。
\end{proposition}

\begin{proof}
Lemma \ref{derham-lemma-distinguished-triangle-blowup} で構成した三角を考える。
写像
$$
Rb_*(\Omega^\bullet_{X'/S}) \oplus i_*\Omega^\bullet_{Z/S} \to
i_*Rp_*(\Omega^\bullet_{E/S})
$$
は、その像が直和因子 $i_*\Omega^\bullet_{Z/S}$ を含む分裂をもつと主張する。
Derived Categories, Lemma \ref{derived-lemma-split} により、これは三角の第一の矢印が
直和因子 $i_*\Omega^\bullet_{Z/S}$ 上で消える分裂をもつことを示し、補題が従う。
この主張は $Rp_*\Omega^\bullet_{E/S}$ を直和に分解して証明する。
その第一の部分は $\Omega^\bullet_{Z/S}$ に対応し、第二の部分は
$Rb_*\Omega^\bullet_{X'/S}$ を経て持ち上げられる。

\medskip\noindent
主張を証明する。$X$ を $S$ 上の相対次元が一定の開閉部分スキームに分解できる。
Morphisms, Lemma \ref{morphisms-lemma-smooth-omega-finite-locally-free} を参照せよ。
導来圏 $D(X, f^{-1}\mathcal{O})_S)$ も対応する圏の積に分解するので、
$X$ は相対次元 $N$ が $S$ 上で一定であると仮定してよい。さらに
$Z = \coprod Z_m$ を相対次元 $m \geq 0$ の $S$ 上の開閉部分スキームに分解する。
$i_m : Z_m \to X$ は $i$ の $Z_m$ への制限であり、余次元 $N - m$ の正則埋め込みである。
Divisors, Lemma
\ref{divisors-lemma-immersion-smooth-into-smooth-regular-immersion}.
$E = \coprod E_m$ を対応する分解とし、$E_m = p^{-1}(Z_m)$ とおく。
制限を $p_m : E_m \to Z_m$ と書けば、これは $p$ の $E_m$ への制限であり、標準同型
$$
\tilde \xi_m :
\bigoplus\nolimits_{t = 0, \ldots, N - m - 1}
\Omega^\bullet_{Z_m/S}[-2t]
\longrightarrow
Rp_{m, *}\Omega^\bullet_{E_m/S}
$$
が $D(Z_m, (Z_m \to S)^{-1}\mathcal{O}_S)$ に存在する。次数 $0$ では標準写像
$\Omega^\bullet_{Z_m/S} \to Rp_{m, *}\Omega^\bullet_{E_m/S}$ をとる。
Remark \ref{derham-remark-projective-space-bundle-formula} を参照せよ。
したがって同型
$$
\tilde \xi :
\bigoplus\nolimits_m
\bigoplus\nolimits_{t = 0, \ldots, N - m - 1}
\Omega^\bullet_{Z_m/S}[-2t]
\longrightarrow
Rp_*(\Omega^\bullet_{E/S})
$$
が $D(Z, (Z \to S)^{-1}\mathcal{O}_S)$ に存在する。その始域の直和因子
$\Omega^\bullet_{Z/S} = \bigoplus \Omega^\bullet_{Z_m/S}$ への制限は、標準写像
$\Omega^\bullet_{Z/S} \to Rp_*(\Omega^\bullet_{E/S})$ である。
部分複体 $M_m$ と $K_m$ を、Remark \ref{derham-remark-projective-space-bundle-formula} で導入した複体
$\bigoplus\nolimits_{t = 0, \ldots, N - m - 1} \Omega^\bullet_{Z_m/S}[-2t]$ の中に考える。
次のようにおく。
$$
M = \bigoplus M_m
\quad\text{および}\quad
K = \bigoplus K_m
$$
$M = K[-2]$ であり、構成により写像
$$
c_{E/Z} \oplus \tilde \xi|_M :
\Omega^\bullet_{Z/S} \oplus M
\longrightarrow
Rp_*(\Omega^\bullet_{E/S})
$$
は同型である（上記の注意を参照せよ）。

\medskip\noindent
Lemma \ref{derham-lemma-gysin-via-log-complex} の写像
$$
\delta : \Omega^\bullet_{E/S}[-2] \longrightarrow \Omega^\bullet_{X'/S}
$$
を $D(X', (X' \to S)^{-1}\mathcal{O}_S)$ において考える。
この写像は、次の合成が
$$
\Omega^\bullet_{E/S}[-2] \longrightarrow \Omega^\bullet_{X'/S}
\longrightarrow
\Omega^\bullet_{E/S}
$$
が Remark \ref{derham-remark-cup-product-as-a-map} の $\theta'$ 写像となる性質をもつ。
ここで $c_1^{dR}(\mathcal{O}_{X'}(-E))|_E) = c_1^{dR}(\mathcal{O}_E(1))$ である。
Remark \ref{derham-remark-projective-space-bundle-formula} の最後の主張により、図式
$$
\xymatrix{
K[-2] \ar[d]_{(\tilde \xi|_K)[-2]} \ar[r]_{\text{id}} &
M \ar[d]^{\tilde x|_M} \\
Rp_*\Omega^\bullet_{E/S}[-2] \ar[r]^-{Rp_*\theta'} &
Rp_*\Omega^\bullet_{E/S}
}
$$
は可換である。したがって、主張の望ましい分裂は次の写像として得られる。
\begin{align*}
Rp_*(\Omega^\bullet_{E/S})
& \xrightarrow{(c_{E/Z} \oplus \tilde \xi|_M)^{-1}}
\Omega^\bullet_{Z/S} \oplus M \\
& \xrightarrow{\text{id} \oplus \text{id}^{-1}}
\Omega^\bullet_{Z/S} \oplus K[-2] \\
& \xrightarrow{\text{id} \oplus (\tilde \xi|_K)[-2]}
\Omega^\bullet_{Z/S} \oplus Rp_*\Omega^\bullet_{E/S}[-2] \\
& \xrightarrow{\text{id} \oplus Rb_*\delta}
\Omega^\bullet_{Z/S} \oplus Rb_*\Omega^\bullet_{X'/S}
\end{align*}
上で述べた $\theta'$ と $\delta$ の関係、および $\theta'$、$\tilde \xi|_K$、
$\tilde \xi|_M$ を含む上の可換図式により、これが標準写像
$\Omega^\bullet_{Z/S} \oplus Rb_*(\Omega^\bullet_{X'/S}) \to
Rp_*(\Omega^\bullet_{E/S})$ の切断であることが分かる。これで主張の証明は完了する。
\end{proof}

\noindent
Lemma \ref{derham-lemma-splitting-on-omega-a}
は、Hodge コホモロジー上で分裂を構成する方が Proposition
\ref{derham-proposition-blowup-split} の結果よりもはるかに容易であることを示す。
読者には次の節へ先に進むことを勧める。

\begin{lemma}
\label{derham-lemma-ext-zero}
$i : Z \to X$ を余次元 $c$ の正則なスキームの閉埋め込みとする。このとき
$\Ext^q_{\mathcal{O}_X}(i_*\mathcal{F}, \mathcal{E}) = 0$ は
$q < c$ に対して、$\mathcal{E}$ が $X$ 上局所自由で、$\mathcal{F}$ が任意の
$\mathcal{O}_Z$-加群ならば成り立つ。
\end{lemma}

\begin{proof}
局所から大域への $\Ext$ のスペクトル系列により、$X$ 上アフィン局所に
これを示せば十分である。Cohomology, Section \ref{cohomology-section-ext} を参照せよ。
したがって $X = \Spec(A)$ と仮定でき、$f_1, \ldots, f_c$ という正則列が $A$ に存在し、
$Z = \Spec(A/(f_1, \ldots, f_c))$ となると仮定できる。さらに $c \geq 1$ としてよい。
すると $f_1 : \mathcal{E} \to \mathcal{E}$ は単射である。$i_*\mathcal{F}$ は $f_1$ で消えるので、
これは $i = 0$ のとき補題が成り立つことを示し、全射
$$
\Ext^{q - 1}_{\mathcal{O}_X}(i_*\mathcal{F}, \mathcal{E}/f_1\mathcal{E})
\longrightarrow
\Ext^q_{\mathcal{O}_X}(i_*\mathcal{F}, \mathcal{E})
$$
を得る。この矢印の始域が零であることを示せば十分である。
次に同じ議論を繰り返す。$c \geq 2$ ならば写像
$f_2 : \mathcal{E}/f_1\mathcal{E} \to \mathcal{E}/f_1\mathcal{E}$
は単射である。$i_*\mathcal{F}$ は $f_2$ で消えるので、これは $q = 1$ のとき補題が成り立つことを示し、
全射
$$
\Ext^{q - 2}_{\mathcal{O}_X}(i_*\mathcal{F},
\mathcal{E}/f_1\mathcal{E} + f_2\mathcal{E})
\longrightarrow
\Ext^{q - 1}_{\mathcal{O}_X}(i_*\mathcal{F}, \mathcal{E}/f_1\mathcal{E})
$$
を得る。この手続きを続ければ補題が証明される。
\end{proof}

\begin{lemma}
\label{derham-lemma-splitting-on-omega-a}
More on Morphisms, Lemma \ref{more-morphisms-lemma-blowup} と同じ記号を用いる。
$a \geq 0$ に対して、一意な射
$Rb_*\Omega^a_{X'/S} \to \Omega^a_{X/S}$ が $D(\mathcal{O}_X)$ に存在し、
$\Omega^a_{X/S} \to Rb_*\Omega^a_{X'/S}$ との合成は
$\Omega^a_{X/S}$ 上の恒等射である。
\end{lemma}

\begin{proof}
まず $X$ を $S$ 上の相対次元が一定の開閉部分スキームに分解できる。
Morphisms, Lemma \ref{morphisms-lemma-smooth-omega-finite-locally-free} を参照せよ。
導来圏 $D(X, f^{-1}\mathcal{O})_S)$ も対応する圏の積に分解するので、
$X$ は相対次元 $N$ が $S$ 上で一定であると仮定してよい。さらに
$Z = \coprod Z_m$ を相対次元 $m \geq 0$ の $S$ 上の開閉部分スキームに分解する。
$i_m : Z_m \to X$ は $i$ の $Z_m$ への制限であり、余次元 $N - m$ の正則埋め込みである。
Divisors, Lemma
\ref{divisors-lemma-immersion-smooth-into-smooth-regular-immersion}.
$E = \coprod E_m$ を対応する分解とし、$E_m = p^{-1}(Z_m)$ とおく。
次の自然な写像が存在すると主張する。
$$
b^*\Omega^a_{X/S} \to \Omega^a_{X'/S} \to
b^*\Omega^a_{X/S} \otimes_{\mathcal{O}_{X'}}
\mathcal{O}_{X'}(\sum (N - m - 1)E_m)
$$
その合成は包含
$\mathcal{O}_{X'} \to \mathcal{O}_{X'}(\sum (N - m - 1)E_m)$ によって誘導される。
これを示すには、第一の矢印の余核が $X'$ 上局所的に有効 Cartier 因子
$\sum (N - m - 1)E_m$ の局所方程式で消えることを示せば十分である。
このため、Lemma \ref{derham-lemma-comparison-general} の証明と同様に $X$ 上エタール局所的に考え、
Lemma \ref{derham-lemma-comparison-bis} を適用すればよい。
Lemma \ref{derham-lemma-blowup-twist-same-cohomology} を用いてエタール局所的に計算すると、
誘導された合成
$$
\Omega^a_{X/S} \to Rb_*\Omega^a_{X'/S} \to
Rb_*\left(b^*\Omega^a_{X/S} \otimes_{\mathcal{O}_{X'}}
\mathcal{O}_{X'}(\sum (N - m - 1)E_m)\right)
$$
は $D(\mathcal{O}_X)$ における同型である。これにより、この補題の写像の存在が得られる。

\medskip\noindent
一意性を示すには、$\tau_{\geq 1}Rb_*\Omega_{X'/S}$ から
$\Omega^a_{X/S}$ への $D(\mathcal{O}_X)$ における非零写像が存在しないことを示せば十分である。
そのためには、$R^qb_*\Omega_{X'/s}[-q]$ から $\Omega^a_{X/S}$ への
$D(\mathcal{O}_X)$ における非零写像が $q \geq 1$ に対して存在しないことを示せばよい
（細部は省略する）。Lemma \ref{derham-lemma-comparison-general} により、
$R^qb_*\Omega_{X'/s} \cong i_*R^qp_*\Omega^a_{E/S}$ は
$Z = \coprod Z_m$ 上の加群の押し出しである。
さらに、$R^qp_*\Omega^a_{E/S}$ の $Z_m$ への制限は $q < N - m$ のときに限り非零であることに注意する。
実際、$E_m \to Z_m$ のファイバーの次元は $N - m - 1$ なので、Limits, Lemma
\ref{limits-lemma-higher-direct-images-zero-above-dimension-fibre} を適用できる。
したがって、望ましい消滅は Lemma \ref{derham-lemma-ext-zero} から従う。
\end{proof}














\section{微分形式の層の比較}
\label{derham-section-quasi-finite-syntomic}

\noindent
この節の目的は、すべての局所準有限シントミックなスキーム射
$f : Y \to X$ に対して写像
$$
c^p_{Y/X} :
\Omega^p_{Y/\mathbf{Z}}
\longrightarrow
f^*\Omega^p_{X/\mathbf{Z}} \otimes_{\mathcal{O}_Y} \det(\NL_{Y/X})
$$
を、次の性質を満たすすべての $p \geq 0$ について構成することである。
\begin{enumerate}
\item
\label{derham-item-degree-zero}
$c_{Y/X}^0(1) = \delta(\NL_{Y/X})$ である。Discriminants, Section
\ref{discriminant-section-tate-map} を参照せよ。
\item
\label{derham-item-multiplicative}
$c_{Y/X}^{q + p}(\omega \wedge \eta) = \omega \wedge c_{Y/X}^p(\eta)$
は、$\omega$ が $f^*\Omega^q_{X/\mathbf{Z}}$ の局所切断であり、
$\eta$ が $\Omega^p_{Y/\mathbf{Z}}$ の局所切断であるときに成り立つ。
\item
\label{derham-item-base-change}
$c^p_{Y/X}$ の構成は開集合への制限および基底変換と両立する
（Lemma \ref{derham-lemma-base-change-Garel-upstairs} で定式化される形で）。
\end{enumerate}
これらの条件から、写像 $c^p_{Y/X}$ は
$\mathcal{O}_Y$-線形であり、さらに
$f^*\Omega^p_{X/\mathbf{Z}} \to \Omega^p_{Y/\mathbf{Z}}$
との合成は $\delta(\NL_{Y/X})$ による乗法であることに注意せよ。
実際、性質 (1), (2), (3) をもつ作用素 $c^p_{Y/X}$ の一意な族が存在することを示す。
この結果は、$f$ が有限射である場合にド・ラーム複体上のトレース写像を構成するため、
Section \ref{derham-section-trace} で適用される。

\begin{lemma}
\label{derham-lemma-funny-map}
$R$ を環とし、次の可換図式を考える。
$$
\xymatrix{
0 \ar[r] &
K^0 \ar[r] &
L^0 \ar[r] &
M^0 \ar[r] & 0 \\
& & L^{-1} \ar[u]_\partial \ar@{=}[r] &
M^{-1} \ar[u]
}
$$
これを $R$-加群の図式とし、上段は完全で、$M^0$ と $M^{-1}$ は
同じランクの有限自由加群であるとする。$p \geq 0$ に対して標準写像
$$
c^p : 
\wedge^p(H^0(L^\bullet))
\longrightarrow
\wedge^p(K^0) \otimes_R \det(M^\bullet)
$$
が存在し、次の性質を満たす。
\begin{enumerate}
\item $c^0(1) = \delta(M^\bullet)$ であり、
\item $c^{q + p}(\omega \wedge \eta) = \omega \wedge c^p(\eta)$
 $\omega \in \wedge^q(K^0)$ および $\eta \in \wedge^p(H^0(L^\bullet))$ に対して成り立つ。
\end{enumerate}
特に、$c^p$ と写像
$\wedge^p(K^0) \to \wedge^p(H^0(L^\bullet))$ との合成は
$\delta(M^\bullet)$ による乗法である。
\end{lemma}

\begin{proof}
$M^0$ と $M^{-1}$ がランク $n$ の自由加群であるとする。各 $p \geq 0$
に対して一意な全射
$$
\pi_p :
\wedge^{p + n}(L^0)
\longrightarrow
\wedge^p(K^0) \otimes \wedge^n(M^0)
$$
が存在し、次の二つの性質をもつ。(i)
$
\pi_p(k_1 \wedge \ldots \wedge k_p \wedge l_1 \wedge \ldots \wedge l_n) =
k_1 \wedge \ldots \wedge k_p \otimes
m_1 \wedge \ldots \wedge m_n
$
が、$k_1, \ldots, k_p \in K^0$ であり、$l_1, \ldots, l_n \in L^0$ が
$m_1, \ldots, m_n \in M^0$ に写るときに成り立つ。(ii) $\pi_p$ の核は、
外積 $l_1 \wedge \ldots \wedge l_{p + n}$ であって、そのうち
$> p$ 個の $l_j$ が $K^0$ に属するものによって生成される部分加群である。
基底 $e_1, \ldots, e_n$ を $L^{-1} = M^{-1}$ に対して選ぶ。
そこで、写像を次の規則で定める。
$$
\eta \mapsto
\pi_p(\tilde \eta \wedge \partial(e_1) \wedge \ldots \wedge \partial(e_n))
\otimes (e_1 \wedge \ldots \wedge e_n)^{\otimes -1}
$$
$\eta \in \wedge^p(H^0(L^\bullet))$ とし、$\tilde \eta \in \wedge^p(L^0)$ は
$\eta$ の持ち上げとする。右辺が意味をもつのは、
$\det(M^\bullet) = \wedge^n(M^0) \otimes \wedge^n(M^{-1})^{\otimes -1}$.
だからである。右辺の式は $\tilde \eta$ の選び方に依存しない。実際、
$\partial(\xi) \wedge \partial(e_1) \wedge \ldots \wedge \partial(e_n) = 0$
が任意の $\xi \in L^{-1}$ に対して成り立つ。直接計算すれば、この写像は
$L^{-1} = M^{-1}$ の基底の選び方にも依存しないことが分かる。
More on Algebra, Section \ref{more-algebra-section-determinants-complexes}
における $\delta(M^\bullet)$ の定義から $c^0(1) = \delta(M^\bullet)$ は直ちに従う。
最後に、主張 (2) の規則は定義から直接確認できる。
\end{proof}

\begin{lemma}
\label{derham-lemma-funny-map-functorial}
$R_1 \to R_2$ を環準同型とする。$i = 1, 2$ に対して可換図式
$$
\xymatrix{
0 \ar[r] &
K_i^0 \ar[r] &
L_i^0 \ar[r] &
M_i^0 \ar[r] & 0 \\
& & L_i^{-1} \ar[u]_{\partial_i} \ar@{=}[r] &
M_i^{-1} \ar[u]
}
$$
が Lemma \ref{derham-lemma-funny-map} と同様に $R_i$-加群からなるものとする。
写像 $K_1^0 \to K_2^0$, $L_1^j \to L_2^j$, $M_1^j \to M_2^j$ が与えられ、
所与の環準同型 $R_1 \to R_2$ および表示した図式内の写像と両立すると仮定する。
写像 $M_1^j \otimes_{R_1} R_2 \to M_2^j$ が同型ならば、図式
$$
\xymatrix{ 
\wedge^p(H^0(L_1^\bullet)) \ar[d]
\ar[r]_-{c^p} &
\wedge^p(K_1^0) \otimes_{R_1} \det(M_1^\bullet) \ar[d] \\
\wedge^p(H^0(L_2^\bullet))
\ar[r]^-{c^p} &
\wedge^p(K_2^0) \otimes_{R_2} \det(M_2^\bullet)
}
$$
は可換である。
\end{lemma}

\begin{proof}
これは Lemma \ref{derham-lemma-funny-map} の証明で与えた写像の明示的記述から従う。
$M_1^j \otimes_{R_1} R_2 \to M_2^j$ が同型であるという仮定は、
$M_1^{-1}$ に用いた基底が $M_2^{-1}$ の基底に写ることを保証するために必要である
（したがって、二つの構成の両立性を確認する際に「同じ」基底を使える）。
\end{proof}

\begin{remark}
\label{derham-remark-local-description}
$A$ を環とする。$P = A[x_1, \ldots, x_n]$ とおく。
$f_1, \ldots, f_n \in P$ を取り、$B = P/(f_1, \ldots, f_n)$ とおく。
$A \to B$ が準有限であると仮定する。このとき $B$ は $A$ 上の相対大域完全交叉であり
（Algebra, Definition \ref{algebra-definition-relative-global-complete-intersection}）、
Algebra, Lemma \ref{algebra-lemma-relative-global-complete-intersection-conormal}
により $(f_1, \ldots, f_n)/(f_1, \ldots, f_n)^2$ は類
$\overline{f}_i$ を生成元とする自由加群である。次の図式を考える。
$$
\xymatrix{
\Omega_{A/\mathbf{Z}} \otimes_A B \ar[r] &
\Omega_{P/\mathbf{Z}} \otimes_P B \ar[r] &
\Omega_{P/A} \otimes_P B \\
&
(f_1, \ldots, f_n)/(f_1, \ldots, f_n)^2 \ar[u] \ar@{=}[r] &
(f_1, \ldots, f_n)/(f_1, \ldots, f_n)^2 \ar[u]
}
$$
右列は $\NL_{B/A}$ を $D(B)$ において表すので、コホモロジーは
$\Omega_{B/A}$ を次数 $0$ にもつ。上段は分裂短完全列
$0 \to \Omega_{A/\mathbf{Z}} \otimes_A B \to
\Omega_{P/\mathbf{Z}} \otimes_P B \to \Omega_{P/A} \otimes_P B \to 0$.
である。Algebra, Lemma \ref{algebra-lemma-differential-seq} により、
中列のコホモロジーは $\Omega_{B/\mathbf{Z}}$ を次数 $0$ にもつ。
したがって Lemma \ref{derham-lemma-funny-map} により標準的な $B$-加群写像
$$
\Omega^p_{B/\mathbf{Z}} \longrightarrow
\Omega^p_{A/\mathbf{Z}} \otimes_A \det(\NL_{B/A})
$$
を $p \geq 0$ に対して得る。これは性質 (1) と (2) を満たし、特に
$\Omega^p_{A/\mathbf{Z}} \to \Omega^p_{B/\mathbf{Z}}$
との合成が $\delta(\NL_{B/A})$ による乗法であるという性質を満たす。
\end{remark}

\begin{lemma}
\label{derham-lemma-base-change-Garel-upstairs}
可換図式
$$
\xymatrix{
Y' \ar[d]_{f'} \ar[r]_b & Y \ar[d]^f \\
X' \ar[r]^a & X
}
$$
を考え、これが $Y'$ と $X' \times_X Y$ のある開部分スキームとの同型を誘導するとする。
$f$ は局所準有限かつシントミックであり、写像
$c^p_{Y/X} : \Omega^p_{Y/\mathbf{Z}} \to
f^*\Omega^p_{X/\mathbf{Z}} \otimes_{\mathcal{O}_Y} \det(\NL_{Y/X})$
が $p \geq 0$ に対して与えられ、
(\ref{derham-item-degree-zero}) と (\ref{derham-item-multiplicative}) を満たすと仮定する。
このとき、写像の族
$c^p_{Y'/X'} : \Omega^p_{Y'/\mathbf{Z}} \to
(f')^*\Omega^p_{X'/\mathbf{Z}} \otimes_{\mathcal{O}_{Y'}} \det(\NL_{Y'/X'})$
であって、$p \geq 0$ に対して (\ref{derham-item-degree-zero}) と
(\ref{derham-item-multiplicative}) を満たし、さらに図式
$$
\xymatrix{
b^*\Omega^p_{Y/\mathbf{Z}} \ar[rr]_-{b^*c^p_{Y/X}} \ar[d] & &
b^*(f^*\Omega^p_{X/\mathbf{Z}} \otimes_{\mathcal{O}_Y} \det(\NL_{Y/X}))
\ar@{=}[r] &
(f')^*a^*\Omega^p_{X/\mathbf{Z}} \otimes_{\mathcal{O}_Y} b^*\det(\NL_{Y/X}))
\ar[d] \\
\Omega^p_{Y'/\mathbf{Z}} \ar[rrr]^-{c^p_{Y'/X'}} & & &
(f')^*\Omega^p_{X'/\mathbf{Z}} \otimes_{\mathcal{O}_{Y'}} \det(\NL_{Y'/X'})
}
$$
がすべての $p \geq 0$ に対して可換となるものは高々一つである。ここで縦の矢印には、写像
$b^*\Omega^p_{Y/\mathbf{Z}} \to \Omega^p_{Y'/\mathbf{Z}}$ および
$a^*\Omega^p_{X/\mathbf{Z}} \to \Omega^p_{X'/\mathbf{Z}}$
すなわち Section \ref{derham-section-de-rham-complex} の写像と、
Discriminants, Section \ref{discriminant-section-tate-map} の同一視
$b^*\det(\NL_{Y/X}) = \det(\NL_{Y'/X'})$ を用いる。
\end{lemma}

\begin{proof}
写像
$$
(f')^*\Omega_{X'/\mathbf{Z}} \oplus b^*\Omega_{Y/\mathbf{Z}}
\longrightarrow
\Omega_{Y'/\mathbf{Z}}
$$
は全射である。実際、$Y'$ はファイバー積の開部分スキームであり、対応する代数的主張は、
$\Omega_{B \otimes_A C/\mathbf{Z}}$ が $\Omega_{B/\mathbf{Z}}$ と
$\Omega_{C/\mathbf{Z}}$ の像によって加群として生成されるということである。
したがってすべての $p$ に対して写像
$$
\bigoplus\nolimits_{i = 0, \ldots, p}
(f')^*\Omega^i_{X'/\mathbf{Z}} \otimes
b^*\Omega^{p - i}_{Y/\mathbf{Z}}
\longrightarrow
\Omega^p_{Y'/\mathbf{Z}}
$$
は全射である。条件 (\ref{derham-item-degree-zero}) と (\ref{derham-item-multiplicative}) を
補題中の図式の可換性と合わせると、写像 $c^p_{Y'/X'}$ は一意に定まる
（ただし存在は直ちには従わない）。
\end{proof}

\begin{lemma}
\label{derham-lemma-existence-base-change-Garel-upstairs}
可換図式
$$
\xymatrix{
Y' \ar[d]_{f'} \ar[r]_b & Y \ar[d]^f \\
X' \ar[r]^a & X
}
$$
を考え、これが $Y'$ と $X' \times_X Y$ のある開部分スキームとの同型を誘導するとする。
$f$ は局所準有限かつシントミックであると仮定する。このとき、各 $y' \in Y'$ に対して
開集合 $V' \subset Y'$, $V \subset Y$, $U' \subset X$, $U \subset X$ を、
$y' \in V'$, $f'(V') \subset U'$, $b(V') \subset V$, $a(U') \subset U$,
$f(V) \subset U$ となるように選べる。さらに $p \geq 0$ に対して写像
$c^p_{V/U}$ と $c^p_{V'/U'}$ が存在し、(\ref{derham-item-degree-zero}) と
(\ref{derham-item-multiplicative}) を満たし、図式
$$
\xymatrix{
V' \ar[d] \ar[r]_b & V \ar[d] \\
U' \ar[r] & U
}
$$
と Lemma \ref{derham-lemma-base-change-Garel-upstairs} で説明した意味で両立する。
\end{lemma}

\begin{proof}
証明のため $Y'$ を $X' \times_X Y$ で置き換えてよいことは明らかである。
アフィン開集合 $V \subset Y$ と $U \subset X$ を選び、$V$ が $y'$ の像を含むようにする。
これは Discriminants, Lemma \ref{discriminant-lemma-syntomic-quasi-finite} (5) による。
アフィン開集合 $U' \subset X'$ を選び、それが $y'$ の像を含み $U$ に写るようにする。
$X, Y, X', Y'$ を $U, V, U', U' \times_U V$ で置き換えれば、
次の段落で述べる状況にあると仮定してよい。

\medskip\noindent
$X' = \Spec(A')$, $X = \Spec(A)$, $Y = \Spec(B)$, および
$Y' = \Spec(A' \otimes_A B)$ とし、$B = A[x_1, \ldots, x_n]/(f_1, \ldots, f_n)$
は $A$ 上の相対大域完全交叉であると仮定する。このとき
$B' = A'[x_1, \ldots, x_n]/(f'_1, \ldots, f'_n)$ も $A'$ 上の相対大域完全交叉である。
ここで $f'_i$ は $f_i$ の $A'[x_1, \ldots, x_n]$ における像である。
Algebra, Lemma \ref{algebra-lemma-base-change-relative-global-complete-intersection} を参照せよ。
Remark \ref{derham-remark-local-description} の構成から写像 $c^p_{V/U}$ と $c^p_{V'/U'}$ を得る。
これらの写像は、Lemma \ref{derham-lemma-funny-map-functorial} と、$B$ および $B'$ の
$A$ および $A'$ 上の表示の両立性により両立する。細部は省略する。
\end{proof}

\begin{lemma}
\label{derham-lemma-Garel-upstairs}
すべての局所準有限シントミックなスキーム射 $f : Y \to X$ に
$\mathcal{O}_Y$-加群写像
$$
c^p_{Y/X} :
\Omega^p_{Y/\mathbf{Z}}
\longrightarrow
f^*\Omega^p_{X/\mathbf{Z}} \otimes_{\mathcal{O}_Y} \det(\NL_{Y/X})
$$
を $p \geq 0$ に対して対応させ、(\ref{derham-item-degree-zero}),
(\ref{derham-item-multiplicative}), (\ref{derham-item-base-change}) を満たす規則が一意に存在する。
特に、$c^p_{Y/X}$ と $f^*\Omega^p_{X/\mathbf{Z}} \to \Omega^p_{Y/\mathbf{Z}}$
との合成は $\delta(\NL_{Y/X})$ による乗法である。
\end{lemma}

\begin{proof}
この証明は Discriminants, Proposition \ref{discriminant-proposition-tate-map}
の証明と非常によく似ているので、まずそちらの証明を見ることを勧める。

\medskip\noindent
主張を言い換えよう。圏 $\mathcal{C}$ を考える。その対象は $Y/X$ と書き、
局所準有限シントミックなスキーム射 $f : Y \to X$ である。また射
$b/a : Y'/X' \to Y/X$ は可換図式
$$
\xymatrix{
Y' \ar[d]_{f'} \ar[r]_b & Y \ar[d]^f \\
X' \ar[r]^a & X
}
$$
であって、$Y'$ と $X' \times_X Y$ のある開部分スキームとの同型を誘導するものとする。
補題の意味は、各対象 $Y/X$ が $\mathcal{C}$ に属するとき写像 $c^p_{Y/X}$,
$p \geq 0$ があり、性質 (\ref{derham-item-degree-zero}) と
(\ref{derham-item-multiplicative}) を満たし、さらに各射
$b/a : Y'/X' \to Y/X$ が $\mathcal{C}$ に属するとき $c^p_{Y/X}$ と $c^p_{Y'/X'}$ が
Lemma \ref{derham-lemma-base-change-Garel-upstairs} の意味で両立するということである。

\medskip\noindent
$Y/X$ を $\mathcal{C}$ から取り、$y \in Y$ が与えられたとする。
アフィン開集合 $V \subset Y$ と $U \subset X$ を $f(V) \subset U$ となるように選び、
写像
$$
\Omega^p_{Y/\mathbf{Z}}|_V
\longrightarrow
\left(
f^*\Omega^p_{X/\mathbf{Z}} \otimes_{\mathcal{O}_Y} \det(\NL_{Y/X})
\right)|_V
$$
で性質 (\ref{derham-item-degree-zero}) と (\ref{derham-item-multiplicative}) を満たすものを得られる。
これは Discriminants, Lemma \ref{discriminant-lemma-syntomic-quasi-finite} (5)
のようにアフィン開集合を選び、Remark \ref{derham-remark-local-description} の構成を用いれば従う。

\medskip\noindent
Discriminants, Section \ref{discriminant-section-tate-map} の議論を参照すると、
$f$ のエタール軌跡は $\delta(\NL_{Y/X})$ が消えない点の集合にちょうど等しい。
特に、
$f^*\Omega^p_{X/\mathbf{Z}} \to \Omega^p_{Y/\mathbf{Z}}$
は $\delta(\NL_{Y/X})$ を逆にすると同型になる。
したがって、$\Omega^p_{X/\mathbf{Z}}$ が有限局所自由であり、
$\delta(\NL_{Y/X})$ の $\mathcal{O}_Y$ における零化イデアルが零ならば、
前段落で構成した局所写像は一意であり、自動的に貼り合わさる。

\medskip\noindent
$\mathcal{C}_{nice} \subset \mathcal{C}$ を、次を満たす $Y/X$ からなる充満部分圏とする。
\begin{enumerate}
\item $\Omega_{X/\mathbf{Z}}$ は局所自由であり、
\item $\delta(\NL_{Y/X})$ の $\mathcal{O}_Y$ における零化イデアルは零である。
\end{enumerate}
前段落の考察から、任意の対象 $Y/X$ が $\mathcal{C}_{nice}$ に属するとき、
条件 (\ref{derham-item-degree-zero}) と (\ref{derham-item-multiplicative}) を満たす一意な写像
$c^p_{Y/X}$, $p \geq 0$ がある。$b/a : Y'/X' \to Y/X$ が
$\mathcal{C}_{nice}$ の射ならば、写像 $c^p_{Y/X}$ と $c^p_{Y'/X'}$ は
Lemma \ref{derham-lemma-base-change-Garel-upstairs} の意味で両立する。
実際、Lemma \ref{derham-lemma-existence-base-change-Garel-upstairs} により局所的には両立する写像が存在し、
先ほどの一意性により、それらは $c^p_{Y/X}$ および $c^p_{Y'/X'}$ と一致する。
言い換えれば、充満部分圏 $\mathcal{C}_{nice}$ 上では問題を解いたことになる。
$Y/X$ が $\mathcal{C}_{nice}$ に属するとき、得られた解を引き続き $c^p_{Y/X}$ と書く。

\medskip\noindent
実際、より一般に、射 $b/a : Y'/X' \to Y/X$ が $\mathcal{C}$ にあり、
$Y/X$ は $\mathcal{C}_{nice}$ の対象であるとする。
同じ議論で、今度は Lemma \ref{derham-lemma-base-change-Garel-upstairs} の一意性を用いると、
写像 $c^p_{Y'/X'}$, $p \geq 0$ が存在し、条件 (\ref{derham-item-degree-zero}) と
(\ref{derham-item-multiplicative}) を満たし、すでに構成した写像 $c^p_{Y/X}$ と両立する。
これを基底変換 $(b/a)^*c^p_{Y/X}$ と呼ぶ。

\medskip\noindent
射
$$
Y_1/X_1 \xleftarrow{b_1/a_1} Y/X \xrightarrow{b_2/a_2} Y_2/X_2
$$
を $\mathcal{C}$ において考え、$Y_1/X_1$ と $Y_2/X_2$ は
$\mathcal{C}_{nice}$ の対象であるとする。
{\bf 主張。} 二つの基底変換 $(b_i/a_i)^*c^p_{Y_i/X_i}$, $i = 1, 2$ は等しい。
まずこの主張から補題が従うことを示し、その後で主張を証明する。

\medskip\noindent
$d, n \geq 1$ とし、Discriminants, Example
\ref{discriminant-example-universal-quasi-finite-syntomic} で構成された局所準有限シントミック射
$Y_{n, d} \to X_{n, d}$ を考える。このとき $Y_{n, d}$ と $Y_{n, d}$ は
$\mathbf{Z}$ 上有限型かつ滑らかな既約スキームである。実際、$X_{n, d}$ は
$\mathbf{Z}$ 上の多項式環のスペクトルであり、$Y_{n, d}$ はその開部分スキームである。
射 $Y_{n, d} \to X_{n, d}$ は局所準有限シントミックであり、稠密開集合上エタールである。
Discriminants, Lemma \ref{discriminant-lemma-universal-quasi-finite-syntomic-etale} を参照せよ。
したがって $\delta(\NL_{Y_{n, d}/X_{n, d}})$ は零でない。例えば、
Discriminants, Remark \ref{discriminant-remark-local-description-delta} には
$\delta(\NL_{Y/X})$ の局所記述があり、Morphisms, Lemma
\ref{morphisms-lemma-etale-at-point} (8) にはエタール射の局所記述がある。
既約正則スキーム上の可逆加群の零でない切断は、零化イデアルが零である。
したがって $Y_{n, d}/X_{n, d}$ は $\mathcal{C}_{nice}$ の対象である。

\medskip\noindent
$Y/X$ を $\mathcal{C}$ の任意の対象とし、$y \in Y$ とする。
Discriminants, Lemma \ref{discriminant-lemma-locally-comes-from-universal} により、
$n, d \geq 1$ と射
$$
Y/X \leftarrow V/U \xrightarrow{b/a} Y_{n, d}/X_{n, d}
$$
を $\mathcal{C}$ の中に取れ、$V \subset Y$ と $U \subset X$ は開集合である。
このとき基底変換 $c^p_{V/U} = (b/a)^*c^p_{Y_{n, d}/X_{n, d}}$ がある。
主張により、これらの局所的に構成された写像 $c^p_{V/U}$ は貼り合わさる。
こうして、性質 (\ref{derham-item-degree-zero}) と (\ref{derham-item-multiplicative}) をもつ
整合的な大域写像 $c^p_{Y/X}$ を得る。最後に、$b/a : Y'/X' \to Y/X$ を
$\mathcal{C}$ の任意の射とする。この段落で構成した写像
$c^p_{Y'/X'}$, $p \geq 0$ と $c^p_{Y/X}$, $p \geq 0$ がある。
これらが Lemma \ref{derham-lemma-base-change-Garel-upstairs} の意味で両立することを確認するには、
$Y'/X'$ と $Y/X$ 上で局所的に議論してよい。したがって射
$Y/X \to Y_{n, d}/X_{n, d}$ が存在すると仮定してよい。構成により、族
$c^p_{Y'/X'}$ と族 $c^p_{X/Y}$ はいずれも $Y_{n, d}/X_{n, d}$ への写像と両立する。
このことと Lemma \ref{derham-lemma-base-change-Garel-upstairs} の一意性から、
$c^p_{Y'/X'}$ は $c^p_{Y/X}$ と両立することが容易に従う。
ゆえに残るのは主張の証明である。

\medskip\noindent
証明の残りでは主張を示す。点 $y \in Y$ を取り、二つの写像が $y$ の
ある開近傍上で一致することを示せばよい。したがって $Y_1$, $Y_2$ を、
$y$ の像の $Y_1$ および $Y_2$ における開近傍で置き換えてよい。
ゆえに $Y, X, Y_1, X_1, Y_2, X_2$ はアフィンであり、それぞれ環
$B, A, B_1, A_1, B_2, A_2$ のスペクトルであると仮定してよい。図式は
$$
\xymatrix{
B_1 \ar[r] & B & B_2 \ar[l] \\
A_1 \ar[u] \ar[r] & A \ar[u] & A_2 \ar[l] \ar[u]
}
$$
である。仮定により、$B$ のスペクトルは $A \otimes_{A_1} B_1$ と
$A \otimes_{A_2} B_2$ の両方のスペクトルのアフィン開集合である。
さらに縮小して、元 $g_i \in A \otimes_{A_i} B_i$ が存在し、所与の写像が同型
$(A \otimes_{A_i} B_i)_{g_i} = B$ を与えると仮定してよい。
Properties, Lemma \ref{properties-lemma-standard-open-two-affines} を参照せよ。
十分大きな変数の族 $x_\alpha$, $y_\beta$ を取り、全射
$$
A'_1 = A_1[x_\alpha] \to A
\quad\text{および}\quad
A'_2 = A_2[y_\beta] \to A
$$
を $A_1$-代数および $A_2$-代数の全射として選べるようにする。
次に、持ち上げ $h_i \in A'_i \otimes_{A_i} B_i$ を、
$g_i \in A \otimes_{A_i} B_i$ に対して選び、図式
$$
\xymatrix{
(A'_1 \otimes_{A_1} B_1)_{h_1} \ar[r] & B &
(A'_2 \otimes_{A_1} B_2)_{h_2} \ar[l] \\
A'_1 \ar[u] \ar[r] & A \ar[u] & A'_2 \ar[l] \ar[u]
}
$$
を考える。構成により二つの正方形は余カルテシアンである。次に環準同型
$$
A' = A'_1 \times_A A'_2 \longrightarrow
B' =
(A'_1 \otimes_{A_1} B_1)_{h_1} \times_B
(A'_2 \otimes_{A_1} B_2)_{h_2}
$$
を考える。More on Algebra, Lemma
\ref{more-algebra-lemma-module-over-fibre-product} により
$A'_1 \otimes_{A'} B' = (A'_1 \otimes_{A_1} B_1)_{h_1}$ および
$A'_2 \otimes_{A'} B' = (A'_2 \otimes_{A_2} B_2)_{h_2}$ が成り立つ。
特に、射 $Y' = \Spec(B') \to \Spec(A') = X'$ のファイバーは、
写像 $Y_i \to X_i$ のファイバーの基底変換の開部分スキームであり、
したがって有限な局所完全交叉である
（Discriminants, Lemma \ref{discriminant-lemma-syntomic-quasi-finite} を参照）。
More on Algebra, Lemma
\ref{more-algebra-lemma-properties-algebras-over-fibre-product} により、
環準同型 $A' \to B'$ は平坦かつ有限表示である。ゆえに Discriminants, Lemma
\ref{discriminant-lemma-syntomic-quasi-finite} (3) により、
$Y'/X'$ は $\mathcal{C}$ の対象であると結論する。ここで可換図式
$$
\xymatrix{
& & Y/X \ar[ld] \ar@/_2pc/[lldd]_{b_1/a_1} \ar[rd] \ar@/^2pc/[rrdd]^{b_2/a_2} \\
& Y'_1/X'_1 \ar[rd] \ar[ld] & & Y'_2/X'_2 \ar[ld] \ar[rd] \\
Y_1/X_1 & & Y'/X' & & Y_2/X_2
}
$$
を考える。ここで $Y'_i/X'_i$ は
$A'_i \to (A'_i \otimes_{A_i} B_i)_{h_i}$ に対応する。
$Y'_i/X'_i$ は $\mathcal{C}_{nice}$ の対象であることに注意せよ。実際これは、
無限次元アフィン空間と $Y_i/X_i$ との積の開部分スキームを取ることで得られる。
小さな細部は省略する。特に、$c^p_{Y_i/X_i}$ を $(b_i/a_i)$ で引き戻したものは、
$c^p_{Y'_i/X'_i}$ を $Y/X \to Y'_i/X'_i$ で引き戻したものと等しい。
$Y'/X'$ が $\mathcal{C}_{nice}$ の対象ならば証明は終わるが、ほとんどの場合そうではない。
実際、その場合には $c^p_{Y'/X'}$ を正方形の二辺に沿って引き戻すことで
望む結論を得る。$Y'/X'$ が $\mathcal{C}_{nice}$ に属さないという問題を避けるため、
上の議論から、スキーム $X, Y, X'_1, Y'_1, X'_2, Y'_2, X', Y'$ をすべて必要に応じて
縮小した後、ある $n, d \geq 1$ を取り、図式を次のように拡張できることに注意する。
$$
\xymatrix{
& Y/X \ar[ld] \ar[rd] \\
Y'_1/X'_1 \ar[rd] & & Y'_2/X'_2 \ar[ld] \\
& Y'/X' \ar[d] \\
& Y_{n, d}/X_{n, d}
}
$$
そこで $c^p_{Y_{n, d}/X_{n, d}}$ から引き戻すことにより、
すでに与えた議論を適用できる。これで証明は完了する。
\end{proof}








\section{ド・ラーム複体上のトレース写像}
\label{derham-section-trace}

\noindent
この節の内容の一部については \cite{Garel} を参照せよ。
$S$ をスキームとする。$f : Y \to X$ を $S$ 上のスキームの有限局所自由射とする。
このときトレース写像
$\text{Trace}_f : f_*\mathcal{O}_Y \to \mathcal{O}_X$ が存在する。
Discriminants, Section \ref{discriminant-section-discriminant} を参照せよ。
この状況で、ド・ラーム複体上のトレース写像とは複体の写像
$$
\Theta_{Y/X} : f_*\Omega^\bullet_{Y/S} \longrightarrow \Omega^\bullet_{X/S}
$$
であって、$\Theta_{Y/X}$ は $\text{Trace}_f$ に等しく（次数 $0$）、かつ
$$
\Theta_{Y/X}(\omega \wedge \eta) = \omega \wedge \Theta_{Y/X}(\eta)
$$
を局所切断 $\omega$ が $\Omega^\bullet_{X/S}$ に属し、局所切断
$\eta$ が $f_*\Omega^\bullet_{Y/S}$ に属するときに満たすものをいう。
このようなトレース写像 $\Theta_{Y/X}$ がすべての有限局所自由射
$Y \to X$ に対して存在するかどうかは明らかでない。反例または証明がある場合は
\href{mailto:stacks.project@gmail.com}{stacks.project@gmail.com}
まで電子メールで知らせてほしい。

\begin{example}
\label{derham-example-no-trace}
ド・ラーム複体上のトレース写像が存在しない例を挙げる。
$\mathbf{C}$-代数 $B = \mathbf{C}[x, y]$ に $G = \{\pm 1\}$ の作用を入れ、
$x \mapsto -x$ および $y \mapsto -y$ で定める。
不変式 $A = B^G$ は $\mathbf{C}$ 上有限型の正規整域であり、
$x^2, xy, y^2$ で生成される。包含 $A \subset B$ に対して、合理的なトレース写像
$\Omega_{B/\mathbf{C}} \to \Omega_{A/\mathbf{C}}$
は $1$-形式上に存在しないと主張する。実際、元
$\omega = x \text{d} y \in \Omega_{B/\mathbf{C}}$
を考える。$\omega$ は $G$ の作用で不変だから、「合理的な」トレース写像が存在すれば、
$2\omega$ は $\Omega_{A/\mathbf{C}} \to \Omega_{B/\mathbf{C}}$ の像に属するはずである。
しかしそうではない。$2\omega$ を
$\text{d}(x^2)$, $\text{d}(xy)$, $\text{d}(y^2)$ の線形結合として、
係数を $B$ に取っても書く方法はない。
この例は \cite{Zannier} の主定理と矛盾する。
\end{example}

\begin{lemma}
\label{derham-lemma-Garel}
すべての有限シントミックなスキーム射 $f : Y \to X$ に
$\mathcal{O}_X$-加群写像
$$
\Theta^p_{Y/X} :
f_*\Omega^p_{Y/\mathbf{Z}}
\longrightarrow
\Omega^p_{X/\mathbf{Z}}
$$
を対応させ、次の性質を満たす規則が一意に存在する。
\begin{enumerate}
\item
$\Omega^p_{X/\mathbf{Z}} \otimes_{\mathcal{O}_X} f_*\mathcal{O}_Y
\to f_*\Omega^p_{Y/\mathbf{Z}}$ との合成は
$\text{id} \otimes \text{Trace}_f$ に等しい。ここで
$\text{Trace}_f : f_*\mathcal{O}_Y \to \mathcal{O}_X$ は
Discriminants, Section \ref{discriminant-section-discriminant} の写像である。
\item この規則は基底変換と両立する。
\end{enumerate}
\end{lemma}

\begin{proof}
まず $X$ が局所ネーター的であると仮定する。
Lemma \ref{derham-lemma-Garel-upstairs} により標準写像
$$
c^p_{Y/X} : \Omega_{Y/\mathbf{Z}}^p
\longrightarrow
f^*\Omega_{X/\mathbf{Z}}^p \otimes_{\mathcal{O}_Y} \det(\NL_{Y/X})
$$
がある。Discriminants, Proposition \ref{discriminant-proposition-tate-map}
により標準同型
$$
c_{Y/X} : \det(\NL_{Y/X}) \to \omega_{Y/X}
$$
があり、$\delta(\NL_{Y/X})$ を $\tau_{Y/X}$ に写す。これらを合わせると
$$
c^p_{Y/X} \otimes c_{Y/X} :
\Omega_{Y/\mathbf{Z}}^p
\longrightarrow
f^*\Omega_{X/\mathbf{Z}}^p \otimes_{\mathcal{O}_Y} \omega_{Y/X}
$$
を得る。Discriminants, Section \ref{discriminant-section-finite-morphisms} により、
これは写像
$$
\Theta_{Y/X}^p :
f_*\Omega_{Y/\mathbf{Z}}^p
\longrightarrow
\Omega_{X/\mathbf{Z}}^p
$$
を与えることと同じである。$c^p_{Y/X} \otimes c_{Y/X}$ と
$\Theta_{Y/X}^p$ の関係には評価写像
$f_*\omega_{Y/X} \to \mathcal{O}_X$
を用い、これは $\tau_{Y/X}$ を $\text{Trace}_f(1)$ に写すことを思い出そう。
Discriminants, Section \ref{discriminant-section-finite-morphisms} を参照せよ。
したがって性質 (1) が成り立つ。$X' \to X$ による基底変換で
$X'$ が局所ネーター的なものについても性質 (2) が成り立つ。実際、
$c^p_{Y/X}$ と $c_{Y/X}$ はいずれもそのような基底変換と両立する。
$f : Y \to X$ が有限シントミックで $X$ が局所ネーター的である場合、
得られた解を引き続き $\Theta^p_{Y/X}$ と書く。

\medskip\noindent
一意性。有限シントミック射 $f: Y \to X$ があり、$X$ は
$\Spec(\mathbf{Z})$ 上滑らかで、$f$ は $X$ の稠密開集合上エタールであるとする。
この場合、$\Theta^p_{Y/X}$ は性質 (1) によって一意に定まると主張する。
実際、写像
$$
\Omega^p_{X/\mathbf{Z}} \otimes_{\mathcal{O}_X} f_*\mathcal{O}_Y \to
f_*\Omega^p_{Y/\mathbf{Z}} \to
\Omega^p_{X/\mathbf{Z}}
$$
を考える。層 $\Omega^p_{X/\mathbf{Z}}$ は、仮定した滑らかさにより捩れなしである。
したがって、$\Theta^p_{Y/X}$ の制限が、$f$ がエタールとなる稠密開集合上で
一意に定まることを確認すれば十分であり、すなわち $f$ がエタールであると仮定してよい。
ところが $f$ がエタールならば
$f^*\Omega_{X/\mathbf{Z}} = \Omega_{Y/\mathbf{Z}}$
なので、表示した式の最初の矢印は同型である。合成はすでに定められているから、
これで一意性が保証される。

\medskip\noindent
$f : Y \to X$ を局所ネーター的スキームの有限シントミック射とする。
$x \in X$ とする。Discriminants, Lemma
\ref{discriminant-lemma-locally-comes-from-universal-finite} により、
$d \geq 1$ と可換図式
$$
\xymatrix{
Y \ar[d] &
V \ar[d] \ar[l] \ar[r] &
V_d \ar[d] \\
X &
U \ar[l] \ar[r] &
U_d
}
$$
を取れ、$x \in U \subset X$ は開集合であり、$V = f^{-1}(U)$ かつ
$V = U \times_{U_d} V_d$ である。したがって $\Theta^p_{Y/X}|_V$ は
写像 $\Theta^p_{V_d/U_d}$ の引き戻しである。
しかし、上の一意性の議論および Discriminants, Lemmas
\ref{discriminant-lemma-universal-finite-syntomic-smooth} および
\ref{discriminant-lemma-universal-finite-syntomic-etale}
により、写像 $\Theta^p_{V_d/U_d}$ は要件 (1) によって一意に定まる。
ゆえに一意性が成り立つ。

\medskip\noindent
ここまでで、すべての有限シントミック射 $Y \to X$ で
$X$ が局所ネーター的なものについて存在と一意性が分かった。ここで Lemma
\ref{derham-lemma-Garel-upstairs} の証明と同様の議論をして一般の $X$ へ拡張してもよい。
しかし、その代わりに絶対ネーター近似を直接用いて証明を完了できる。
すなわち、$\Theta^p_{Y/X}$ を構成するには $X$ 上ザリスキー局所的に構成すれば十分である
（一意性も示すものとする）。したがって $X$ はアフィンであると仮定してよい
（小さな細部は省略する）。このとき $X = \lim_{i \in I} X_i$ と書ける。
これはネーター的アフィンスキームの有向集合 $I$ 上の極限である。
Algebra, Lemma \ref{algebra-lemma-colimit-category-fp-algebras} により、
$0 \in I$ とアフィンの有限表示射 $f_0 : Y_0 \to X_0$ を取れ、その
$X$ への基底変換は $Y \to X$ である。$0$ を大きくすれば
$Y_0 \to X_0$ は有限かつシントミックであると仮定してよい。
Algebra, Lemma \ref{algebra-lemma-colimit-lci} および
\ref{algebra-lemma-colimit-finite} を参照せよ。$i \geq 0$ に対して基底変換
$f_i : Y_i = Y_0 \times_{X_0} X_i \to X_i$ も有限シントミックである。このとき
$$
\Gamma(X, f_*\Omega^p_{Y/\mathbf{Z}}) =
\Gamma(Y, \Omega^p_{Y/\mathbf{Z}}) =
\colim_{i \geq 0} \Gamma(Y_i, \Omega^p_{Y_i/\mathbf{Z}}) =
\colim_{i \geq 0} \Gamma(X_i, f_{i, *}\Omega^p_{Y_i/\mathbf{Z}})
$$
したがって $\Theta^p_{Y/X}$ を写像 $\Theta^p_{Y_i/X_i}$ の余極限として定義できる
（またそう定義せざるを得ない）。この写像は、任意のカルテシアン図式
$$
\xymatrix{
Y' \ar[r] \ar[d] & Y \ar[d] \\
X' \ar[r] & X
}
$$
で $X'$ がアフィンであるものと両立する。実際、上の議論によりネーター的アフィンスキームの
場合にはこれが分かっている（小さな細部は省略する。ヒント：さらに
$X' = \lim_{j \in J} X'_j$ と書けば、すべての $i \in I$ に対して
$j \in J$ と射 $X'_j \to X_i$ が存在し、射 $X' \to X$ と両立する）。
これで証明は完了する。
\end{proof}

\begin{proposition}
\label{derham-proposition-Garel}
\begin{reference}
\cite{Garel}
\end{reference}
$f : Y \to X$ をスキームの有限シントミック射とする。
Lemma \ref{derham-lemma-Garel} の写像 $\Theta^p_{Y/X}$ は複体の写像
$$
\Theta_{Y/X} :
f_*\Omega^\bullet_{Y/\mathbf{Z}}
\longrightarrow
\Omega^\bullet_{X/\mathbf{Z}}
$$
を定め、次の性質をもつ。
\begin{enumerate}
\item 次数 $0$ では
$\text{Trace}_f : f_*\mathcal{O}_Y \to \mathcal{O}_X$
を得る。Discriminants, Section \ref{discriminant-section-discriminant} を参照せよ。
\item
$\Theta_{Y/X}(\omega \wedge \eta) = \omega \wedge \Theta_{Y/X}(\eta)$
が、$\omega$ が $\Omega^\bullet_{X/\mathbf{Z}}$ に属し、$\eta$ が
$f_*\Omega^\bullet_{Y/\mathbf{Z}}$ に属するときに成り立つ。
\item $f$ が基底スキーム $S$ 上の射ならば、$\Theta_{Y/X}$ は複体の写像
$f_*\Omega^\bullet_{Y/S} \to \Omega^\bullet_{X/S}$
を誘導する。
\end{enumerate}
\end{proposition}

\begin{proof}
Discriminants, Lemma
\ref{discriminant-lemma-locally-comes-from-universal-finite} により、
すべての $x \in X$ に対して $d \geq 1$ と可換図式
$$
\xymatrix{
Y \ar[d] &
V \ar[d] \ar[l] \ar[r] &
V_d \ar[d] \ar[r] &
Y_d = \Spec(B_d) \ar[d] \\
X &
U \ar[l] \ar[r] &
U_d \ar[r] &
X_d = \Spec(A_d)
}
$$
を取れ、$x \in U \subset X$ はアフィン開集合であり、$V = f^{-1}(U)$ かつ
$V = U \times_{U_d} V_d$ である。$U = \Spec(A)$ および $V = \Spec(B)$ と書く。
$B = A \otimes_{A_d} B_d$ であることに注意し、
$B_d = A_d e_1 \oplus \ldots \oplus A_d e_d$ を思い出そう。
$a_1, \ldots, a_r \in A$ および $b_1, \ldots, b_s \in B$ があるとする。
$b_j = \sum a_{j, l} e_l$ と書け、ここで $a_{j, l} \in A$ である。
$N = r + sd$ とおき、次の分解を考える。
$$
\xymatrix{
V \ar[r] \ar[d] &
V' = \mathbf{A}^N \times V_d \ar[r] \ar[d] &
V_d \ar[d] \\
U \ar[r]&
U' = \mathbf{A}^N \times U_d \ar[r] &
U_d
}
$$
ここで右下の水平矢印は、前の図式の射 $U \to U_d$ と、射
$U \to \mathbf{A}^N$ から定まり、後者は
$a_1, \ldots, a_r, a_{1, 1}, \ldots, a_{s, d}$ によって与えられる。
このとき関数 $a_1, \ldots, a_r$ は
$\Gamma(U', \mathcal{O}_{U'}) \to \Gamma(U, \mathcal{O}_U)$ の像に属し、
関数 $b_1, \ldots, b_s$ は
$\Gamma(V', \mathcal{O}_{V'}) \to \Gamma(V, \mathcal{O}_V)$ の像に属する。
このようにして、次の群の元の任意の有限集合について次が分かる。\footnote{結局、
これらの元はすべて、次の形の元の有限和
$a_0 \text{d}a_1 \wedge \ldots \wedge \text{d}a_i$ で
$a_0, \ldots, a_i \in A$、または次の形の元の有限和
$b_0 \text{d}b_1 \wedge \ldots \wedge \text{d}b_j$ で
$b_0, \ldots, b_j \in B$ である。}
$$
\Gamma(V, \Omega^i_{Y/\mathbf{Z}}),\quad i = 0, 1, 2, \ldots
\quad\text{および}\quad
\Gamma(U, \Omega^j_{X/\mathbf{Z}}),\quad j = 0, 1, 2, \ldots
$$
上と同様に分解 $V \to V' \to V_d$ および $U \to U' \to U_d$ を取れ、
$V' = \mathbf{A}^N \times V_d$ かつ $U' = \mathbf{A}^N \times U_d$ であり、
これらの切断は次の切断の引き戻しとなる。
$$
\Gamma(V', \Omega^i_{V'/\mathbf{Z}}),\quad i = 0, 1, 2, \ldots
\quad\text{および}\quad
\Gamma(U', \Omega^j_{U'/\mathbf{Z}}),\quad j = 0, 1, 2, \ldots
$$
この帰結として、
$\text{d} \circ \Theta_{Y/X} = \Theta_{Y/X} \circ \text{d}$
を確認するには、これが $\Theta_{V'/U'}$ に対して成り立つことを確認すれば十分である。
補題の性質 (2) についても同様である。

\medskip\noindent
Discriminants, Lemmas
\ref{discriminant-lemma-universal-finite-syntomic-smooth} および
\ref{discriminant-lemma-universal-finite-syntomic-etale}
により、スキーム $U_d$ は滑らかであり、射 $V_d \to U_d$ は
$U_d$ の稠密開集合上エタールである。したがって射 $V' \to U'$ についても同じことが成り立つ。
$\Omega_{U'/\mathbf{Z}}$ は局所自由なので $\Omega^p_{U'/\mathbf{Z}}$ は捩れなしであり、
$V'$ が有限エタールとなる開集合に制限して望む関係を確認すれば十分である。
さらに、全射エタール基底変換の後で関係を確認してよい。
したがって有限エタール被覆を分裂させ、次の形の射を考えていると仮定してよい。
$$
\coprod\nolimits_{i = 1, \ldots, d} W \longrightarrow W
$$
ここで $W$ は $\mathbf{Z}$ 上滑らかである。この場合、構成の任意の局所的性質は
（真である限り）自明に確認できる。これで (1) と (2) の証明が完了する。

\medskip\noindent
最後に、(3) は (2) から従うことに注意する。実際、$\Omega_{Y/S}$ は
$\Omega_{Y/\mathbf{Z}}$ を $\Omega_{S/\mathbf{Z}}$ の局所切断の引き戻しで
生成される部分加群で割った商である。
\end{proof}

\begin{example}
\label{derham-example-Garel}
$A$ を環とする。$f = x^d + \sum_{0 \leq i < d} a_{d - i} x^i \in A[x]$ とおく。
$B = A[x]/(f)$ とおく。Proposition \ref{derham-proposition-Garel} により複体の射
$$
\Theta_{B/A} : \Omega^\bullet_B \longrightarrow \Omega^\bullet_A
$$
がある。特に、$t \in B$ が $x \in A[x]$ の像を表すならば、元
$$
\Theta_{B/A}(t^i\text{d}t) \in \Omega^1_A,\quad i = 0, \ldots, d - 1
$$
を考えられる。これらの元は何か。Proposition \ref{derham-proposition-Garel} の証明で用いたのと
同じ原理により、普遍的な場合、すなわち $A = \mathbf{Z}[a_1, \ldots, a_d]$ の場合に
計算すれば十分である。さらに $A$ を分数体
$\mathbf{Q}(a_1, \ldots, a_d)$ で置き換えてもよい。記号的に
$$
f = \prod\nolimits_{i = 1, \ldots, d} (x - \alpha_i)
$$
と書けば、$\mathbf{Q}(\alpha_1, \ldots, \alpha_d)$ 上で代数 $B$ は分裂する。
$$
\mathbf{Q}(a_0, \ldots, a_{d - 1})[x]/(f) 
\longrightarrow
\prod\nolimits_{i = 1, \ldots, d} \mathbf{Q}(\alpha_1, \ldots, \alpha_d),
\quad
t \longmapsto (\alpha_1, \ldots, \alpha_d)
$$
したがって、例えば
$$
\Theta(\text{d}t) = \sum \text{d} \alpha_i = - \text{d}a_1
$$
次に
$$
\Theta(t\text{d}t) = \sum \alpha_i \text{d}\alpha_i =
a_1 \text{d} a_1 - \text{d}a_2
$$
さらに
$$
\Theta(t^2\text{d}t) = \sum \alpha_i^2 \text{d}\alpha_i =
- a_1^2 \text{d} a_1 + a_1 \text{d}a_2 + a_2 \text{d}a_1 - \text{d}a_3
$$
となる（計算誤差を除く）。以下も同様である。これは、$f(x) = x^d - a$ ならば
$$
\Theta_{B/A}(t^i\text{d}t) =
\left\{
\begin{matrix}
0 & i = 0, \ldots, d - 2 & \text{のとき} \\
\text{d}a & i = d - 1 & \text{のとき}
\end{matrix}
\right.
$$
が $\Omega_A$ において成り立つことを示唆する。これは実際に正しい。
この特別な場合には、拡大 $\mathbf{Q}(a)[x]/(x^d - a)$ について計算すれば
直接確認できる。
\end{example}

\begin{lemma}
\label{derham-lemma-Garel-map-frobenius-smooth-char-p}
$p$ を素数とする。$X \to S$ を相対次元 $d$ の滑らかな射で、標数 $p$ のスキームの射とする。
相対 Frobenius $F_{X/S} : X \to X^{(p)}$ は $X/S$ に対するものであり
（Varieties, Definition \ref{varieties-definition-relative-frobenius}）、有限シントミックである。また、対応する写像
$$
\Theta_{X/X^{(p)}} :
F_{X/S, *}\Omega^\bullet_{X/S} \to \Omega^\bullet_{X^{(p)}/S}
$$
は次数 $d$ を除くすべての次数で零であり、その次数では全射を定める。
\end{lemma}

\begin{proof}
Varieties, Lemma \ref{varieties-lemma-relative-frobenius-finite} により、
$F_{X/S}$ は有限射である。$F_{X/S}$ が平坦であることを示すには、
ファイバー間の射 $F_{X/S, s} : X_s \to X^{(p)}_s$ が
すべての $s \in S$ に対して平坦であることを示せば十分である。
More on Morphisms, Theorem \ref{more-morphisms-theorem-criterion-flatness-fibre} を参照せよ。
$X_s \to X^{(p)}_s$ の平坦性は Algebra, Lemma
\ref{algebra-lemma-CM-over-regular-flat} と、すでに示した有限性から従う。
More on Morphisms, Lemma \ref{more-morphisms-lemma-lci-permanence} により、
射 $F_{X/S}$ は局所完全交叉射である。したがって $F_{X/S}$ は有限シントミックである
（More on Morphisms, Lemma \ref{more-morphisms-lemma-flat-lci} を参照）。

\medskip\noindent
任意の点 $x \in X$ に対し、次の可換図式を選ぶことができる。
$$
\xymatrix{
X \ar[d] & U \ar[l] \ar[d]_\pi \\
S & \mathbf{A}^d_S \ar[l]
}
$$
ここで $\pi$ はエタールであり、$x \in U$ で、この開部分スキームは $X$ に含まれる。
Morphisms, Lemma \ref{morphisms-lemma-smooth-etale-over-affine-space} を参照せよ。
さらに、
$\mathbf{A}^d_S \to \mathbf{A}^d_S$, $(x_1, \ldots, x_d) \mapsto
(x_1^p, \ldots, x_d^p)$ は $\mathcal{A}^d_S$ の $S$ 上の相対 Frobenius である。
次の可換図式を考える。
$$
\xymatrix{
U \ar[d]_\pi \ar[r]_{F_{X/S}} & U^{(p)} \ar[d]^{\pi^{(p)}} \\
\mathbf{A}^d_S \ar[r]^{x_i \mapsto x_i^p} & \mathbf{A}^d_S
}
$$
これは Varieties, Lemma \ref{varieties-lemma-relative-frobenius-endomorphism-identity}
における $\pi : U \to \mathbf{A}^d_S$ に対する図式であり、\'Etale Morphisms, Lemma
\ref{etale-lemma-relative-frobenius-etale} によりカルテシアンである。
$\Theta$ の構成は基底変換と両立し、さらに
$\Omega_{U/S} = \pi^*\Omega_{\mathbf{A}^d_S/S}$
（Lemma \ref{derham-lemma-etale}）であるから、この補題を
$\mathbf{A}^d_S$ に対して示せば十分である。

\medskip\noindent
$A$ を標数 $p$ の環とする。$A$-代数準同型
$A[y_1, \ldots, y_d] \to A[x_1, \ldots, x_d]$
のうち、$y_i$ を $x_i^p$ に送る一意なものを考える。以上の議論により、次の写像を計算すればよい。
$$
\Theta^i : \Omega^i_{A[x_1, \ldots, x_d]/A} \to
\Omega^i_{A[y_1, \ldots, y_d]/A}
$$
この場合の計算を読者自身で行うことを勧める。
Example \ref{derham-example-Garel} と同様に、$\Theta^i$ の公式を普遍的な場合
$$
\mathbf{Z}[y_1, \ldots, y_d] \to \mathbf{Z}[x_1, \ldots, x_d],\quad
y_i \mapsto x_i^p
$$
において計算することに帰着できる。さらに $\Theta^i$ の公式は、複素数の場合、すなわち
$\mathbf{C}$-代数準同型
$$
\mathbf{C}[y_1, \ldots, y_d] \to \mathbf{C}[x_1, \ldots, x_d],\quad
y_i \mapsto x_i^p
$$
に対して計算すれば得られる。さらに $x_1, \ldots, x_d$ と $y_1, \ldots, y_d$ を
可逆にしてもよい。この場合、$\text{d}x_i = p^{-1} x_i^{- p + 1}\text{d}y_i$ である。
したがって、
\begin{align*}
\Theta^i(
x_1^{e_1} \ldots x_d^{e_d} \text{d}x_1 \wedge \ldots \wedge \text{d}x_i)
& =
p^{-i} \Theta^i(
x_1^{e_1 - p + 1} \ldots x_i^{e_i - p + 1} x_{i + 1}^{e_{i + 1}} \ldots
x_d^{e_d} \text{d}y_1 \wedge \ldots \wedge \text{d}y_i ) \\
& =
p^{-i} \text{Trace}(x_1^{e_1 - p + 1} \ldots x_i^{e_i - p + 1}
x_{i + 1}^{e_{i + 1}} \ldots x_d^{e_d})
\text{d}y_1 \wedge \ldots \wedge \text{d}y_i
\end{align*}
が $\Theta^i$ の性質から従う。初等的な計算により、上式のトレースは、
$e_1, \ldots, e_i$ が $-1$ と法 $p$ で合同であり、かつ
$e_{i + 1}, \ldots, e_d$ が $p$ で割り切れる場合を除いて零である。
さらに、この場合には
$$
p^{d - i} y_1^{(e_1 - p + 1)/p} \ldots y_i^{(e_i - p + 1)/p}
y_{i + 1}^{e_{i + 1}/p} \ldots y_d^{e_d/p}
\text{d}y_1 \wedge \ldots \wedge \text{d}y_i
$$
を得る。したがって標数 $p$ では $d = i$ でない限り零となり、
この場合には任意の $d$-形式が像に現れる。
\end{proof}








\section{ポアンカレ双対性}
\label{derham-section-poincare-duality}

\noindent
この節では、体上の固有滑らかなスキームのド・ラームコホモロジーに対する
ポアンカレ双対性を証明する。まず、これが Hodge コホモロジーに対して
どのように働くかを説明する。

\begin{lemma}
\label{derham-lemma-duality-hodge}
$k$ を体とする。$X$ を $k$ 上の空でない滑らかな固有スキームで、
次元 $d$ の等次元なものとする。$k$-線形写像
$$
t : H^d(X, \Omega^d_{X/k}) \longrightarrow k
$$
が存在し、$H^0(X, \mathcal{O}_X)$ の単元による乗法を前合成する違いを除いて一意であり、
次の性質をもつ。すべての $p, q$ に対し、ペアリング
$$
H^q(X, \Omega^p_{X/k}) \times H^{d - q}(X, \Omega^{d - p}_{X/k})
\longrightarrow
k, \quad
(\xi, \xi') \longmapsto t(\xi \cup \xi')
$$
は完全である。
\end{lemma}

\begin{proof}
Duality for Schemes, Lemma \ref{duality-lemma-duality-proper-over-field}
により、$\omega_X^\bullet = \Omega^d_{X/k}[d]$ である。
$\Omega_{X/k}$ は階数 $d$ の局所自由加群であるから
（Morphisms, Lemma \ref{morphisms-lemma-smooth-omega-finite-locally-free}）、
$$
\Omega^d_{X/k} \otimes_{\mathcal{O}_X} (\Omega^p_{X/k})^\vee
\cong
\Omega^{d - p}_{X/k}
$$
を得る。したがって Duality for Schemes, Lemma
\ref{duality-lemma-duality-proper-over-field-perfect} により、主張を満たす
$k$-線形写像 $t : H^d(X, \Omega^d_{X/k}) \to k$ を得る。特にペアリング
$H^0(X, \mathcal{O}_X) \times H^d(X, \Omega^d_{X/k}) \to k$
は完全である。このことから、任意の $k$-線形写像
$t' : H^d(X, \Omega^d_{X/k}) \to k$ は
$\xi \mapsto t(g\xi)$ の形に書ける。ここで $g \in H^0(X, \mathcal{O}_X)$ である。もちろん、$t'$ が引き続き
$H^0(X, \mathcal{O}_X)$ と $H^d(X, \Omega^d_{X/k})$ の間の双対性を与えるためには、
$g$ は単元でなければならない。$\langle -, - \rangle_{p, q}$ を
$t$ を用いて構成したペアリングと書き、$\langle -, - \rangle'_{p, q}$ を
$t'$ を用いて構成したペアリングと書く。明らかに
$$
\langle \xi, \xi' \rangle'_{p, q} =
\langle g\xi, \xi' \rangle_{p, q}
$$
が $\xi \in H^q(X, \Omega^p_{X/k})$ および
$\xi' \in H^{d - q}(X, \Omega^{d - p}_{X/k})$ に対して成り立つ。$g$ は単元、すなわち
可逆であるから、$t'$ を用いても $t$ による場合と同様に、すべての $p, q$ に対して
完全ペアリングが得られる。
\end{proof}

\begin{lemma}
\label{derham-lemma-bottom-part-degenerates}
$k$ を体とする。$X$ を $k$ 上の滑らかな固有スキームとする。写像
$$
\text{d} : H^0(X, \mathcal{O}_X) \to H^0(X, \Omega^1_{X/k})
$$
は零である。
\end{lemma}

\begin{proof}
$X$ は $k$ 上滑らかなので、$k$ 上幾何学的被約である。
Varieties, Lemma \ref{varieties-lemma-smooth-geometrically-normal} を参照せよ。
したがって $H^0(X, \mathcal{O}_X) = \prod k_i$ は有限次分離体拡大
$k_i/k$ の有限直積である。Varieties, Lemma
\ref{varieties-lemma-proper-geometrically-reduced-global-sections} を参照せよ。
ゆえに $\Omega_{H^0(X, \mathcal{O}_X)/k} = \prod \Omega_{k_i/k} = 0$ である
（例えば Algebra, Lemma
\ref{algebra-lemma-characterize-separable-algebraic-field-extensions} を参照）。
補題の写像は、ド・ラーム複体の関手性により
$$
H^0(X, \mathcal{O}_X) \to
\Omega_{H^0(X, \mathcal{O}_X)/k} \to
H^0(X, \Omega_{X/k})
$$
と分解するので（Section \ref{derham-section-de-rham-complex} を参照）、結論を得る。
\end{proof}

\begin{lemma}
\label{derham-lemma-top-part-degenerates}
$k$ を体とする。$X$ を $k$ 上の滑らかな固有スキームで、
次元 $d$ の等次元なものとする。写像
$$
\text{d} : H^d(X, \Omega^{d - 1}_{X/k}) \to H^d(X, \Omega^d_{X/k})
$$
は零である。
\end{lemma}

\begin{proof}
これは Lemmas \ref{derham-lemma-bottom-part-degenerates} と \ref{derham-lemma-duality-hodge}
を組み合わせれば従うと思いたくなる。しかし、この議論は成立しない。なぜなら写像
$\mathcal{O}_X \to \Omega^1_{X/k}$ と
$\Omega^{d - 1}_{X/k} \to \Omega^d_{X/k}$ は $\mathcal{O}_X$-線形ではないからである。
したがって Duality for Schemes, Remark
\ref{duality-remark-coherent-duality-proper-over-field} で論じた関手性を用いて、
Lemma \ref{derham-lemma-bottom-part-degenerates} の写像がこの補題の写像と双対であると
結論することはできない。

\medskip\noindent
$X$ を $X$ の連結成分で置き換えてよい。したがって $X$ は既約であると仮定してよい。
Varieties, Lemmas \ref{varieties-lemma-smooth-geometrically-normal} および
\ref{varieties-lemma-proper-geometrically-reduced-global-sections} により、
$k' = H^0(X, \mathcal{O}_X)$ は有限次分離拡大 $k'/k$ である。
$\Omega_{k'/k} = 0$ なので（例えば Algebra, Lemma
\ref{algebra-lemma-characterize-separable-algebraic-field-extensions} を参照）、
$\Omega_{X/k} = \Omega_{X/k'}$ である
（Morphisms, Lemma \ref{morphisms-lemma-triangle-differentials} を参照）。
よって $k$ を $k'$ で置き換え、$H^0(X, \mathcal{O}_X) = k$ と仮定してよい。

\medskip\noindent
$H^0(X, \mathcal{O}_X) = k$ と仮定する。Lemma \ref{derham-lemma-duality-hodge} により
$\dim H^d(X, \Omega^d_{X/k}) = 1$ を得る。まず $k$ の標数が素数 $p$ であると仮定する。
$F_{X/k} : X \to X^{(p)}$ を $X$ の $k$ 上の相対 Frobenius と書く。
この射について Lemma \ref{derham-lemma-Garel-map-frobenius-smooth-char-p} で証明した事実を
念頭に置く。次の可換図式を考える。
$$
\xymatrix{
H^d(X, \Omega^{d - 1}_{X/k}) \ar[d] \ar[r] &
H^d(X^{(p)}, F_{X/k, *}\Omega^{d - 1}_{X/k}) \ar[d] \ar[r]_{\Theta^{d - 1}} &
H^d(X^{(p)}, \Omega^{d - 1}_{X^{(p)}/k}) \ar[d] \\
H^d(X, \Omega^d_{X/k}) \ar[r] &
H^d(X^{(p)}, F_{X/k, *}\Omega^d_{X/k}) \ar[r]^{\Theta^d} &
H^d(X^{(p)}, \Omega^d_{X^{(p)}/k})
}
$$
左側の二つの水平矢印は、$F_{X/k}$ が有限であるため同型である。
Cohomology of Schemes, Lemma \ref{coherent-lemma-relative-affine-cohomology} を参照せよ。
右側の正方形は、$\Theta_{X^{(p)}/X}$ が複体の射であり
$\Theta^{d - 1}$ が零であるため可換である。したがって
$\Theta^d$ が零でないことを示せば十分である（上の議論により写像
$\Theta^d$ の始域の次元は $1$ である）。一方、
$$
\Theta^d : F_{X/k, *}\Omega^d_{X/k} \to \Omega^d_{X^{(p)}/k}
$$
は全射である。したがって右完全関手 $H^d(X^{(p)}, -)$ を適用した後も全射である
（Cohomology, Proposition \ref{cohomology-proposition-vanishing-Noetherian}
から従う、次数 $d$ を超えるコホモロジーの消滅により右完全である）。
最後に、$H^d(X^{(d)}, \Omega^d_{X^{(d)}/k})$ は零でない。実際、例えば
$H^0(X^{(d)}, \mathcal{O}_{X^{(p)}})$ と双対である。これは
$X^{(p)}$ に $k$ 上の Lemma \ref{derham-lemma-duality-hodge} を適用すればよい。
これでこの場合の証明は完了する。

\medskip\noindent
最後に、$k$ の標数が $0$ であると仮定する。
$k$ を有限型 $\mathbf{Z}$-部分代数 $R$ のフィルター付き余極限として書ける。
そのうちの一つに対して、スキームのカルテシアン図式
$$
\xymatrix{
X \ar[d] \ar[r] & Y \ar[d] \\
\Spec(k) \ar[r] & \Spec(R)
}
$$
を取れ、$Y \to \Spec(R)$ は相対次元 $d$ の滑らかな固有射となる。
Limits, Lemmas \ref{limits-lemma-descend-finite-presentation},
\ref{limits-lemma-descend-smooth}, \ref{limits-lemma-descend-dimension-d}, および
\ref{limits-lemma-eventually-proper} を参照せよ。
加群 $M^{i, j} = H^j(Y, \Omega^i_{Y/R})$ は有限 $R$-加群である。
Cohomology of Schemes, Lemma
\ref{coherent-lemma-proper-over-affine-cohomology-finite} を参照せよ。
したがって $R$ を局所化で置き換えれば、これらの加群はすべて有限自由であると仮定してよい。
$M^{i, j} \otimes_R k = H^j(X, \Omega^i_{X/k})$
が平坦基底変換により成り立つ
（Cohomology of Schemes, Lemma
\ref{coherent-lemma-flat-base-change-cohomology}）。
したがって $M^{d - 1, d} \to M^{d, d}$ が零であることを示せば十分である。
これは整域上の有限自由加群の間の写像だから、素イデアル
$\mathfrak p \subset R$ の稠密集合であって、$\kappa(\mathfrak p)$ と
テンソルを取った後に零となるものを見つければよい。
$R$ は $\mathbf{Z}$ 上有限型なので、剰余体の標数が正である素イデアル
$\mathfrak p$ 全体を取れる（細部は省略する）。ここで
$$
M^{d - 1, d} \otimes_R \kappa(\mathfrak p) =
H^d(Y_{\kappa(\mathfrak p)},
\Omega^{d - 1}_{Y_{\kappa(\mathfrak p)}/\kappa(\mathfrak p)})
$$
が、例えば Limits, Lemma
\ref{limits-lemma-higher-direct-images-zero-above-dimension-fibre} により成り立つ。
$M^{d, d}$ についても同様である。したがって
$M^{d - 1, d} \otimes_R \kappa(\mathfrak p) \to
M^{d, d} \otimes_R \kappa(\mathfrak p)$
は、上で扱った正標数の場合により零である。
\end{proof}

\begin{proposition}
\label{derham-proposition-poincare-duality}
$k$ を体とする。$X$ を $k$ 上の空でない滑らかな固有スキームで、
次元 $d$ の等次元なものとする。$k$-線形写像
$$
t : H^{2d}_{dR}(X/k) \longrightarrow k
$$
が存在し、$H^0(X, \mathcal{O}_X)$ の単元による乗法を前合成する違いを除いて一意であり、
次の性質をもつ。すべての $i$ に対し、ペアリング
$$
H^i_{dR}(X/k) \times H_{dR}^{2d - i}(X/k)
\longrightarrow
k, \quad
(\xi, \xi') \longmapsto t(\xi \cup \xi')
$$
は完全である。
\end{proposition}

\begin{proof}
Hodge-to-de Rham スペクトル系列
（Section \ref{derham-section-hodge-to-de-rham}）、$\Omega^i_{X/k}$ の
$i > d$ に対する消滅、Cohomology, Proposition
\ref{cohomology-proposition-vanishing-Noetherian} の消滅定理、および
Lemmas \ref{derham-lemma-bottom-part-degenerates} と
\ref{derham-lemma-top-part-degenerates} の結果により、
$H^0_{dR}(X/k) = H^0(X, \mathcal{O}_X)$ および
$H^d(X, \Omega^d_{X/k}) = H_{dR}^{2d}(X/k)$ を得る。
より正確には、これらの同一視は複体の写像
$$
\Omega^\bullet_{X/k} \to \mathcal{O}_X[0]
\quad\text{および}\quad
\Omega^d_{X/k}[-d] \to \Omega^\bullet_{X/k}
$$
から得られる。$t : H_{dR}^{2d}(X/k) \to k$ を、この同一視を通じて
Lemma \ref{derham-lemma-duality-hodge} の $t$ に対応するように選ぶ。
すると、いずれにせよ補題に表示したペアリングは $i = 0$ に対して完全である。

\medskip\noindent
$\underline{k}$ を値 $k$ の $X$ 上の定数層とする。
$\Omega^\bullet = \Omega^\bullet_{X/k}$ と略記する。
この状況では写像 (\ref{derham-equation-wedge}) は
$$
\wedge :
\text{Tot}(\Omega^\bullet \otimes_{\underline{k}} \Omega^\bullet)
\longrightarrow
\Omega^\bullet
$$
と書ける。任意の整数 $p = 0, 1, \ldots, d$ に対し、この写像は次数の理由から部分複体
$\text{Tot}(\sigma_{> p} \Omega^\bullet \otimes_{\underline{k}}
\sigma_{\geq d - p} \Omega^\bullet)$ を零化する。
したがって $\wedge$ の部分複体
$\text{Tot}(\Omega^\bullet \otimes_{\underline{k}}
\sigma_{\geq d - p}\Omega^\bullet)$ への制限は、複体の写像
$$
\gamma_p :
\text{Tot}(\sigma_{\leq p} \Omega^\bullet \otimes_{\underline{k}}
\sigma_{\geq d - p} \Omega^\bullet)
\longrightarrow
\Omega^\bullet
$$
を経由する。Section \ref{derham-section-cup-product} と同じ手順により、カップ積
$$
H^i(X, \sigma_{\leq p} \Omega^\bullet) \times
H^{2d - i}(X, \sigma_{\geq d - p}\Omega^\bullet)
\longrightarrow
H_{dR}^{2d}(X, \Omega^\bullet)
$$
$p$ に関する帰納法により、これらのカップ積が $t$ を介して
$H^i(X, \sigma_{\leq p} \Omega^\bullet)$ と
$H^{2d - i}(X, \sigma_{\geq d - p}\Omega^\bullet)$ の間の完全ペアリングを
誘導することを証明する。$p = d$ の場合が命題の主張である。

\medskip\noindent
基底の場合は $p = 0$ である。この場合は単に Lemma \ref{derham-lemma-duality-hodge} の
$H^i(X, \mathcal{O}_X)$ と $H^{d - i}(X, \Omega^d)$ の間のペアリングを得るので、
結果は成り立つ。

\medskip\noindent
帰納段階。結果が $p$ に対して成り立つと仮定する。このとき識別三角
$$
\Omega^{p + 1}[-p - 1] \to
\sigma_{\leq p + 1}\Omega^\bullet \to
\sigma_{\leq p}\Omega^\bullet \to
\Omega^{p + 1}[-p]
$$
および識別三角
$$
\sigma_{\geq d - p}\Omega^\bullet \to
\sigma_{\geq d - p - 1}\Omega^\bullet \to
\Omega^{d - p - 1}[-d + p + 1] \to
(\sigma_{\geq d - p}\Omega^\bullet)[1]
$$
を考える。いずれも $\underline{k}$-加群の層の複体のホモトピー圏における識別三角である。
特に、写像 $\sigma_{\leq p}\Omega^\bullet \to \Omega^{p + 1}[-p]$ および
$\Omega^{d - p - 1}[-d + p + 1] \to (\sigma_{\geq d - p}\Omega^\bullet)[1]$
は実際の複体の写像で与えられ、具体的には微分
$\Omega^p \to \Omega^{p + 1}$ および微分
$\Omega^{d - p - 1} \to \Omega^{d - p}$ を用いる。
これらの識別三角に付随するコホモロジーの長完全系列を考える。
$$
\xymatrix{
H^{i - 1}(X, \sigma_{\leq p}\Omega^\bullet) \ar[d]_a \\
H^i(X, \Omega^{p + 1}[-p - 1]) \ar[d]_b \\
H^i(X, \sigma_{\leq p + 1}\Omega^\bullet) \ar[d]_c \\
H^i(X, \sigma_{\leq p}\Omega^\bullet) \ar[d]_d \\
H^{i + 1}(X, \Omega^{p + 1}[-p - 1])
}
\quad\quad
\xymatrix{
H^{2d - i  + 1}(X, \sigma_{\geq d - p}\Omega^\bullet) \\
H^{2d - i}(X, \Omega^{d - p - 1}[-d + p + 1]) \ar[u]_{a'} \\
H^{2d - i}(X, \sigma_{\geq d - p - 1}\Omega^\bullet) \ar[u]_{b'} \\
H^{2d - i}(X, \sigma_{\geq d - p}\Omega^\bullet) \ar[u]_{c'} \\
H^{2d - i - 1}(X, \Omega^{d - p - 1}[-d + p + 1]) \ar[u]_{d'}
}
$$
帰納法と Lemma \ref{derham-lemma-duality-hodge} により、第一、第二、第四、および第五行の
$k$-ベクトル空間の間に上で構成したペアリングは完全である。$5$-補題により、
中央行のコホモロジー群の間のペアリングが完全であることを示すには、組
$(a, a')$, $(b, b')$, $(c, c')$, および $(d, d')$ が
与えられたペアリングと両立することを示せば十分である（以下を参照）。

\medskip\noindent
まず組 $(c, c')$ についてこれを証明する。この場合、単に次の可換図式がある。
$$
\xymatrix{
\text{Tot}(\sigma_{\leq p} \Omega^\bullet \otimes_{\underline{k}}
\sigma_{\geq d - p} \Omega^\bullet) \ar[d]_{\gamma_p} &
\text{Tot}(\sigma_{\leq p + 1} \Omega^\bullet \otimes_{\underline{k}}
\sigma_{\geq d - p} \Omega^\bullet) \ar[l] \ar[d] \\
\Omega^\bullet &
\text{Tot}(\sigma_{\leq p + 1} \Omega^\bullet \otimes_{\underline{k}}
\sigma_{\geq d - p - 1} \Omega^\bullet) \ar[l]_-{\gamma_{p + 1}}
}
$$
したがって $\alpha \in H^i(X, \sigma_{\leq p + 1}\Omega^\bullet)$ および
$\beta \in H^{2d - i}(X, \sigma_{\geq d - p}\Omega^\bullet)$ があれば、
$\gamma_p(\alpha \cup c'(\beta)) = \gamma_{p + 1}(c(\alpha) \cup \beta)$
をカップ積の関手性により得る。

\medskip\noindent
同様に、組 $(b, b')$ については可換図式
$$
\xymatrix{
\text{Tot}(\sigma_{\leq p + 1} \Omega^\bullet \otimes_{\underline{k}}
\sigma_{\geq d - p - 1} \Omega^\bullet) \ar[d]_{\gamma_{p + 1}} &
\text{Tot}(\Omega^{p + 1}[-p - 1] \otimes_{\underline{k}}
\sigma_{\geq d - p - 1} \Omega^\bullet) \ar[l] \ar[d] \\
\Omega^\bullet &
\Omega^{p + 1}[-p - 1]
\otimes_{\underline{k}}
\Omega^{d - p - 1}[-d + p + 1] \ar[l]_-\wedge
}
$$
を用いて同じように論じる。

\medskip\noindent
組 $(d, d')$ については可換図式
$$
\xymatrix{
\Omega^{p + 1}[-p] \otimes_{\underline{k}}
\Omega^{d - p - 1}[-d + p] \ar[d] &
\text{Tot}(\sigma_{\leq p}\Omega^\bullet \otimes_{\underline{k}}
\Omega^{d - p - 1}[-d + p]) \ar[l] \ar[d] \\
\Omega^\bullet &
\text{Tot}(\sigma_{\leq p}\Omega^\bullet \otimes_{\underline{k}}
\sigma_{\geq d - p}\Omega^\bullet) \ar[l]
}
$$
を用い、$H^i(X, \sigma_{\leq p}\Omega^\bullet)$ および
$H^{2d - i}(X, \Omega^{d - p - 1}[-d + p])$ のコホモロジー類を考える。
$i$ を $i - 1$ に変えれば組 $(a, a')$ に対する結果を得る。
これでペアリングが完全であることの証明が完了する。

\medskip\noindent
$t$ が $H^0(X, \mathcal{O}_X)$ の単元による乗法を前合成する違いを除いて一意であることを
示す議論は省略する。
\end{proof}







\section{チャーン類}
\label{derham-section-chern-classes}

\noindent
ここまでに証明した結果を用いれば、Weil Cohomology Theories,
Section \ref{weil-section-chern} の議論からド・ラームコホモロジーにおける
チャーン類を構成できる。

\begin{lemma}
\label{derham-lemma-chern-classes}
任意の準コンパクトかつ準分離なスキーム $X$ に全チャーン類
$$
c^{dR} :
K_0(\textit{Vect}(X))
\longrightarrow
\prod\nolimits_{i \geq 0} H^{2i}_{dR}(X/\mathbf{Z})
$$
を対応させ、次の性質を満たす規則が一意に存在する。
\begin{enumerate}
\item $c^{dR}(\alpha + \beta) = c^{dR}(\alpha) c^{dR}(\beta)$ が
$\alpha, \beta \in K_0(\textit{Vect}(X))$ に対して成り立つ。
\item $f : X \to X'$ が準コンパクトかつ準分離なスキームの射ならば、
$c^{dR}(f^*\alpha) = f^*c^{dR}(\alpha)$ である。
\item $\mathcal{L} \in \Pic(X)$ が与えられたとき、
$c^{dR}([\mathcal{L}]) = 1 + c_1^{dR}(\mathcal{L})$ である。
\end{enumerate}
\end{lemma}

\noindent
この構成は容易にすべてのスキームへ拡張できるが、そのためには
Weil Cohomology Theories, Section \ref{weil-section-chern} の議論を
少し強化する必要がある。

\begin{proof}
Weil Cohomology Theories, Proposition
\ref{weil-proposition-chern-class} を適用してこれを得る。

\medskip\noindent
$\mathcal{C}$ をすべての準コンパクトかつ準分離なスキームからなる圏とする。
これは明らかに Weil Cohomology Theories, Section \ref{weil-section-chern} の
条件 (1)、(2)、および (3) (a)、(b)、(c) を満たす。

\medskip\noindent
次数付き代数の圏への反変関手 $A$ は $\mathcal{C}$ 上で $X$ を
$A(X) = \bigoplus_{i \geq 0} H_{dR}^{2i}(X/\mathbf{Z})$
へ送り、これにはカップ積を入れる。関手性は
Section \ref{derham-section-de-rham-cohomology} で、カップ積は
Section \ref{derham-section-cup-product} で論じた。
加法的写像 $c_1^A$ として Section \ref{derham-section-first-chern-class} で構成した
$c_1^{dR}$ を取る。

\medskip\noindent
実際、Lemma \ref{derham-lemma-cup-product-graded-commutative} により可換代数を得る。
したがって $A$ に対する公理 (1) が成り立つ。

\medskip\noindent
$A$ に対する公理 (2) を確認するには、
$H^*_{dR}(X \coprod Y/\mathbf{Z}) =  H^*_{dR}(X/\mathbf{Z}) \times
H^*_{dR}(Y/\mathbf{Z})$.
を確認すれば十分である。これは、ド・ラームコホモロジーが
（Zariski 位相における）微分次数付き代数の層のコホモロジーを取ることで
構成されることから従う。

\medskip\noindent
$A$ に対する公理 (3) は、可逆加群の第一チャーン類を取る操作が引き戻しと
両立するという主張にほかならない。これはより一般的な
Lemma \ref{derham-lemma-pullback-c1} から従う。

\medskip\noindent
$A$ に対する公理 (4) は、Proposition
\ref{derham-proposition-projective-space-bundle-formula} で証明した射影空間束公式である。

\medskip\noindent
公理 (5)。$X$ を準コンパクトかつ準分離なスキームとし、
$\mathcal{E} \to \mathcal{F}$ を有限局所自由 $\mathcal{O}_X$-加群の間の全射で、
それぞれの階数が $r + 1$ および $r$ であるものとする。対応する包含射を
$i : P' = \mathbf{P}(\mathcal{F}) \to \mathbf{P}(\mathcal{E}) = P$ と書く。
これは $X$ 上の滑らかな射影スキームの射であり、$P'$ を $P$ 上の
有効 Cartier 因子として表す。したがって Lemma \ref{derham-lemma-check-log-smooth} により、
$P' \subset P$ に対する $\mathbf{Z}$ 上の対数極複体が定義される。
ゆえに $a \in A(P)$ が $i^*a = 0$ を満たすならば、Lemma
\ref{derham-lemma-log-complex-consequence} により
$a \cup c_1^A(\mathcal{O}_P(P')) = 0$ である。これで証明は完了する。
\end{proof}

\begin{remark}
\label{derham-remark-splitting-principle}
Weil Cohomology Theories, Lemmas
\ref{weil-lemma-splitting-principle}（分裂原理）および
\ref{weil-lemma-chern-classes-E-tensor-L}（テンソル積のチャーン類）の類似は、
準コンパクトかつ準分離なスキーム上のド・ラームチャーン類に対して成り立つ。
これは明らかである。実際、Lemma \ref{derham-lemma-chern-classes} の証明で
Weil Cohomology Theories, Section \ref{weil-section-chern} のすべての公理が
満たされることを示した。
\end{remark}

\noindent
$\mathbf{Q}$ 上のスキームを扱えば、チャーン指標を構成できる。

\begin{lemma}
\label{derham-lemma-chern-character}
任意の準コンパクトかつ準分離なスキーム $X$ で $\mathbf{Q}$ 上のものに「チャーン指標」
$$
ch^{dR} : K_0(\textit{Vect}(X)) \longrightarrow
\prod\nolimits_{i \geq 0} H_{dR}^{2i}(X/\mathbf{Q})
$$
を対応させ、次の性質を満たす規則が一意に存在する。
\begin{enumerate}
\item $ch^{dR}$ はすべての $X$ に対して環準同型である。
\item $f : X' \to X$ が $\mathbf{Q}$ 上の準コンパクトかつ準分離なスキームの射ならば、
$f^* \circ ch^{dR} =  ch^{dR} \circ f^*$ である。
\item $\mathcal{L} \in \Pic(X)$ が与えられたとき、
$ch^{dR}([\mathcal{L}]) = \exp(c_1^{dR}(\mathcal{L}))$ である。
\end{enumerate}
\end{lemma}

\noindent
この構成は容易にすべての $\mathbf{Q}$ 上のスキームへ拡張できるが、そのためには
Weil Cohomology Theories, Section \ref{weil-section-chern} の議論を
少し強化する必要がある。

\begin{proof}
Lemma \ref{derham-lemma-chern-classes} の証明とまったく同様に、
$\mathbf{Q}$ 上の準コンパクトかつ準分離なスキームの圏と関手
$A^*(X) = \bigoplus_{i \geq 0} H_{dR}^{2i}(X/\mathbf{Q})$
が Weil Cohomology Theories, Section \ref{weil-section-chern} の公理を
満たすことを示せる。さらにこの場合、$A(X)$ は $\mathbf{Q}$-代数であり、
これはすべての $X$ に対して成り立つ。したがって補題は
Weil Cohomology Theories, Proposition
\ref{weil-proposition-chern-character} から従う。
\end{proof}







\section{Weil コホモロジー理論}
\label{derham-section-weil}

\noindent
$k$ を標数 $0$ の体とする。この節では関手
$$
X \longmapsto H^*_{dR}(X/k)
$$
が、Weil Cohomology Theories, Definition
\ref{weil-definition-weil-cohomology-theory} で定義された意味で、
$k$ 上の係数体 $k$ の Weil コホモロジー理論を定めることを証明する。
この章の前の構成が、Weil Cohomology Theories, Section
\ref{weil-section-c1} の公理 (A1)--(A9) を満たすデータ (D0)、(D1)、(D2') を
与えることを順に確認する。

\medskip\noindent
この節の残りでは体 $k$ で標数 $0$ のものを固定し、$F = k$ とおく。
次のデータを取る。
\begin{enumerate}
\item[(D0)] $1$-次元 $F$-ベクトル空間 $F(1)$ として
$F(1) = F = k$ を取る。
\item[(D1)] 関手 $H^*$ として、滑らかな射影スキーム
$X$ で $k$ 上のものを $H^*_{dR}(X/k)$ に送る関手を取る。
関手性は Section \ref{derham-section-de-rham-cohomology} で、カップ積は
Section \ref{derham-section-cup-product} で論じた。
Lemma \ref{derham-lemma-cup-product-graded-commutative} により次数付き可換
$F$-代数を得る。
\item[(D2')] 写像 $c_1^H : \Pic(X) \to H^2(X)(1)$ として、
Section \ref{derham-section-first-chern-class} で導入したド・ラーム第一チャーン類を用いる。
\end{enumerate}
公理 (A1)--(A9) が成り立つことを示す。

\medskip\noindent
この段落では Weil Cohomology Theories, Lemma
\ref{weil-lemma-check-over-extension} を用いて、公理の確認を $k$ が
代数閉である場合へ帰着する。$k'$ を $k$ の代数閉包と書く。
$F' = k'$ とおく。上とまったく同様に、$k'$ 上で係数体 $F'$ の
データ (D0)、(D1)、(D2') を得る。Lemma
\ref{derham-lemma-proper-smooth-de-Rham} により関手的同型
$$
H_{dR}^{2d}(X/k) \otimes_k k'
\longrightarrow
H_{dR}^{2d}(X_{k'}/k')
$$
が、滑らかかつ射影的な $X$ で $k$ 上のものに対して存在する。さらに図式
$$
\xymatrix{
\Pic(X) \ar[r]_{c^{dR}_1} \ar[d] & H_{dR}^2(X/k) \ar[d] \\
\Pic(X_{k'}) \ar[r]^{c^{dR}_1} & H_{dR}^2(X_{k'}/k')
}
$$
は Lemma \ref{derham-lemma-pullback-c1} により可換である。これで帰着の証明は完了する。

\medskip\noindent
$k$ は標数零の代数閉体であると仮定する。上で与えたデータ
(D0)、(D1)、(D2') に対して公理 (A1)--(A9) を示す。

\medskip\noindent
公理 (A1)。ここでは
$H^*_{dR}(X \coprod Y/k) =  H^*_{dR}(X/k) \times H^*_{dR}(Y/k)$.
を確認しなければならない。これはド・ラームコホモロジーが
（Zariski 位相における）微分次数付き代数の層のコホモロジーを取ることで
構成されることから従う。

\medskip\noindent
公理 (A2)。これは可逆加群の第一チャーン類を取る操作が引き戻しと
両立するという主張にほかならない。これはより一般的な
Lemma \ref{derham-lemma-pullback-c1} から従う。

\medskip\noindent
公理 (A3)。これはより一般的な Proposition
\ref{derham-proposition-projective-space-bundle-formula} から従う。

\medskip\noindent
公理 (A4)。これはより一般的な Lemma
\ref{derham-lemma-log-complex-consequence} から従う。

\medskip\noindent
この時点ですでに Weil Cohomology Theories, Lemmas
\ref{weil-lemma-chern-classes} および
\ref{weil-lemma-cycle-classes} を用いると、チャーン指標とサイクル類写像
$$
\gamma :
\CH^*(X)
\longrightarrow
\bigoplus\nolimits_{i \geq 0} H^{2i}_{dR}(X/k)
$$
を得る。これは滑らかかつ射影的な $X$ で $k$ 上のものに対して次数付き環準同型であり、
滑らかな射影スキームの射 $f : X \to Y$ で $k$ 上のものによる引き戻しと両立する。

\medskip\noindent
公理 (A5)。$H_{dR}^*(\Spec(k)/k) = k = F$ は次数 $0$ に集中する。
Lemma \ref{derham-lemma-kunneth-de-rham} により二つの滑らかな射影
$k$-スキームの積に対する K\"unneth 公式がある
（体 $k$ 上でテンソル積を取るので、主張中の導来テンソル積に問題はない）。

\medskip\noindent
公理 (A7)。これは Proposition \ref{derham-proposition-blowup-split} から従う。

\medskip\noindent
公理 (A8)。$X$ を $k$ 上の滑らかな射影スキームとする。
Weil Cohomology Theories, Section \ref{weil-section-c1} にあるこの公理の
説明により、$k' = H^0(X, \mathcal{O}_X)$ は有限分離 $k$-代数である。
$H_{dR}^*(\Spec(k')/k) = k'$ は次数 $0$ に集中する。これは
$\Omega_{k'/k} = 0$ だからである。
Lemma \ref{derham-lemma-bottom-part-degenerates} により
$H_{dR}^0(X, \mathcal{O}_X) = k'$ も成り立つので、公理を得る。

\medskip\noindent
公理 (A6)。$X$ を $k$ 上の空でない滑らかな射影スキームで、
次元 $d$ の等次元なものとする。
$\Delta : X \to X \times_{\Spec(k)} X$ を $X$ の $k$ 上の対角射と書く。
$k$-線形写像
$$
\lambda : H_{dR}^{2d}(X/k) \longrightarrow k
$$
で $(1 \otimes \lambda)\gamma([\Delta]) = 1$ が $H^0_{dR}(X/k)$ で成り立つものが
存在することを示さなければならない。次のように書く。
$$
\gamma = \gamma([\Delta]) = \gamma_0 + \ldots + \gamma_{2d}
$$
ここで $\gamma_i \in H_{dR}^i(X/k) \otimes_k H_{dR}^{2d - i}(X/k)$ は
K\"unneth 成分である。問題は、線形写像
$\lambda : H_{dR}^{2d}(X/k) \to k$ で
$(1 \otimes \lambda)\gamma_0 = 1$ が $H^0_{dR}(X/k)$ において成り立つものの
存在を示すことである。

\medskip\noindent
$X = \coprod X_i$ を $X$ の連結、したがって既約な成分への分解とする。
これに対応して $\Delta = \coprod \Delta_i$ かつ
$\Delta_i \subset X_i \times X_i$ である。したがって
$$
\gamma([\Delta]) = \sum \gamma([\Delta_i])
$$
であり、さらに $\gamma([\Delta_i])$ は
$\Delta_i \subset X_i \times X_i$ の類に対応する。その際の分解は
$$
H^*_{dR}(X \times X) = \prod\nolimits_{i, j} H^*_{dR}(X_i \times X_j)
$$
である。細部は省略する。これを示す一つの方法は、$\CH^0(X \times X)$ に
冪等元 $e_{i, j}$ があり、それらが開かつ閉な部分スキーム
$X_i \times X_j$ に対応し、
$\gamma$ が $e_{i, j}$ を表示したコホモロジーの直積分解における対応する
冪等元へ送る環準同型であることを用いることである。
$\lambda_i : H_{dR}^{2d}(X_i/k) \to k$ で
$(1 \otimes \lambda_i)\gamma([\Delta_i]) = 1$ が $H^0_{dR}(X_i/k)$ で
成り立つものを見つけられれば、$\lambda = \sum \lambda_i$ と取ることで
$X$ に対する問題が解ける。したがって $X$ は既約であると仮定してよい。

\medskip\noindent
$X$ が既約な場合の公理 (A6) の証明。$k$ は代数閉であるから
$H^0_{dR}(X/k) = k$ である。実際、$H^0(X, \mathcal{O}_X) = k$ である。
これは $X$ が代数閉体上の射影多様体だからである（例えば Varieties, Lemma
\ref{varieties-lemma-proper-geometrically-reduced-global-sections} を参照）。
$x \in X$ を任意の閉点とする。カルテシアン図式
$$
\xymatrix{
x \ar[d] \ar[r] & X \ar[d]^\Delta \\
X \ar[r]^-{x \times \text{id}} & X \times_{\Spec(k)} X
}
$$
を考える。$\gamma$ と引き戻しの両立性から、$\gamma([\Delta])$ は
$\gamma([x])$ へ写る。これは $H_{dR}^{2d}(X/k)$ における主張である。言い換えると
$\gamma_0 = 1 \otimes \gamma([x])$ である。ここから二つのことを得る。
(a) 類 $\gamma([x])$ は $x$ に依存しない。(b) 類 $\gamma([x])$ が零でないことを
示せば十分である。したがって (c) 零サイクル $\alpha$ で $X$ 上のもののうち、
$\gamma(\alpha) \not = 0$ を満たすものを一つ見つければ十分である。
このため有限射
$$
f : X \longrightarrow \mathbf{P}^d_k
$$
を選ぶ。このような射が存在することについては Intersection Theory,
Section \ref{intersection-section-projection}、特に
Lemma \ref{intersection-lemma-projection-generically-finite} を参照せよ。
$f$ は有限シントミックであることに注意する
（More on Morphisms, Lemma \ref{more-morphisms-lemma-lci-permanence} により
局所完全交叉射であり、Algebra, Lemma
\ref{algebra-lemma-CM-over-regular-flat} により平坦である）。
Proposition \ref{derham-proposition-Garel} によりトレース写像
$$
\Theta_f :
f_*\Omega^\bullet_{X/k}
\longrightarrow
\Omega^\bullet_{\mathbf{P}^d_k/k}
$$
があり、標準写像
$$
\Omega^\bullet_{\mathbf{P}^d_k/k}
\longrightarrow
f_*\Omega^\bullet_{X/k}
$$
との合成は $f$ の次数による乗法である。したがって写像
$$
\Theta : H_{dR}^{2d}(X/k) \to H_{dR}^{2d}(\mathbf{P}^d_k/k)
$$
で、$\Theta \circ f^*$ が正の整数による乗法となるものを得る。
したがって $\mathbf{P}^d_k$ 上に類が零でない零サイクルを見つけられれば、
$\gamma$ と引き戻しの両立性により結論を得る。これは
Lemma \ref{derham-lemma-de-rham-cohomology-projective-space} により成り立ち、
公理 (A6) の証明が完了する。

\medskip\noindent
以下では、次の事実を断りなく用いる。第一に、Weil Cohomology Theories,
Remark \ref{weil-remark-trace} により写像
$\lambda_X : H^{2d}_{dR}(X/k) \to k$ は一意である。
第二に、公理 (A6) の証明において、$\lambda_X(\gamma([x])) = 1$ が
$X$ の既約な場合に成り立つことを見た。すなわち、サイクル類写像
$\gamma : \CH^d(X) \to H_{dR}^{2d}(X/k)$ と $\lambda_X$ の合成は
次数写像である。

\medskip\noindent
公理 (A9)。$Y \subset X$ を、空でない滑らかな等次元射影スキーム $X$ で
$k$ 上かつ次元 $d$ のものにおける、空でない滑らかな因子とする。図式
$$
\xymatrix{
H_{dR}^{2d - 2}(X/k)
\ar[rrr]_{c^{dR}_1(\mathcal{O}_X(Y)) \cap -} \ar[d]_{\text{制限}} & & &
H_{dR}^{2d}(X) \ar[d]^{\lambda_X} \\
H_{dR}^{2d - 2}(Y/k) \ar[rrr]^-{\lambda_Y} & & & k
}
$$
が可換であることを示さなければならない。ここで $\lambda_X$ と $\lambda_Y$ は
公理 (A6) におけるものである。上で、$X = \coprod X_i$ を連結成分
（同じことだが既約成分）に分解すると、それに応じて
$\lambda_X = \sum \lambda_{X_i}$ となることを見た。同様に、
$Y = \coprod Y_j$ を連結成分（同じことだが既約成分）に分解すると、
$\lambda_Y = \sum \lambda_{Y_j}$ となる。さらに、この場合
$\mathcal{O}_X(Y) = \otimes_j \mathcal{O}_X(Y_j)$ であり、したがって
$$
c_1^{dR}(\mathcal{O}_X(Y)) = \sum\nolimits_j
c^{dR}_1(\mathcal{O}_X(Y_j))
$$
が $H_{dR}^2(X/k)$ において成り立つ。直接的な図式追跡により、$X$ と $Y$ が
ともに既約な場合に図式の可換性を証明すれば十分である。Weil Cohomology Theories,
Lemma \ref{weil-lemma-poincare-duality} により、そのとき
$H_{dR}^{2d - 2}(Y/k)$ は $1$ 次元である。実際、$Y$ に対する
Poincar\'e 双対性がある。
公理 (A4) により、制限写像（左の垂直矢印）の核は
$c^{dR}_1(\mathcal{O}_X(Y))$ とのカップ積の核に含まれる。
したがって、コホモロジー類 $a \in H_{dR}^{2d - 2}(X)$ を一つ見つけ、
その $Y$ への制限が零でなく、その $a$ に対して図式が可換であることを
示せば十分である。任意の豊富な可逆加群 $\mathcal{L}$ を取り、
$$
a = c^{dR}_1(\mathcal{L})^{d - 1}
$$
と置く。このとき $a|_Y = c^{dR}_1(\mathcal{L}|_Y)^{d - 1}$ であり、したがって
$$
\lambda_Y(a|_Y) = \deg(c_1(\mathcal{L}|_Y)^{d - 1} \cap [Y])
$$
であることが、上の $\lambda_Y$ の記述から分かる。これは Chow Homology,
Lemma \ref{chow-lemma-degrees-and-numerical-intersections} と
Varieties, Lemma \ref{varieties-lemma-ample-positive} を合わせれば正の整数である。
同様に
$$
\lambda_X(c^{dR}_1(\mathcal{O}_X(Y)) \cap a) =
\deg(c_1(\mathcal{O}_X(Y)) \cap c_1(\mathcal{L})^{d - 1} \cap [X])
$$
$c_1(\mathcal{O}_X(Y)) \cap [X] = [Y]$ は本質的に定義から分かるので、零サイクルの等式
$$
(Y \to X)_*\left(c_1(\mathcal{L}|_Y)^{d - 1} \cap [Y]\right) =
c_1(\mathcal{O}_X(Y)) \cap c_1(\mathcal{L})^{d - 1} \cap [X]
$$
が $X$ 上で成り立つ。したがって、これらのサイクルは同じ次数をもち、証明が完了する。

\begin{proposition}
\label{derham-proposition-de-rham-is-weil}
$k$ を標数零の体とする。滑らかな射影スキーム $X$ で $k$ 上のものを
$H_{dR}^*(X/k)$ に送る函子は、Weil Cohomology Theories, Definition
\ref{weil-definition-weil-cohomology-theory} の意味で Weil コホモロジー理論である。
\end{proposition}

\begin{proof}
上の議論において、データ (D0), (D1), (D2') が Weil Cohomology Theories,
Section \ref{weil-section-c1} の公理 (A1) -- (A9) を満たすことを示した。
したがって Weil Cohomology Theories, Proposition
\ref{weil-proposition-get-weil} により結論を得る。

\medskip\noindent
以下は読まないでいただきたい。上記の主張の証明では、さらに Lemmas
\ref{derham-lemma-proper-smooth-de-Rham},
\ref{derham-lemma-pullback-c1},
\ref{derham-lemma-log-complex-consequence},
\ref{derham-lemma-kunneth-de-rham},
\ref{derham-lemma-bottom-part-degenerates},
\ref{derham-lemma-de-rham-cohomology-projective-space},
Propositions
\ref{derham-proposition-projective-space-bundle-formula},
\ref{derham-proposition-blowup-split}, および
\ref{derham-proposition-Garel},
Weil Cohomology Theories, Lemmas
\ref{weil-lemma-check-over-extension},
\ref{weil-lemma-chern-classes},
\ref{weil-lemma-cycle-classes}, および
\ref{weil-lemma-poincare-duality},
Weil Cohomology Theories, Remark \ref{weil-remark-trace},
Varieties, Lemmas
\ref{varieties-lemma-proper-geometrically-reduced-global-sections} および
\ref{varieties-lemma-ample-positive},
Intersection Theory, Section \ref{intersection-section-projection} および
Lemma \ref{intersection-lemma-projection-generically-finite},
More on Morphisms, Lemma \ref{more-morphisms-lemma-lci-permanence},
Algebra, Lemma \ref{algebra-lemma-CM-over-regular-flat}, および
Chow Homology, Lemma
\ref{chow-lemma-degrees-and-numerical-intersections} も用いた。
\end{proof}

\begin{remark}
\label{derham-remark-hodge-cohomology-is-weil}
上とまったく同じ方法で、Hodge コホモロジー $X \mapsto H_{Hodge}^*(X/k)$ が
$c_1^{Hodge}$ を備えると Weil コホモロジー理論を定めることを示せる。
これが必要になれば、ここで正確に定式化して証明する予定である。
そこから次の興味深い帰結が得られる。Weil コホモロジー理論の Betti 数が
（体 $k$ の標数が $0$ であるとして）選んだ Weil コホモロジー理論に依存しないならば、
Hodge-to-de Rham スペクトル系列は $E_1$ で退化する！もちろん、
Hodge-to-de Rham スペクトル系列の退化は知られている
（例えば見事な代数的証明について \cite{Deligne-Illusie} を参照）が、
決して容易な結果ではない。このことは Betti 数の独立性を証明することも
難しい問題であることを示唆しており、知る限りでは今なお未解決である。
関連する問題については Weil Cohomology Theories, Remark
\ref{weil-remark-betti-numbers-in-some-sense} を参照せよ。
\end{remark}







\section{閉埋め込みに対する Gysin 写像}
\label{derham-section-gysin}

\noindent
本節では閉埋め込みに対する Gysin 写像を定義する。

\begin{remark}
\label{derham-remark-gysin-equations}
$X \to S$ をスキームの射とする。また
$f_1, \ldots, f_c \in \Gamma(X, \mathcal{O}_X)$ とする。
$Z \subset X$ を $f_1, \ldots, f_c$ で切り出される閉部分スキームとする。
以下では {\it Gysin 写像}
\begin{equation}
\label{derham-equation-gysin}
\gamma^p_{f_1, \ldots, f_c} :
\Omega^p_{Z/S}
\longrightarrow
\mathcal{H}_Z^c(\Omega^{p + c}_{X/S})
\end{equation}
を考える。これは次のように定義される。局所切断 $\omega$ で
$\Omega^p_{Z/S}$ のものが、切断 $\tilde \omega$ で $\Omega^p_{X/S}$ のものの
制限であるとき、
$$
\gamma^p_{f_1, \ldots, f_c}(\omega) =
c_{f_1, \ldots, f_c}(\tilde \omega|_Z) \wedge
\text{d}f_1 \wedge \ldots \wedge \text{d}f_c
$$
と置く。ここで $c_{f_1, \ldots, f_c} : \Omega^p_{X/S} \otimes \mathcal{O}_Z \to
\mathcal{H}_Z^c(\Omega^p_{X/S})$ は Derived Categories of Schemes, Remark
\ref{perfect-remark-supported-map-c-equations} で構成された写像である。
これは良定義である。実際、$\omega$ を固定して $\tilde \omega$ の選び方を
$\sum f_i \omega'_i + \sum \text{d}(f_i) \wedge \omega''_i$
の形の元だけ変えても、この構成では零へ写る。
\end{remark}

\begin{lemma}
\label{derham-lemma-gysin-differential}
Gysin 写像 (\ref{derham-equation-gysin}) は $\Omega^\bullet_{X/S}$ および
$\Omega^\bullet_{Z/S}$ 上の de Rham 微分と両立する。
\end{lemma}

\begin{proof}
正しく解釈すれば、ほとんど自明な計算から従う。まず、
$\mathcal{H}^c_Z$ を $\mathcal{O}_X$-加群の圏上で計算して得られる函子は、
$X$ 上のアーベル層の圏上で同様に定義した函子と一致する。Cohomology,
Lemma \ref{cohomology-lemma-sections-support-abelian-unbounded} を参照せよ。
したがって微分 $\text{d} : \Omega^p_{X/S} \to \Omega^{p + 1}_{X/S}$ は写像
$\mathcal{H}^c_Z(\Omega^p_{X/S}) \to \mathcal{H}^c_Z(\Omega^{p + 1}_{X/S})$ を誘導する。
さらに Derived Categories of Schemes, Remark
\ref{perfect-remark-support-c-equations} の拡張交代 {\v C}ech 複体の構成は、
アーベル層の圏上でも機能する。写像
$$
\Coker\left(\bigoplus \mathcal{F}_{1 \ldots \hat i \ldots c} \to
\mathcal{F}_{1 \ldots c}\right)
\longrightarrow
i_*\mathcal{H}^c_Z(\mathcal{F})
$$
は $c_{f_1, \ldots, f_c}$ の構成で用いられるものである。Derived Categories of
Schemes, Remark \ref{perfect-remark-supported-map-c-equations} を参照せよ。
これは $X$ 上のすべてのアーベル層の圏で良定義かつ函子的である。
したがって、本補題は等式
$$
\text{d}\left(
\frac{\tilde \omega \wedge \text{d}f_1 \wedge \ldots \wedge
\text{d}f_c}{f_1 \ldots f_c}\right) =
\frac{\text{d}(\tilde \omega) \wedge
\text{d}f_1 \wedge \ldots \wedge \text{d}f_c}{f_1 \ldots f_c}
$$
から従い、これは明らかである。
\end{proof}

\begin{lemma}
\label{derham-lemma-gysin-global}
$X \to S$ をスキームの射とする。$Z \to X$ を有限表示の閉埋め込みとし、
その余法層 $\mathcal{C}_{Z/X}$ が階数 $c$ の局所自由層であるとする。
このとき標準的な写像
$$
\gamma^p : \Omega^p_{Z/S} \to \mathcal{H}^c_Z(\Omega^{p + c}_{X/S})
$$
が存在し、これは局所的には写像 $\gamma^p_{f_1, \ldots, f_c}$ で与えられる。
Remark \ref{derham-remark-gysin-equations} を参照せよ。
\end{lemma}

\begin{proof}
仮定により、$x \in Z \subset X$ が与えられると、開近傍 $U$ であって $x$ を含み、
$Z$ が $c$ 個の元 $f_1, \ldots, f_c \in \mathcal{O}_X(U)$ で切り出されるものが
存在する。したがって、$f_1, \ldots, f_c$ と $g_1, \ldots, g_c$ が
$\mathcal{O}_X(U)$ に属し、ともに $Z \cap U$ を切り出すならば、写像
$\gamma^p_{f_1, \ldots, f_c}$ と $\gamma^p_{g_1, \ldots, g_c}$ が一致することを
示せば十分である。そのため $U$ を縮小した後、$g_j = \sum a_{ji} f_i$ が、ある
$a_{ji} \in \mathcal{O}_X(U)$ に対して成り立つと仮定してよい。このとき
$c_{f_1, \ldots, f_c} = \det(a_{ji}) c_{g_1, \ldots, g_c}$ である。
Derived Categories of Schemes, Lemma
\ref{perfect-lemma-supported-map-determinant} を参照せよ。一方、
$$
\text{d}(g_1) \wedge \ldots \wedge \text{d}(g_c) \equiv
\det(a_{ji}) \text{d}(f_1) \wedge \ldots \wedge \text{d}(f_c)
\bmod (f_1, \ldots, f_c)\Omega^c_{X/S}
$$
である。これらの関係を合わせれば、直接的な計算から所望の等式を得る。
\end{proof}

\begin{lemma}
\label{derham-lemma-gysin-differential-global}
$X \to S$ および $i : Z \to X$ は Lemma \ref{derham-lemma-gysin-global} の通りとする。
Gysin 写像 $\gamma^p$ は $\Omega^\bullet_{X/S}$ および
$\Omega^\bullet_{Z/S}$ 上の de Rham 微分と両立する。
\end{lemma}

\begin{proof}
これは局所的に確認でき、その場合 Lemma \ref{derham-lemma-gysin-differential} から従う。
\end{proof}

\begin{lemma}
\label{derham-lemma-gysin-projection}
$X \to S$ および $i : Z \to X$ は Lemma \ref{derham-lemma-gysin-global} の通りとする。
$\alpha \in H^q(X, \Omega^p_{X/S})$ が与えられると、
$\gamma^p(\alpha|_Z) = i^{-1}\alpha \wedge \gamma^0(1)$ が
$H^q(Z, \mathcal{H}^c_Z(\Omega^{p + c}_{X/S}))$ において成り立つ。
記法については証明を参照せよ。
\end{lemma}

\begin{proof}
制限 $\alpha|_Z$ は Hodge コホモロジーの函子性によって得られる
$H^q(Z, \Omega^p_{Z/S})$ の元である。コホモロジーの函子性を写像
$\gamma^p : \Omega^p_{Z/S} \to \mathcal{H}^c_Z(\Omega^{p + c}_{X/S})$
に適用すると、$\gamma^p(\alpha|_Z)$ を
$H^q(Z, \mathcal{H}^c_Z(\Omega^{p + c}_{X/S}))$ の元として得る。
これで公式の左辺が説明される。

\medskip\noindent
右辺を説明するため、まず環付き空間の写像
$i : (Z, i^{-1}\mathcal{O}_X) \to (X, \mathcal{O}_X)$ で引き戻し、元
$i^{-1}\alpha \in H^q(Z, i^{-1}\Omega^p_{X/S})$ を得る。
$\gamma^0(1) \in H^0(Z, \mathcal{H}_Z^c(\Omega^c_{X/S}))$ を、
$1 \in H^0(Z, \mathcal{O}_Z) = H^0(Z, \Omega^0_{Z/S})$ の
$\gamma^0$ による像とする。カップ積を用いて元
$$
i^{-1}\alpha \cup \gamma^0(1)
\in
H^{q + c}(Z, 
i^{-1}\Omega^p_{X/S} \otimes_{i^{-1}\mathcal{O}_X}
\mathcal{H}^c_Z(\Omega^c_{X/S}))
$$
を得る。Cohomology, Remark \ref{cohomology-remark-support-cup-product} と
外積を用いると、標準的な写像
$$
i^{-1}\Omega^p_{X/S} \otimes_{i^{-1}\mathcal{O}_X}^\mathbf{L}
R\mathcal{H}_Z(\Omega^c_{X/S}) \to
R\mathcal{H}_Z(\Omega^p_{X/S} \otimes_{\mathcal{O}_X}^\mathbf{L}
\Omega^c_{X/S}) \to
R\mathcal{H}_Z(\Omega^{p + c}_{X/S})
$$
がある。Derived Categories of Schemes, Lemma
\ref{perfect-lemma-supported-trivial-vanishing} により、対象
$R\mathcal{H}_Z(\Omega^j_{X/S})$ のコホモロジー層は次数 $> c$ で消える。
したがって次数 $c$ のコホモロジー層上で写像
$$
i^{-1}\Omega^p_{X/S} \otimes_{i^{-1}\mathcal{O}_X}
\mathcal{H}^c_Z(\Omega^c_{X/S}) \longrightarrow
\mathcal{H}^c_Z(\Omega^{p + c}_{X/S})
$$
を得る。式 $i^{-1}\alpha \wedge \gamma^0(1)$ は、カップ積
$i^{-1}\alpha \cup \gamma^0(1)$ のコホモロジーの函子性による像である。

\medskip\noindent
以上のように公式の内容を説明したので、カップ積の一般的性質
（Cohomology, Section \ref{cohomology-section-cup-product}）により、図式
$$
\xymatrix{
i^{-1}\Omega^p_X \otimes \Omega^0_Z \ar[rr]_{\text{id} \otimes \gamma^0}
\ar[d] & &
i^{-1}\Omega^p_X \otimes \mathcal{H}^c_Z(\Omega^c_X) \ar[d]^\wedge \\
\Omega^p_Z \otimes \Omega^0_Z \ar[r]^\wedge &
\Omega^p_Z \ar[r]^{\gamma^p} &
\mathcal{H}^c_Z(\Omega^{p + c}_X)
}
$$
が $Z$ 上の層の圏で可換であることを証明すれば十分である
（記法を明らかな仕方で濫用している）。これは写像
$\gamma^j_{f_1, \ldots, f_c}$ に対する簡単な計算に帰着するので省略する。
実際、これらの写像は、まさにこの可換性が成り立ち、かつ $1$ が
$\frac{\text{d}f_1 \wedge \ldots \wedge \text{d}f_c}{f_1 \ldots f_c}$ へ
写るように選ばれている。
\end{proof}

\begin{lemma}
\label{derham-lemma-gysin-transverse}
$c \geq 0$ を整数とする。また
$$
\xymatrix{
Z' \ar[d]_h \ar[r] & X' \ar[d]_g \ar[r] & S' \ar[d] \\
Z \ar[r] & X \ar[r] & S
}
$$
をスキームの可換図式とする。次を仮定する。
\begin{enumerate}
\item $Z \to X$ および $Z' \to X'$ は
Lemma \ref{derham-lemma-gysin-global} の仮定を満たす。
\item 図式の左の正方形はカルテシアンである。
\item $h^*\mathcal{C}_{Z/X} \to \mathcal{C}_{Z'/X'}$
(Morphisms, Lemma \ref{morphisms-lemma-conormal-functorial})
は同型である。
\end{enumerate}
このとき図式
$$
\xymatrix{
h^*\Omega^p_{Z/S} \ar[rr]_-{h^{-1}\gamma^p} \ar[d] & &
\mathcal{O}_{X'}|_{Z'} \otimes_{h^{-1}\mathcal{O}_X|_Z}
h^{-1}\mathcal{H}^c_Z(\Omega^{p + c}_{X/S}) \ar[d] \\
\Omega^p_{Z'/S'} \ar[rr]^{\gamma^p} & &
\mathcal{H}^c_{Z'}(\Omega^{p + c}_{X'/S'})
}
$$
は可換である。左の垂直矢印は微分加群の函子性によるものであり、右の垂直矢印には
Cohomology, Remark \ref{cohomology-remark-support-functorial} を用いる。
\end{lemma}

\begin{proof}
より正確には、合成
\begin{align*}
\mathcal{O}_{X'}|_{Z'} \otimes_{h^{-1}\mathcal{O}_X|_Z}^\mathbf{L}
h^{-1}R\mathcal{H}_Z(\Omega^{p + c}_{X/S})
& \to
R\mathcal{H}_{Z'}(Lg^*\Omega^{p + c}_{X/S}) \\
& \to
R\mathcal{H}_{Z'}(g^*\Omega^{p + c}_{X/S}) \\
& \to
R\mathcal{H}_{Z'}(\Omega^{p + c}_{X'/S'})
\end{align*}
を考える。第一の矢印は Cohomology, Remark
\ref{cohomology-remark-support-functorial} により、最後の矢印は微分の函子性により
与えられる。Derived Categories of Schemes, Lemma
\ref{perfect-lemma-supported-trivial-vanishing} により、コホモロジー層は次数 $> c$ で
消えるので、これは右の垂直矢印を誘導する。可換性は局所的に確認できる。
したがって、$Z$ は $f_1, \ldots, f_c \in \Gamma(X, \mathcal{O}_X)$ により
切り出されると仮定してよい。このとき $Z'$ は $f'_i = g^\sharp(f_i)$ により
切り出される。写像 $c_{f_1, \ldots, f_c}$ および
$c_{f'_1, \ldots, f'_c}$ は可換図式
$$
\xymatrix{
h^*i^*\Omega^p_{X/S} \ar[rr]_-{h^{-1}c_{f_1, \ldots, f_c}} \ar[d] & &
\mathcal{O}_{X'}|_{Z'} \otimes_{h^{-1}\mathcal{O}_X|_Z}
h^{-1}\mathcal{H}^c_Z(\Omega^p_{X/S}) \ar[d] \\
(i')^*\Omega^p_{X'/S'} \ar[rr]^{c_{f'_1, \ldots, f'_c}} & &
\mathcal{H}^c_{Z'}(\Omega^p_{X'/S'})
}
$$
に入る。Derived Categories of Schemes, Remark
\ref{perfect-remark-supported-functorial} を参照せよ。$p$-形式 $\omega$ で
$Z$ 上のものが与えられると、$\gamma^p(\omega)$ は、$X$ および $Z$ 上で局所的に
$p$-形式 $\tilde \omega$ で $X$ 上のものを選んで $\omega$ を持ち上げ、
$\gamma^p(\omega) =
c_{f_1, \ldots, f_c}(\tilde \omega) \wedge
\text{d}f_1 \wedge \ldots \wedge \text{d}f_c$.
と取ることにより定義されることを思い出そう。形式
$\text{d}f_1 \wedge \ldots \wedge \text{d}f_c$ は
$\text{d}f'_1 \wedge \ldots \wedge \text{d}f'_c$ に引き戻されるので結論を得る。
\end{proof}

\begin{remark}
\label{derham-remark-how-to-use}
$X \to S$, $i : Z \to X$, および $c \geq 0$ は
Lemma \ref{derham-lemma-gysin-global} の通りとする。$p \geq 0$ とし、
$\mathcal{H}^i_Z(\Omega^{p + c}_{X/S}) = 0$ が $i = 0, \ldots, c - 1$ に対して
成り立つと仮定する。$X \to S$ が滑らかで $Z \to X$ が Koszul 正則埋め込みならば、
この消滅は成り立つ。Derived Categories of Schemes, Lemma
\ref{perfect-lemma-supported-vanishing} を参照せよ。このとき写像
$$
\gamma^{p, q} :
H^q(Z, \Omega^p_{Z/S})
\longrightarrow
H^{q + c}(X, \Omega^{p + c}_{X/S})
$$
を得る。まず
$\gamma^p : \Omega^p_{Z/S} \to \mathcal{H}^c_Z(\Omega^{p + c}_{X/S})$ を用いて
$$
H^q(Z, \mathcal{H}^c_Z(\Omega^{p + c}_{X/S})) =
H^q(Z, R\mathcal{H}_Z(\Omega^{p + c}_{X/S})[c]) =
H^q(X, i_*R\mathcal{H}_Z(\Omega^{p + c}_{X/S})[c])
$$
へ写し、次に随伴写像
$i_*R\mathcal{H}_Z(\Omega^{p + c}_{X/S}) \to \Omega^{p + c}_{X/S}$ を用いて、
所望の Hodge コホモロジー加群へ進む。
\end{remark}

\begin{lemma}
\label{derham-lemma-gysin-differential-hodge}
$X \to S$ および $i : Z \to X$ は Lemma \ref{derham-lemma-gysin-global} の通りとする。
$X \to S$ は滑らかで、$Z \to X$ は Koszul 正則であると仮定する。
Gysin 写像 $\gamma^{p, q}$ は $\Omega^\bullet_{X/S}$ および
$\Omega^\bullet_{Z/S}$ 上の de Rham 微分と両立する。
\end{lemma}

\begin{proof}
これは Lemma \ref{derham-lemma-gysin-differential-global} から直ちに従う。
\end{proof}

\begin{lemma}
\label{derham-lemma-gysin-projection-global}
$X \to S$, $i : Z \to X$, および $c \geq 0$ は
Lemma \ref{derham-lemma-gysin-global} の通りとする。$X \to S$ は滑らかで
$Z \to X$ は Koszul 正則であると仮定する。
$\alpha \in H^q(X, \Omega^p_{X/S})$ が与えられると、
$\gamma^{p, q}(\alpha|_Z) = \alpha \cup \gamma^{0, 0}(1)$ が
$H^{q + c}(X, \Omega^{p + c}_{X/S})$ において成り立つ。ここで $\gamma^{a, b}$ は
Remark \ref{derham-remark-how-to-use} の通りである。
\end{lemma}

\begin{proof}
本補題は Lemma \ref{derham-lemma-gysin-projection} と Cohomology, Lemma
\ref{cohomology-lemma-support-cup-product} から従う。以下のより詳しい議論は
読み飛ばすことを勧める。

\medskip\noindent
以下では、
$R\mathcal{H}_Z(\Omega^j_{X/S}) = \mathcal{H}^c_Z(\Omega^j_{X/S})[-c]$
がすべての $j$ に対して成り立つことを断りなく用いる。これは
Remark \ref{derham-remark-how-to-use} で指摘した通りである。また同一視
$H^{q + c}_Z(X, \Omega^j_{X/S}) = H^{q + c}(Z, R\mathcal{H}_Z(\Omega^j_{X/S}) =
H^q(Z, \mathcal{H}^c_Z(\Omega^j_{X/S}))$
も暗黙に用いる。最初の同一視については Cohomology, Lemma
\ref{cohomology-lemma-local-to-global-sections-with-support} を参照せよ。
これらの同一視のもとで、
\begin{enumerate}
\item $\gamma^0(1) \in H^c_Z(X, \Omega^c_{X/S})$ は
$\gamma^{0, 0}(1)$ へ写り、これは $H^c(X, \Omega^c_{X/S})$ に属する。
\item Lemma \ref{derham-lemma-gysin-projection} の等式の右辺
$i^{-1}\alpha \wedge \gamma^0(1)$ は、Cohomology, Remark
\ref{cohomology-remark-support-cup-product-global} における元 $\alpha$ と
$\gamma^0(1)$ のカップ積の外積による像である。言い換えると、
Lemma \ref{derham-lemma-gysin-projection} の証明における構成と Cohomology, Remark
\ref{cohomology-remark-support-cup-product-global} の構成は一致する。
\item Cohomology, Lemma \ref{cohomology-lemma-support-cup-product} により、これは
$\alpha \cup \gamma^{0, 0}(1)$ へ写り、これは
$H^{q + c}(X, \Omega^p_{X/S} \otimes \Omega^c_{X/S})$ に属する。
\item Lemma \ref{derham-lemma-gysin-projection} の等式の左辺
$\gamma^p(\alpha|_Z)$ は $\gamma^{p, q}(\alpha|_Z)$ へ写る。
\end{enumerate}
これで証明が完了する。
\end{proof}

\begin{lemma}
\label{derham-lemma-gysin-transverse-global}
$c \geq 0$ とし、
$$
\xymatrix{
Z' \ar[d]_h \ar[r] & X' \ar[d]_g \ar[r] & S' \ar[d] \\
Z \ar[r] & X \ar[r] & S
}
$$
が Lemma \ref{derham-lemma-gysin-transverse} の仮定を満たすとする。さらに、
$X \to S$ および $X' \to S'$ は滑らかで、$Z \to X$ および $Z' \to X'$ は
Koszul 正則埋め込みであると仮定する。このとき図式
$$
\xymatrix{
H^q(Z, \Omega^p_{Z/S}) \ar[rr]_-{\gamma^{p, q}} \ar[d] & &
H^{q + c}(X, \Omega^{p + c}_{X/S}) \ar[d] \\
H^q(Z', \Omega^p_{Z'/S'}) \ar[rr]^{\gamma^{p, q}} & &
H^{q + c}(X', \Omega^{p + c}_{X'/S'})
}
$$
は可換である。ここで $\gamma^{p, q}$ は Remark \ref{derham-remark-how-to-use} の通りである。
\end{lemma}

\begin{proof}
これは Lemma \ref{derham-lemma-gysin-transverse} と Cohomology, Lemma
\ref{cohomology-lemma-support-functorial} を合わせれば従う。
\end{proof}

\begin{lemma}
\label{derham-lemma-class-of-a-point}
$k$ を体とする。$X$ を $k$ 上の既約で滑らかな固有スキームで、次元 $d$ の
ものとする。$Z \subset X$ をただ一つの $k$-有理点 $x$ からなる被約閉部分スキームと
する。このとき、元 $1 \in k = H^0(Z, \mathcal{O}_Z) = H^0(Z, \Omega^0_{Z/k})$ の
Remark \ref{derham-remark-how-to-use} の写像
$H^0(Z, \Omega^0_{Z/k}) \to H^d(X, \Omega^d_{X/k})$ による像は零でない。
\end{lemma}

\begin{proof}
写像 $\gamma^0 : \mathcal{O}_Z \to
\mathcal{H}^d_Z(\Omega^d_{X/k}) = R\mathcal{H}_Z(\Omega^d_{X/k})[d]$ は写像
$$
g^0 : i_*\mathcal{O}_Z \longrightarrow \Omega^d_{X/k}[d]
$$
の $D(\mathcal{O}_X)$ における随伴である。$\Omega^d_{X/k} = \omega_X$ は
$X/k$ の双対化層であることを思い出そう。Duality for Schemes, Lemma
\ref{duality-lemma-duality-proper-over-field} を参照せよ。したがって、本補題の
主張にある写像の $k$-線形双対は写像
$$
H^0(X, \mathcal{O}_X) \to \Ext^d_X(i_*\mathcal{O}_Z, \omega_X)
$$
であり、$1$ を $g^0$ へ送る。したがって $g^0$ が零でないことを示せば十分である。
これは任意の近傍 $U$ で点 $x$ を含むものにおいて確認してよい。そのような $U$ を、
$f_1, \ldots, f_d \in \mathcal{O}_X(U)$ で $x$ においてのみ消え、極大イデアル
$\mathfrak m_x \subset \mathcal{O}_{X, x}$ を生成するものが存在するように選ぶ。
$U = \Spec(R)$ はアフィンであると仮定してよい。$\gamma^0$ の構成を見直すと、
この拡大は
$$
k \to
(R \to \bigoplus\nolimits_{i_0} R_{f_{i_0}} \to
\bigoplus\nolimits_{i_0 < i_1} R_{f_{i_0}f_{i_1}} \to
\ldots \to R_{f_1\ldots f_r})[d] \to R[d]
$$
で与えられ、第一の写像で $1$ は $1/f_1 \ldots f_c$ へ写る。
これは零でない。実際、この場合 $1/f_1 \ldots f_c$ は局所コホモロジー群
$H^d_{(f_1, \ldots, f_d)}(R)$ の零でない元である。
\end{proof}






\section{相対 Poincar\'e 双対性}
\label{derham-section-relative-poincare-duality}

\noindent
本節では、底上の固有滑らかなスキームの相対 de Rham コホモロジーに対する
Poincar\'e 双対性を証明する。まず Section \ref{derham-section-poincare-duality}
を参照することを強く勧める。

\begin{situation}
\label{derham-situation-relative-duality}
ここで $S$ は準コンパクトかつ準分離なスキームであり、
$f : X \to S$ はスキームの固有滑らかな射で、そのすべての
ファイバーは空でなく、次元 $n$ の等次元である。
\end{situation}

\begin{lemma}
\label{derham-lemma-relative-bottom-part-degenerates}
Situation \ref{derham-situation-relative-duality} において、順像
$f_*\mathcal{O}_X$ は有限 \'etale $\mathcal{O}_S$-代数であり、
$S$ 上局所的に $Rf_*\mathcal{O}_X = f_*\mathcal{O}_X \oplus P$ が
$D(\mathcal{O}_S)$ において成り立つ。ここで
$P$ は完全で、その Tor 振幅は $[1, \infty)$ に含まれる。
写像 $\text{d} : f_*\mathcal{O}_X \to f_*\Omega_{X/S}$ は零である。
\end{lemma}

\begin{proof}
主張の前半は Derived Categories of Schemes, Lemma
\ref{perfect-lemma-proper-flat-geom-red} から従う。
$S' = \underline{\Spec}_S(f_*\mathcal{O}_X)$ と置くと、分解
$X \to S' \to S$ を得る（これは Stein 分解である。More on Morphisms,
Section \ref{more-morphisms-section-stein-factorization} を参照せよ。ただし、
ここではこの事実を必要としない）。さらに、たとえば Morphisms, Lemma
\ref{morphisms-lemma-triangle-differentials} および
\ref{morphisms-lemma-etale-at-point} により
$\Omega_{X/S} = \Omega_{X/S'}$ であることが分かる。したがって当然、
$\text{d} : f_*\mathcal{O}_X \to f_*\Omega_{X/S}$ は零である。
\end{proof}

\begin{lemma}
\label{derham-lemma-relative-duality-hodge}
Situation \ref{derham-situation-relative-duality} において、
$\mathcal{O}_S$-加群の写像
$$
t : Rf_*\Omega^n_{X/S}[n] \longrightarrow \mathcal{O}_S
$$
で次の性質をもつものが存在し、$H^0(X, \mathcal{O}_X)$ の単元による
乗法写像を前合成することを除いて一意である。すなわち、すべての $p$ に対し、
$$
Rf_*\Omega^p_{X/S}
\otimes_{\mathcal{O}_S}^\mathbf{L}
Rf_*\Omega^{n - p}_{X/S}[n]
\longrightarrow
\mathcal{O}_S
$$
で与えられる、相対カップ積と $t$ との合成は、$S$ 上の完全複体の
完全なペアリングである。
\end{lemma}

\begin{proof}
まず $S$ が Noether スキームであると仮定する。
Duality for Schemes, Remark \ref{duality-remark-relative-dualizing-complex}
におけるように、$\omega^\bullet_{X/S}$ を $X$ の $S$ 上の相対双対化複体とし、
$Rf_*\omega_{X/S}^\bullet \to \mathcal{O}_S$ をその跡写像とする。
Duality for Schemes, Lemma \ref{duality-lemma-smooth-proper}
（ここで $S$ が Noether であることを用いる）により、同型
$\omega^\bullet_{X/S} \cong \Omega^n_{X/S}[n]$ が存在し、この同型を用いて
$t$ を得る。Lemma \ref{derham-lemma-proper-smooth-de-Rham} により、複体
$Rf_*\Omega^p_{X/S}$ は完全である。
$\Omega^p_{X/S}$ は局所自由であり、さらに写像
$\Omega^p_{X/S} \otimes_{\mathcal{O}_X} \Omega^{n - p}_{X/S} \to
\Omega^n_{X/S}$ が同型 $\Omega^p_{X/S} \cong
\SheafHom_{\mathcal{O}_X}(\Omega^{n - p}_{X/S}, \Omega^n_{X/S})$
を与えるので、相対カップ積によって誘導されるペアリングは
Duality for Schemes, Remark
\ref{duality-remark-relative-dualizing-complex-relative-cup-product} により
完全であることが分かる。

\medskip\noindent
一般の場合における存在の証明。
絶対 Noether 近似により、Cartesian 図式
$$
\xymatrix{
X \ar[r] \ar[d]^f & X_0 \ar[d]^{f_0} \\
S \ar[r]^g & S_0
}
$$
で、$S_0$ が Noether であり、$f_0$ が固有かつ滑らかで、そのすべての
ファイバーが空でなく次元 $n$ の等次元であるものを見いだせる。
Limits, Lemmas \ref{limits-lemma-descend-finite-presentation},
\ref{limits-lemma-descend-smooth}, \ref{limits-lemma-descend-dimension-d},
\ref{limits-lemma-descend-surjective}, および
\ref{limits-lemma-eventually-proper} を参照せよ。補題の主張にあるような写像
$$
t_0 : Rf_{0, *}\Omega^n_{X_0/S_0}[n] \longrightarrow \mathcal{O}_{S_0}
$$
を選ぶ。コホモロジーと基底変換により、複体 $Rf_*(\Omega^p_{X/S})$ は
$Lg^*(Rf_{0, *}\Omega^p_{X_0/S_0})$ に、すべての $p \geq 0$ に対して等しい。
Lemma \ref{derham-lemma-proper-smooth-de-Rham} を参照せよ。特に、$t_0$ の
$g$ による引き戻しから、補題の主張にあるような写像 $t$ を得る。
実際、Cohomology, Lemma
\ref{cohomology-lemma-base-change-cup-product} により、カップ積写像
$$
Rf_*\Omega^p_{X/S}
\otimes_{\mathcal{O}_S}^\mathbf{L}
Rf_*\Omega^{n - p}_{X/S}[n]
\longrightarrow
Rf_*\Omega^n_{X/S}
$$
は $g$ による、$f_0$ に対する対応するカップ積写像の引き戻しである。
したがって、$t_0$ が完全なペアリングを与えるという事実を引き戻すと、
$t$ に対して $S$ 上で同じ性質が得られる。

\medskip\noindent
写像 $t$ の一意性。区別三角形
$f_*\mathcal{O}_X \to Rf_*\mathcal{O}_X \to P \to f_*\mathcal{O}_X[1]$
を選ぶ。Lemma \ref{derham-lemma-relative-bottom-part-degenerates} により、対象 $P$ は
完全で、その Tor 振幅は $[1, \infty)$ に含まれ、この三角形は
$S$ 上局所的に分裂する。したがって
$R\SheafHom_{\mathcal{O}_X}(P, \mathcal{O}_X)$ は完全で、その Tor 振幅は
$(-\infty, -1]$ に含まれる。よって上の双対性から、$S$ 上局所的に
$$
Rf_*\Omega^n_{X/S}[n] \cong
R\SheafHom_{\mathcal{O}_S}(f_*\mathcal{O}_X, \mathcal{O}_S)
\oplus R\SheafHom_{\mathcal{O}_X}(P, \mathcal{O}_X)
$$
であることが分かる。これにより、$R^nf_*\Omega^n_{X/S}$ は有限局所自由であり、
完全な $\mathcal{O}_S$-双線形ペアリング
$$
f_*\mathcal{O}_X \times R^nf_*\Omega^n_{X/S} \longrightarrow \mathcal{O}_S
$$
を $t$ を用いて得る。したがって、任意の $\mathcal{O}_S$-線形写像
$t' : R^nf_*\Omega^n_{X/S} \to \mathcal{O}_S$ は
$t' = t \circ g$ という形であり、ここで
$g \in \Gamma(S, f_*\mathcal{O}_X) = \Gamma(X, \mathcal{O}_X)$ である。
$t'$ がなお完全なペアリングを定めるためには、
$g$ は単元でなければならない。これで証明が完了する。
\end{proof}

\begin{lemma}
\label{derham-lemma-relative-top-part-degenerates}
Situation \ref{derham-situation-relative-duality} において、写像
$\text{d} : R^nf_*\Omega^{n - 1}_{X/S} \to R^nf_*\Omega^n_{X/S}$
は零である。
\end{lemma}

\noindent
Lemma \ref{derham-lemma-top-part-degenerates} の証明で述べたように、この補題は
Lemmas \ref{derham-lemma-relative-duality-hodge} および
\ref{derham-lemma-relative-bottom-part-degenerates} の容易な帰結ではない。

\begin{proof}[$S$ が被約である場合の証明]
$S$ が被約であると仮定する。写像
$\text{d} : R^nf_*\Omega^{n - 1}_{X/S} \to R^nf_*\Omega^n_{X/S}$
は（準連接）$\mathcal{O}_S$-加群の $\mathcal{O}_S$-線形写像であることに注意する。
$\mathcal{O}_S$-加群 $R^nf_*\Omega^n_{X/S}$ は有限局所自由である
（Lemmas \ref{derham-lemma-relative-duality-hodge} および
\ref{derham-lemma-relative-bottom-part-degenerates} により、有限局所自由
$\mathcal{O}_S$-加群 $f_*\mathcal{O}_X$ の双対である）。
$S$ は被約なので、$\text{d}$ の、すべての生成点 $\eta \in S$ における
茎が零であることを示せば十分である。これは、アフィン開集合上の
切断を考え、$\text{d}$ の終域が局所自由であることと、Algebra, Lemma
\ref{algebra-lemma-reduced-ring-sub-product-fields} part (2) を用いれば従う。
$S$ は被約なので $\mathcal{O}_{S, \eta} = \kappa(\eta)$ である。
Algebra, Lemma \ref{algebra-lemma-minimal-prime-reduced-ring} を参照せよ。
したがって $\text{d}_\eta$ は写像
$$
\text{d} :
H^n(X_\eta, \Omega^{n - 1}_{X_\eta/\kappa(\eta)})
\longrightarrow
H^n(X_\eta, \Omega^n_{X_\eta/\kappa(\eta)})
$$
と同一視され、これは Lemma \ref{derham-lemma-top-part-degenerates} により零である。
\end{proof}

\begin{proof}[一般の場合の証明]
問題は $S$ 上平坦局所的であることに注意する。すなわち、$S' \to S$ が
スキームの全射平坦射で、$S'$ への引き戻し後に写像が零ならば、もとの写像も
零である。また、平坦基底変換により、この写像の構成は平坦射による基底変換と
可換である（Cohomology of Schemes, Lemma
\ref{coherent-lemma-flat-base-change-cohomology}）。

\medskip\noindent
More on Morphisms, Theorem
\ref{more-morphisms-theorem-stein-factorization-general} における Stein 分解
$X \to S' \to S$ を考える。Lemma
\ref{derham-lemma-relative-bottom-part-degenerates} により、射
$\pi : S' \to S$ は有限 \'etale である。射 $f : X \to S'$ は、同定理により
固有であり、More on Morphisms, Lemma
\ref{more-morphisms-lemma-smooth-etale-permanence} により滑らかであり、
Stein 分解の定理により幾何学的に連結なファイバーをもつ。
Lemma \ref{derham-lemma-relative-bottom-part-degenerates} の証明で、
$\Omega_{X/S} = \Omega_{X/S'}$ となることを見た。これは $S' \to S$ が
\'etale だからである。
したがって $\Omega^\bullet_{X/S} = \Omega^\bullet_{X/S'}$ である。また
$$
R^qf_*\Omega^p_{X/S} = \pi_*R^qf'_*\Omega^p_{X/S'}
$$
がすべての $p, q$ に対して成り立つ。これは Leray スペクトル系列
（Cohomology, Lemma \ref{cohomology-lemma-relative-Leray}）、$\pi$ が有限、
従ってアフィンであること、および Cohomology of Schemes, Lemma
\ref{coherent-lemma-relative-affine-vanishing} による（もちろん、
$R^qf'_*\Omega^p_{X'/S}$ が準連接であることも用いる）。
したがって補題の写像は、$\pi_*$ を
$\text{d} : R^nf'_*\Omega^{n - 1}_{X/S'} \to R^nf'_*\Omega^n_{X/S'}$
に適用したものである。言い換えると、この補題を証明するには
$f : X \to S$ を $f' : X \to S'$ で置き換え、次段落で扱う場合へ帰着してよい。

\medskip\noindent
$f$ が幾何学的に連結なファイバーをもち、
$f_*\mathcal{O}_X = \mathcal{O}_S$ であると仮定する。
すべての $s \in S$ に対し、\'etale 近傍
$(S', s') \to (S, s)$ で、基底変換 $X' \to S'$（$S$ の基底変換）が
切断をもつものを選べる。More on Morphisms, Lemma
\ref{more-morphisms-lemma-etale-nbhd-dominates-smooth} を参照せよ。
証明冒頭の所見により、これで次段落で扱う場合へ帰着する。

\medskip\noindent
$f$ が幾何学的に連結なファイバーをもち、
$f_*\mathcal{O}_X = \mathcal{O}_S$ であり、$s : S \to X$ が与えられ、
これが $f$ の切断であると仮定する。$S = \Spec(A)$ は
アフィンであると仮定してよく、そう仮定する。写像
$s^* : R\Gamma(X, \mathcal{O}_X) \to R\Gamma(S, \mathcal{O}_S) = A$
は写像 $A \to R\Gamma(X, \mathcal{O}_X)$ の分裂である。したがって
$$
R\Gamma(X, \mathcal{O}_X) = A \oplus P
$$
と書ける。ここで $P$ は $s^*$ の「核」である。Lemma
\ref{derham-lemma-relative-bottom-part-degenerates} により、対象 $P$ は $D(A)$ において完全で、
その Tor 振幅は $[1, n]$ に含まれる。Lemma
\ref{derham-lemma-relative-duality-hodge} の証明と同様に、
$H^n(X, \Omega^n_{X/S})$ は局所自由 $A$-加群で階数 $1$ であることが分かる
（実際には $A$ の双対なので階数 $1$ の自由加群である。すぐに生成元を選ぶが、
それが同じ生成元であることを確認したくはなく、その必要もない）。

\medskip\noindent
$Z \subset X$ を $s$ の像とする。Schemes, Lemma
\ref{schemes-lemma-section-immersion} により、これは $X$ の閉部分スキームである。
$Z \to X$ は正則閉埋め込みであり（従って Divisors, Lemma
\ref{divisors-lemma-regular-quasi-regular-immersion} により Koszul 正則でもある）、
これは Divisors, Lemma
\ref{divisors-lemma-section-smooth-regular-immersion} による。
もちろん $Z \to X$ の余次元は $n$ である。したがって Remark
\ref{derham-remark-how-to-use} により、写像
$$
\gamma^{0, 0} : H^0(Z, \Omega^0_{Z/S}) \longrightarrow H^n(X, \Omega^n_{X/S})
$$
を考えることができ、$\xi = \gamma^{0, 0}(1) \in H^n(X, \Omega^n_{X/S})$ と置く。

\medskip\noindent
$\xi$ は基底元であると主張する。実際、最高次において基底変換が成り立つので
（たとえば Limits, Lemma
\ref{limits-lemma-higher-direct-images-zero-above-dimension-fibre} を参照せよ）、
$H^n(X, \Omega^n_{X/S}) \otimes_A k = H^n(X_k, \Omega^n_{X_k/k})$
が任意の環準同型 $A \to k$ に対して成り立つ。$\xi$ の構成は基底変換と両立するので
（Lemma \ref{derham-lemma-gysin-transverse-global} を参照せよ）、
$\xi$ の $H^n(X_k, \Omega^n_{X_k/k})$ における像は、$k$ が体ならば
Lemma \ref{derham-lemma-class-of-a-point} により非零である。したがって $\xi$ は
可逆加群の至る所で消えない切断であり、従って生成元である。

\medskip\noindent
$\theta \in H^n(X, \Omega^{n - 1}_{X/S})$ とする。
$\text{d}(\theta)$ が $H^n(X, \Omega^n_{X/S})$ において零であることを示す必要がある。
$\text{d}(\theta) = a \xi$ と、ある $a \in A$ に対して書ける。これは $\xi$ が
基底元だからである。従って $a = 0$ を示せばよい。

\medskip\noindent
閉埋め込み
$$
\Delta : X \to X \times_S X
$$
を考える。これも滑らかな射（すなわち、いずれかの射影）の切断であり、
従って余次元 $n$ の正則かつ Koszul な埋め込みでもある。よって写像
$$
\gamma^{p, q} :
H^q(X, \Omega^p_{X/S})
\longrightarrow
H^{q + n}(X \times_S X, \Omega^{p + n}_{X \times_S X/S})
$$
を Remark \ref{derham-remark-how-to-use} により考えることができる。像
$$
\gamma^{n - 1, n}(\theta) \in
H^{2n}(X \times_S X, \Omega^{2n - 1}_{X \times_S X})
$$
を考える。Lemma \ref{derham-lemma-de-rham-complex-product} により
$$
\Omega^{2n - 1}_{X \times_S X} =
\Omega^{n - 1}_{X/S} \boxtimes \Omega^n_{X/S} \oplus
\Omega^n_{X/S} \boxtimes \Omega^{n - 1}_{X/S}
$$
である。K\"unneth 公式（Derived Categories of Schemes, Lemma
\ref{perfect-lemma-kunneth} または Derived Categories of Schemes, Lemma
\ref{perfect-lemma-kunneth-single-sheaf}）により
$$
H^{2n}(X \times_S X, \Omega^{n - 1}_{X/S} \boxtimes \Omega^n_{X/S}) =
H^n(X, \Omega^{n - 1}_{X/S}) \otimes_A H^n(X, \Omega^n_{X/S})
$$
であり、また
$$
H^{2n}(X \times_S X, \Omega^n_{X/S} \boxtimes \Omega^{n - 1}_{X/S}) =
H^n(X, \Omega^n_{X/S}) \otimes_A H^n(X, \Omega^{n - 1}_{X/S})
$$
であることが分かる。実際、最高次を見ているので、関与する高次 Tor 群はない。
$\xi$ が生成元であることと併せると、これは
$$
\gamma^{n - 1, n}(\theta) = \theta_1 \otimes \xi + \xi \otimes \theta_2
$$
と書けることを意味する。ここで
$\theta_1, \theta_2 \in H^n(X, \Omega^{n - 1}_{X/S})$ である。
まったく同様に論じると
$$
\gamma^{n, n}(\xi) = b \xi \otimes \xi
$$
と書ける。これは
$H^{2n}(X \times_S X, \Omega^{2n}_{X \times_S X/S}) =
H^n(X, \Omega^n_{X/S}) \otimes_A H^n(X, \Omega^n_{X/S})$
における等式である。
を、ある $b \in H^0(S, \mathcal{O}_S)$ に対して得る。

\medskip\noindent
{\bf 主張：} $\theta_1 = \theta$, $\theta_2 = \theta$, かつ $b = 1$ である。
この主張から所望の結論 $a = 0$ が従うことを示そう。実際、Lemma
\ref{derham-lemma-gysin-differential-hodge} により
$$
\gamma^{n, n}(\text{d}(\theta)) = \text{d}(\gamma^{n - 1, n}(\theta))
$$
である。上の選択により、これは
$$
a \xi \otimes \xi =
\gamma^{n, n}(a\xi) =
\text{d}(\theta \otimes \xi + \xi \otimes \theta) =
a \xi \otimes \xi + (-1)^n a \xi \otimes \xi
$$
を与える。右端の等式は、写像
$\text{d} : \Omega^{2n - 1}_{X \otimes_S X/S} \to \Omega^{2n}_{X \times_S X/S}$
が Lemma \ref{derham-lemma-de-rham-complex-product} により微分
$\text{d} \boxtimes 1 : \Omega^{n - 1}_{X/S} \boxtimes \Omega^n_{X/S}
\to \Omega^n_{X/S} \boxtimes \Omega^n_{X/S}$
と微分
$(-1)^n 1 \boxtimes \text{d} : \Omega^n_{X/S} \boxtimes \Omega^{n - 1}_{X/S}
\to \Omega^n_{X/S} \boxtimes \Omega^n_{X/S}$ の和だからである。詳しくは
Section \ref{derham-section-kunneth} および Derived Categories of Schemes, Section
\ref{perfect-section-kunneth-complexes} の議論を参照せよ。
$\xi \otimes \xi$ は階数 $1$ の自由 $A$-加群
$H^n(X, \Omega^n_{X/S}) \otimes_A H^n(X, \Omega^n_{X/S})$
の基底なので、
$$
a = a + (-1)^n a \Rightarrow a = 0
$$
と結論し、所望の結果を得る。

\medskip\noindent
証明の残りで上の主張を証明する。
$\eta = \gamma^{0, 0}(1) \in H^n(X \times_S X, \Omega^n_{X \times_S X/S})$.
と記す。$\Omega^n_{X \times_S X/S} =
\bigoplus_{p + p' = n} \Omega^p_{X/S} \boxtimes \Omega^{p'}_{X/S}$
なので、
$$
\eta = \eta_0 + \eta_1 + \ldots + \eta_n
$$
と書ける。ここで $\eta_p$ は
$H^n(X \times_S X, \Omega^p_{X/S} \boxtimes \Omega^{n - p}_{X/S})$.
の元である。$p = 0$ に対しては
\begin{align*}
H^n(X \times_S X, \mathcal{O}_X \boxtimes \Omega^n_{X/S})
& =
H^n(R\Gamma(X, \mathcal{O}_X) \otimes_A^\mathbf{L}
R\Gamma(X, \Omega^n_{X/S})) \\
& =
A \otimes_A H^n(X, \Omega^n_{X/S}) \oplus
H^n(P \otimes_A^\mathbf{L} R\Gamma(X, \Omega^n_{X/S}))
\end{align*}
と書ける。これは先に与えた分解 $R\Gamma(X, \mathcal{O}_X) = A \oplus P$
による。射 $(s, \text{id}) : X \to X \times_S X$ を考える。このとき
スキーム論的に $(s, \text{id})^{-1}(\Delta) = Z$ である。従って Lemma
\ref{derham-lemma-gysin-transverse-global} により
$(s, \text{id})^*\eta = \xi$ であることが分かる。これは
$$
\xi = (s, \text{id})^*\eta = (s^* \otimes \text{id})(\eta_0)
$$
を意味する。すなわち、上の直和分解における $\eta_0$ の第一成分は
$\xi$ である。言い換えると、
$$
\eta_0 = 1 \otimes \xi + \eta'_0
$$
と書ける。ここで
$\eta'_0 \in H^n(P \otimes_A^\mathbf{L} R\Gamma(X, \Omega^n_{X/S}))$ である。
まったく同様に、$p = n$ に対して
\begin{align*}
H^n(X \times_S X, \Omega^n_{X/S} \boxtimes \mathcal{O}_X)
& =
H^n(R\Gamma(X, \Omega^n_{X/S}) \otimes_A^\mathbf{L}
R\Gamma(X, \mathcal{O}_X)) \\
& =
H^n(X, \Omega^n_{X/S}) \otimes_A A \oplus
H^n(R\Gamma(X, \Omega^n_{X/S}) \otimes_A^\mathbf{L} P)
\end{align*}
と書け、さらに
$$
\eta_n = \xi \otimes 1 + \eta'_n
$$
と書ける。ここで
$\eta'_n \in H^n(R\Gamma(X, \Omega^n_{X/S}) \otimes_A^\mathbf{L} P)$ である。

\medskip\noindent
$\text{pr}_1^*\theta = \theta \otimes 1$ および
$\text{pr}_2^*\theta = 1 \otimes \theta$ は $X \times_S X$ 上の
Hodge コホモロジー類であり、$\theta$ に引き戻され、その引き戻しは $\Delta$ によることに
注意する。従って Lemma \ref{derham-lemma-gysin-projection-global} により
$$
\theta_1 \otimes \xi + \xi \otimes \theta_2 =
\gamma^{n - 1, n}(\theta) =
(\theta \otimes 1) \cup \eta =
(1 \otimes \theta) \cup \eta
$$
が $X \times_S X$ の $S$ 上の Hodge コホモロジー環において成り立つ。
$X \times_S X/S$ の微分加群の直和分解に関して、次を得る：
$$
\theta_1 \otimes \xi =
(\theta \otimes 1) \cup \eta_0
\quad\text{and}\quad
\xi \otimes \theta_2 =
(1 \otimes \theta) \cup \eta_n
$$
上で得た公式 $\eta_0 = 1 \otimes \xi + \eta'_0$ を見ると、
$\theta_1 = \theta$ を示すには
$$
(\theta \otimes 1) \cup \eta'_0 = 0
$$
を証明すれば十分であることが分かる。そのために、$\theta \otimes 1$ との
カップ積は、写像
$$
(P \otimes_A^\mathbf{L} R\Gamma(X, \Omega^n_{X/S}))[-n]
\xrightarrow{\theta \otimes 1}
R\Gamma(X, \Omega^{n - 1}_{X/S}) \otimes_A^\mathbf{L}
R\Gamma(X, \Omega^n_{X/S})
$$
が導来圏においてコホモロジーに誘導する作用で与えられることに注意する。
Cohomology, Remark
\ref{cohomology-remark-cup-with-element-map-total-cohomology} を参照せよ。
この写像は二つの写像
$$
\theta : P[-n] \to R\Gamma(X, \Omega^{n - 1}_{X/S})
\quad\text{and}\quad
1 : R\Gamma(X, \Omega^n_{X/S}) \to R\Gamma(X, \Omega^n_{X/S})
$$
の導来テンソル積である。これは Derived Categories of Schemes, Remark
\ref{perfect-remark-annoying-compatibility} による。しかし前者は $D(A)$ において
零である。なぜなら、これは Tor 振幅が $[n + 1, 2n]$ に含まれる完全複体から、
コホモロジーが次数 $0, 1, \ldots, n$ のみにある複体への写像だからである。
More on Algebra, Lemma \ref{more-algebra-lemma-splitting-unique} を参照せよ。
同様の議論により $(1 \otimes \theta) \cup \eta'_n$ の消滅も示せる。最後に、
まったく同様に
$$
b \xi \otimes \xi = \gamma^{n, n}(\xi) = (\xi \otimes 1) \cup \eta_0
$$
を得る。そして上と同様に $(\xi \otimes 1) \cup \eta'_0 = 0$ を示すことで、
先と同じく結論を得る。これで証明が完了する。
\end{proof}

\begin{proposition}
\label{derham-proposition-relative-poincare-duality}
$S$ を準コンパクトかつ準分離なスキームとする。$f : X \to S$ を
スキームの固有滑らかな射で、そのすべてのファイバーが空でなく、
次元 $n$ の等次元であるものとする。$\mathcal{O}_S$-加群の写像
$$
t : R^{2n}f_*\Omega^\bullet_{X/S} \longrightarrow \mathcal{O}_S
$$
で次の性質をもつものが存在し、$H^0(X, \mathcal{O}_X)$ の単元による乗法写像を
前合成することを除いて一意である。すなわち、ペアリング
$$
Rf_*\Omega^\bullet_{X/S}
\otimes_{\mathcal{O}_S}^\mathbf{L}
Rf_*\Omega^\bullet_{X/S}[2n]
\longrightarrow
\mathcal{O}_S, \quad
(\xi, \xi') \longmapsto t(\xi \cup \xi')
$$
は $S$ 上の完全複体の完全なペアリングである。
\end{proposition}

\begin{proof}
証明は Proposition \ref{derham-proposition-poincare-duality} の証明とまったく同じである。

\medskip\noindent
相対 Hodge-to-de Rham スペクトル系列
$$
E_1^{p, q} = R^qf_*\Omega^p_{X/S} \Rightarrow R^{p + q}f_*\Omega^\bullet_{X/S}
$$
（Section \ref{derham-section-hodge-to-de-rham}）、$\Omega^i_{X/S}$ の
$i > n$ に対する消滅、たとえば Limits, Lemma
\ref{limits-lemma-higher-direct-images-zero-above-dimension-fibre} における消滅、
および Lemmas \ref{derham-lemma-relative-bottom-part-degenerates} と
\ref{derham-lemma-relative-top-part-degenerates} の結果により、
$R^0f_*\Omega_{X/S} = R^0f_*\mathcal{O}_X$ および
$R^nf_*\Omega^n_{X/S} = R^{2n}f_*\Omega^\bullet_{X/S}$ であることが分かる。
より正確には、これらの同一視は複体の写像
$$
\Omega^\bullet_{X/S} \to \mathcal{O}_X[0]
\quad\text{and}\quad
\Omega^n_{X/S}[-n] \to \Omega^\bullet_{X/S}
$$
から生じる。$t : R^{2n}f_*\Omega_{X/S} \to \mathcal{O}_S$ を選び、これが
この同一視を通じて Lemma \ref{derham-lemma-relative-duality-hodge} における
$t$ に対応するものとする。

\medskip\noindent
$\Omega^\bullet = \Omega^\bullet_{X/S}$ と略記する。
写像 (\ref{derham-equation-wedge}) は、この状況では
$$
\wedge :
\text{Tot}(\Omega^\bullet \otimes_{f^{-1}\mathcal{O}_S} \Omega^\bullet)
\longrightarrow
\Omega^\bullet
$$
と書ける。任意の整数 $p = 0, 1, \ldots, n$ に対し、この写像は次数の理由から部分複体
$\text{Tot}(\sigma_{> p} \Omega^\bullet \otimes_{f^{-1}\mathcal{O}_S}
\sigma_{\geq n - p} \Omega^\bullet)$ を零にする。
従って、$\wedge$ の部分複体
$\text{Tot}(\Omega^\bullet \otimes_{f^{-1}\mathcal{O}_S}
\geq_{n - p}\Omega^\bullet)$ への制限は、複体の写像
$$
\gamma_p :
\text{Tot}(\sigma_{\leq p} \Omega^\bullet \otimes_{f^{-1}\mathcal{O}_S}
\sigma_{\geq n - p} \Omega^\bullet)
\longrightarrow
\Omega^\bullet
$$
を経由することが分かる。Section \ref{derham-section-cup-product} と同じ手順を用いると、
相対カップ積
$$
Rf_*\sigma_{\leq p} \Omega^\bullet
\otimes_{\mathcal{O}_S}^\mathbf{L}
Rf_*\sigma_{\geq n - p}\Omega^\bullet
\longrightarrow
Rf_*\Omega^\bullet
$$
を得る。$p$ に関する帰納法で、これらのカップ積が $t$ を介して
$Rf_*\sigma_{\leq p} \Omega^\bullet$ と
$Rf_*\sigma_{\geq n - p}\Omega^\bullet[2n]$ の間に完全なペアリングを誘導することを
証明する。$p = n$ の場合が命題の主張である。

\medskip\noindent
基礎段階は $p = 0$ である。この場合、
$$
Rf_*\sigma_{\leq p}\Omega^\bullet = Rf_*\mathcal{O}_X
\quad\text{and}\quad
Rf_*\sigma_{\geq n - p}\Omega^\bullet[2n] = Rf_*(\Omega^n[-n])[2n] =
Rf_*\Omega^n[n]
$$
である。この場合、Lemma \ref{derham-lemma-relative-duality-hodge} における
$Rf_*\mathcal{O}_X$ と $Rf_*\Omega^n[n]$ の間のペアリングがそのまま得られ、
結果は真である。

\medskip\noindent
帰納段階。結果が $p$ に対して真であることが分かっているとする。このとき区別三角形
$$
\Omega^{p + 1}[-p - 1] \to
\sigma_{\leq p + 1}\Omega^\bullet \to
\sigma_{\leq p}\Omega^\bullet \to
\Omega^{p + 1}[-p]
$$
と区別三角形
$$
\sigma_{\geq n - p}\Omega^\bullet \to
\sigma_{\geq n - p - 1}\Omega^\bullet \to
\Omega^{n - p - 1}[-n + p + 1] \to
(\sigma_{\geq n - p}\Omega^\bullet)[1]
$$
を考える。両者は $f^{-1}\mathcal{O}_S$-加群の層の複体のホモトピー圏における
区別三角形であることに注意する。特に写像
$\sigma_{\leq p}\Omega^\bullet \to \Omega^{p + 1}[-p]$ と
$\Omega^{n - p - 1}[-d + p + 1] \to (\sigma_{\geq n - p}\Omega^\bullet)[1]$
は実際の複体の写像で与えられ、具体的には微分
$\Omega^p \to \Omega^{p + 1}$ と微分
$\Omega^{n - p - 1} \to \Omega^{n - p}$ を用いる。
これらの区別三角形に $Rf_*$ を適用して得られる付随する区別三角形を考える：
$$
\xymatrix{
Rf_*\sigma_{\leq p}\Omega^\bullet \ar[d]_a \\
Rf_*\Omega^{p + 1}[-p - 1] \ar[d]_b \\
Rf_*\sigma_{\leq p + 1}\Omega^\bullet \ar[d]_c \\
Rf_*\sigma_{\leq p}\Omega^\bullet \ar[d]_d \\
Rf_*\Omega^{p + 1}[-p - 1]
}
\quad\quad
\xymatrix{
Rf_*\sigma_{\geq n - p}\Omega^\bullet \\
Rf_*\Omega^{n - p - 1}[-n + p + 1] \ar[u]_{a'} \\
Rf_*\sigma_{\geq n - p - 1}\Omega^\bullet \ar[u]_{b'} \\
Rf_*\sigma_{\geq n - p}\Omega^\bullet \ar[u]_{c'} \\
Rf_*\Omega^{n - p - 1}[-n + p + 1] \ar[u]_{d'}
}
$$
対 $(a, a')$, $(b, b')$, $(c, c')$, および $(d, d')$ が与えられた
ペアリングと両立することを以下で示す。これは、左の区別三角形から、右の
区別三角形に $R\SheafHom(-, \mathcal{O}_S)$ を適用して得られる区別三角形への
写像を得ることを意味する。帰納法と Lemma \ref{derham-lemma-duality-hodge} により、
第一、第二、第四、第五行の複体の間に上で構成したペアリングは完全、すなわち、
双対を取ると同型を定めることが分かっている。Derived Categories, Lemma
\ref{derived-lemma-third-isomorphism-triangle} により、中段の複体の間の
ペアリングも所望どおり完全であると結論する。

\medskip\noindent
$e : K \to K'$ と $e' : M' \to M$ を $D(\mathcal{O}_S)$ の対象の写像とし、
$K \otimes_{\mathcal{O}_S}^\mathbf{L} M \to \mathcal{O}_S$ と
$K' \otimes_{\mathcal{O}_S}^\mathbf{L} M' \to \mathcal{O}_S$
をペアリングとする。このとき、図式
$$
\xymatrix{
K' \otimes_{\mathcal{O}_S}^\mathbf{L} M' \ar[d] &
K \otimes_{\mathcal{O}_S}^\mathbf{L} M'
\ar[l]^{e \otimes 1} \ar[d]^{1 \otimes e'} \\
\mathcal{O}_S &
K \otimes_{\mathcal{O}_S}^\mathbf{L} M \ar[l]
}
$$
が可換であるとき、これらのペアリングは両立すると言う。これは実際、図式
$$
\xymatrix{
K \ar[r] \ar[d]_e & R\SheafHom(M, \mathcal{O}_S) \ar[d]^{R\SheafHom(e', -)} \\
K' \ar[r] & R\SheafHom(M', \mathcal{O}_S)
}
$$
が可換であることを意味するので、我々の目的には十分である。

\medskip\noindent
これを対 $(c, c')$ に対して証明しよう。ここでは単に可換図式
$$
\xymatrix{
\text{Tot}(\sigma_{\leq p} \Omega^\bullet \otimes_{f^{-1}\mathcal{O}_S}
\sigma_{\geq n - p} \Omega^\bullet) \ar[d]_{\gamma_p} &
\text{Tot}(\sigma_{\leq p + 1} \Omega^\bullet \otimes_{f^{-1}\mathcal{O}_S}
\sigma_{\geq n - p} \Omega^\bullet) \ar[l] \ar[d] \\
\Omega^\bullet &
\text{Tot}(\sigma_{\leq p + 1} \Omega^\bullet \otimes_{f^{-1}\mathcal{O}_S}
\sigma_{\geq n - p - 1} \Omega^\bullet) \ar[l]_-{\gamma_{p + 1}}
}
$$
を得ることに注意する。カップ積の関手性により、所望の図式の可換性を得る。

\medskip\noindent
同様に、対 $(b, b')$ に対して可換図式
$$
\xymatrix{
\text{Tot}(\sigma_{\leq p + 1} \Omega^\bullet \otimes_{f^{-1}\mathcal{O}_S}
\sigma_{\geq n - p - 1} \Omega^\bullet) \ar[d]_{\gamma_{p + 1}} &
\text{Tot}(\Omega^{p + 1}[-p - 1] \otimes_{f^{-1}\mathcal{O}_S}
\sigma_{\geq n - p - 1} \Omega^\bullet) \ar[l] \ar[d] \\
\Omega^\bullet &
\Omega^{p + 1}[-p - 1]
\otimes_{f^{-1}\mathcal{O}_S}
\Omega^{n - p - 1}[-n + p + 1] \ar[l]_-\wedge
}
$$
を用いる。

\medskip\noindent
対 $(d, d')$ と $(a, a')$ に対しては可換図式
$$
\xymatrix{
\Omega^{p + 1}[-p] \otimes_{f^{-1}\mathcal{O}_S}
\Omega^{n - p - 1}[-n + p] \ar[d] &
\text{Tot}(\sigma_{\leq p}\Omega^\bullet \otimes_{f^{-1}\mathcal{O}_S}
\Omega^{n - p - 1}[-n + p]) \ar[l] \ar[d] \\
\Omega^\bullet &
\text{Tot}(\sigma_{\leq p}\Omega^\bullet \otimes_{f^{-1}\mathcal{O}_S}
\sigma_{\geq n - p}\Omega^\bullet) \ar[l]
}
$$
を用いる。

\medskip\noindent
$t$ が $H^0(X, \mathcal{O}_X)$ の単元による乗法写像を前合成することを除いて
一意であることを示す議論は省略する。
\end{proof}
